\chapter{Denominator identities, replication and genus zero}\label{ch:19}

\providecommand{\chnineteenH}{\mathbb{H}}
\providecommand{\chnineteenR}{\mathbb{R}}
\providecommand{\chnineteenSL}{\mathrm{SL}_2(\mathbb{Z})}
\providecommand{\chnineteenSLR}{\mathrm{SL}_2(\mathbb{R})}
\providecommand{\chnineteenGLR}{\mathrm{GL}_2(\mathbb{R})}
\providecommand{\chnineteenP}{\mathbb{P}^1}
\providecommand{\chnineteenm}{\mathfrak{m}}
\providecommand{\chnineteenu}{\mathfrak{u}}
\providecommand{\chnineteenmat}[4]{\left(\begin{smallmatrix}#1&#2\\#3&#4\end{smallmatrix}\right)}
\providecommand{\chnineteenlean}[1]{\texttt{\directlua{ local s = "\luaescapestring{\detokenize{#1}}" for piece, sep in string.gmatch(s, "([^_.]*)([_.]?)") do if not (piece == "" and sep == "") then tex.sprint(-2, piece .. sep) if not (sep == "") then tex.sprint(string.char(92) .. "allowbreak{}") end end end}}}
\chaptermark{Denominator identities, replication, genus zero}
\begingroup\setlength{\emergencystretch}{3em}\sloppy

This chapter documents the release of the library named in the front matter; every Lean quotation is checked by
\texttt{scripts/textbook/sync.py} to occur verbatim (up to whitespace and comments) in the source files of that
release.  The chapter documents the sixty files that \texttt{Notes/Textbook/assignment.txt} assigns to it: the files
\texttt{MoonshineCNTwisted*}, \texttt{MoonshineCNReplication*} and \texttt{MoonshineCNGenusZero*} written on Fedor
Popov's side of the project, and the files \texttt{MoonshineLRM5*}, \texttt{MoonshineLRHecke*},
\texttt{MoonshineLRFaber*}, \texttt{MoonshineLRJUnbounded*}, \texttt{MoonshineLRNontrivial*},
\texttt{MoonshineLRPropertyP}, \texttt{MoonshineLRTower*} and \texttt{MoonshineLRCoeffBound*} written on Leonardo
Rastelli's side, all in the folder \leanpath{VOA/Moonshine/}.

\paragraph{Where we are.}  Chapter~\ref{ch:16} stated the Conway--Norton conjecture in Lean
(\lean{ConwayNortonConjecture}), defined complete replicability (\lean{CompletelyReplicable}), the abstract interface
\lean{ReplicableFamily}, and proved the step from groups containing $\Gamma_1(N)$ to groups containing $\Gamma_0(N)$.
Chapters~\ref{ch:17} and~\ref{ch:18} built the Monster Lie algebra $\chnineteenm$ from $\Vnat\otimes V_{I\!I_{1,1}}$, proved
the no-ghost theorem $\chnineteenm_{(m,n)}\cong\Vnat_{1+mn}$, and showed that $\chnineteenm$ is a Borcherds datum with
$T(U)\cong U(\chnineteenu^+)$, \emph{given} that $\chnineteenm$ is generated by its simple root spaces (the proposition
\lean{MonsterLieSimpleRootGeneration}, proved in Chapter~\ref{ch:20}).  The present chapter connects these two ends.
It proves, in Lean:
\begin{enumerate}[label=(\Roman*)]
\item \textbf{the product identity for $J$} (the denominator identity of the Monster Lie algebra, but stated and
  proved in Lean as an identity about $J$ alone; no theorem of this chapter mentions $\chnineteenm$ in it)
  \begin{equation}\label{ch19:eq:knz-intro}
    \prod_{m,n\ge1}\bigl(1-p^mq^n\bigr)^{c(mn)}=1-\sum_{a,b\ge1}c(a+b-1)\,p^aq^b ,
  \end{equation}
  for the coefficients $c(n)$ of $J=j-744$, as an identity of formal series (\lean{knz_product}; the product is read as
  \lean{coneProd}, the series with the product's logarithmic derivative, Definition~\ref{ch19:def:coneprod}), by a route that
  never mentions Lie algebras: it is derived from the replicability of $J$, which in turn is proved from the
  $\chnineteenSL$-invariance of $j$;
\item \textbf{the twisted denominator identities}: for every Borcherds datum with a finite-order symmetry $\sigma$,
  and in particular for every finite-order automorphism of $\Vnat$ acting on $\chnineteenm$, a logarithm-free identity
  of two-variable formal series between the McKay--Thompson coefficients of $\sigma$ and of its powers
  (\lean{twistedDenominatorIdentity_of_finOrder}, \lean{twistedDenominatorIdentity_of_simpleRootGeneration});
\item \textbf{replication}: the twisted identities imply, formally and then analytically, that every finite-order
  McKay--Thompson series is completely replicable (\lean{faberPoly_replication}, \lean{faberPoly_eval_mckayThompsonFun},
  \lean{completelyReplicable_of_simpleRootGeneration});
\item \textbf{genus zero}: for a replicable family of positive period that satisfies the prime Hecke identity and
  nontriviality (explicit hypotheses \lean{hB1}, \lean{hD7}), every member is a normalised Hauptmodul of a group of
  determinant one that contains some $\Gamma_1(N')$ and is commensurable with $\chnineteenSL$
  (\lean{ReplicableFamily.exists_hauptmodul_gammaOne}).  The two inputs are separate theorems, proved for every family
  (\lean{ReplicableFamily_primeHeckeIdentity}, \lean{ReplicableFamily_notInjectiveModTranslations}); the composite
  ``for every family of positive period'' is \emph{not} packaged as a single Lean declaration.  Hence the $\Gamma_1$ form of
  Conway--Norton follows from complete replicability (\lean{conwayNortonGammaOne_of_completelyReplicable}), and the full
  conjecture follows from complete replicability and finite order (\lean{conwayNorton_of_completelyReplicable}).
\end{enumerate}
Part (IV) is the longest part of the chapter, and mathematically the most interesting.  It is a \emph{case-free} proof
that complete replicability forces the genus-zero property.  It is modelled on the work of Cummins and Gannon (1997)
and Cummins (2002), but it is a different proof: the analytic continuation of local symmetries, which is the heart of
Cummins--Gannon, is replaced by a real linear ordinary differential equation satisfied by every local symmetry, and
several steps of the published proofs that fail as printed are replaced by corrected arguments (\S\ref{ch19:sec:compare},
\S\ref{ch19:sec:traps}).

\paragraph{What is conditional and what is not.}  Every theorem of this chapter is \texttt{\#print axioms}-clean
(only \lean{propext}, \lean{Classical.choice}, \lean{Quot.sound}).  Very few of them are hypothesis-free; the
\emph{unconditional} ones (no premises at all) are the replicability of $J$, the product identity \lean{knz_product}
and its logarithmic form, the prime Hecke identity, the unboundedness of $J$'s coefficients, and (given only $0<N$)
nontriviality and property P.  All others carry explicit premises, listed in each status box: analytic lemmas assume
\lean{IsJDominated c}; the abstract twisted identity assumes eigenbases or finite order; the genus-zero steps assume
\lean{hN : 0 < N}, the prime Hecke identity (\lean{hB1}), nontriviality (\lean{hD7}), property P (\lean{hP}) or
$\Gamma_1(M)\le G$, depending on the step.  The premises \lean{hB1}, \lean{hD7}, \lean{hP} are discharged by theorems of this
same chapter, but the composites are written out in Lean only for the McKay--Thompson families (Chapter~\ref{ch:20}).
The Monster instances of parts (II) and (III)
take the hypothesis \lean{hgen : MonsterLieSimpleRootGeneration} (and, for the statement over all of
$\Aut(\Vnat)$, \lean{hfin : ∀ g : AutVnat, IsOfFinOrder g}); the Conway--Norton consequences also take
\lean{hrep : CompletelyReplicable}.  These are discharged by theorems of
Chapter~\ref{ch:20}, and the hypothesis-free consequences \lean{completelyReplicable_holds} and
\lean{conwayNortonGammaOne_holds} are recorded in Chapter~\ref{ch:22}.  We say ``conditional on \lean{hgen}'' whenever
a statement of this chapter carries that hypothesis, even though it is later discharged.

\paragraph{How to read this chapter.}  Section~\ref{ch19:sec:overview} gives the logical chain and fixes conventions.
Sections~\ref{ch19:sec:faber}--\ref{ch19:sec:denominator} are part (I), Section~\ref{ch19:sec:twisted} is part (II),
Section~\ref{ch19:sec:replication} is part (III), and Sections~\ref{ch19:sec:gz}--\ref{ch19:sec:main} are part (IV).
Section~\ref{ch19:sec:compare} compares the genus-zero proof with the literature; Section~\ref{ch19:sec:traps} collects
the recorded traps; the chapter ends with the audit modules, the referee's checklist and the references.  As
everywhere in this book, the concern is the \emph{meaning} of the Lean statements; proofs are explained mathematically,
and we do not describe how the Lean proof scripts are organised.

%======================================================================
\section{The logical chain and the conventions}\label{ch19:sec:overview}
%======================================================================

\subsection{The chain of implications}

The following diagram shows the implications proved in this chapter.  Arrows are Lean theorems; a label in brackets
is an explicit hypothesis of the Lean theorem.
\begin{center}\small
\begin{tabular}{p{0.44\textwidth}p{0.48\textwidth}}
\toprule
statement & proved from\\
\midrule
$J$ replicable: $P_M(J)=\sum_{ad=M}\sum_{b<d}J(\frac{a\tau+b}d)$ (\chnineteenlean{faberPoly_cJ_eval_eq_qHeckeSum}) &
  $\chnineteenSL$-invariance of $j$; Hecke cosets; Faber polynomials\\
product identity \eqref{ch19:eq:knz-intro} (\chnineteenlean{knz_product}) & replicability of $J$; Faber generating identity\\
twisted identity $D\tilde u=(1-\tilde u)\tilde P$ for a symmetry of a Borcherds datum
  (\chnineteenlean{twistedDenominatorIdentity_of_finOrder}) & $T(U)\cong U(\chnineteenu^+)$ (Ch.~\ref{ch:18}); PBW; finite order\\
the same for $\sigma\in\Aut(\Vnat)$ on $\chnineteenm$ (\chnineteenlean{twistedDenominatorIdentity_of_simpleRootGeneration}) &
  [\chnineteenlean{hgen}], finite order of $\sigma$, the no-ghost isomorphisms (Ch.~\ref{ch:18})\\
formal replication (\chnineteenlean{faberPoly_replication}) & a twisted identity\\
analytic replication of $T_g$ (\chnineteenlean{faberPoly_eval_mckayThompsonFun}) & twisted identity of $g$; finite order\\
\chnineteenlean{CompletelyReplicable} (\chnineteenlean{completelyReplicable_of_simpleRootGeneration}) & [\chnineteenlean{hgen}], [\chnineteenlean{hfin}]\\
prime Hecke identity for every family (\chnineteenlean{ReplicableFamily_primeHeckeIdentity}) & periodicity of the family\\
nontriviality for every family (\chnineteenlean{ReplicableFamily_notInjectiveModTranslations}) & replication; $J$'s coefficients
  are unbounded\\
genus zero for every family (\chnineteenlean{ReplicableFamily.exists_hauptmodul_gammaOne}) & [prime Hecke identity],
  [nontriviality]\\
\chnineteenlean{ConwayNortonGammaOne} (\chnineteenlean{conwayNortonGammaOne_of_completelyReplicable}) & [\chnineteenlean{hrep}], [prime Hecke],
  [nontriviality]\\
\chnineteenlean{ConwayNortonConjecture} (\chnineteenlean{conwayNorton_of_completelyReplicable}) & the above, [\chnineteenlean{hfin}], and the
  $\Gamma_1\Rightarrow\Gamma_0$ bridge of Ch.~\ref{ch:16}\\
\bottomrule
\end{tabular}
\end{center}
The bracketed inputs ``prime Hecke'' and ``nontriviality'' of the last three rows are discharged for \emph{every}
replicable family by the two theorems of the eighth and ninth rows; the file \leanpath{VOA/Moonshine/MoonshineLRConwayNorton.lean}
(Chapter~\ref{ch:20}) composes them into \lean{M8.conwayNorton_of_finOrder_of_completelyReplicable (hfin) (hrep)}.

\begin{physics}
The chain is Borcherds' proof read as string theory.  The product identity is the statement that the physical states
of the bosonic string compactified on the moonshine orbifold times a two-dimensional Lorentzian torus have the
multiplicities $c(mn)$ (the no-ghost theorem) and that the second-quantised string Fock space they generate is
``free up to one relation''.  The twisted identities are the same statement with a symmetry $g$ inserted in the trace.
Replication says that the $g$-twisted torus partition functions transform among themselves under the Hecke
correspondences (passing to $m$-sheeted covers of the torus).  Genus zero is the rigid analytic consequence of these
covering relations.  None of these physical readings is used as a proof.
\end{physics}

\subsection{Conventions for \texorpdfstring{$q$}{q}-series}\label{ch19:sec:qconv}

Throughout, $q=e^{2\pi i\tau}$ for $\tau\in\chnineteenH$ (the function \lean{qPar} of Chapter~\ref{ch:16}), and
\begin{equation*}
  J(\tau)=j(\tau)-744=\sum_{n\ge-1}c(n)q^n=q^{-1}+196884\,q+21493760\,q^2+\cdots,\qquad c(-1)=1,\ c(0)=0 .
\end{equation*}
Chapter~\ref{ch:05} defined the formal series \lean{JShift} $=qJ$, so that $[q^n]\mathtt{JShift}=c(n-1)$.  Two
encodings of a normalised $q$-series $f=q^{-1}+\sum_{n\ge1}c(n)q^n$ occur:
\begin{itemize}
\item as a \emph{function} on $\chnineteenH$: \lean{qSeriesFun c τ} $=q^{-1}+\sum_{n\ge0}c(n+1)q^{n+1}$
  (Chapter~\ref{ch:16}, the definition of replicable families); the value \lean{c 0} is ignored;
\item as a \emph{formal Laurent series}: \lean{qSeriesLaurent c} $\in\C((q))$ (Definition~\ref{ch19:def:faber} below);
  again \lean{c 0} is ignored.
\end{itemize}
The coefficient sequence of $J$ in this encoding is
\begin{leancode}
noncomputable def cJ (n : ℕ) : ℂ := PowerSeries.coeff (n + 1) JShift
\end{leancode}
(\leanpath{VOA/Moonshine/MoonshineLRM5Replicable.lean}, namespace \lean{VOA.Moonshine.ConwayNorton.M8}), so
\lean{cJ n} $=c(n)$ for $n\ge1$ and \lean{cJ 0} $=c(0)=0$.  The analytic results need a uniform bound on the
coefficients; the library's notion is domination by $J$:
\begin{leancode}
noncomputable def jCoeffMajorant (n : ℕ) : ℝ := ‖PowerSeries.coeff (n + 1 + 1) JShift‖
noncomputable def jTail (y : ℝ) : ℝ :=
  ∑' n : ℕ, jCoeffMajorant n * Real.exp (-2 * π * (n + 1) * y)
def IsJDominated (c : ℕ → ℂ) : Prop :=
  ∀ n : ℕ, ‖c (n + 1)‖ ≤ jCoeffMajorant n
\end{leancode}
(\leanpath{VOA/Moonshine/MoonshineCNGenusZeroAnalytic.lean}, namespace \lean{VOA.Moonshine.ConwayNorton}.)
\begin{dictionary}
\lean{jCoeffMajorant n} $=\|[q^{n+2}]\mathtt{JShift}\|=|c(n+1)|=c(n+1)$ (the coefficients of $J$ are positive integers
for $n\ge1$, but only their absolute values enter).  \lean{IsJDominated c} says $|c(n)|\le c_J(n)$ for every $n\ge1$,
where $c_J(n)=\mathtt{cJ}\ n$ are the coefficients of $J$.  \lean{jTail y}
$=\sum_{n\ge1}c(n)e^{-2\pi ny}$ is the tail of $J$ on the imaginary axis: $J(iy)=e^{2\pi y}+\mathtt{jTail}(y)$.  The
\lean{tsum} (\lean{∑'}) is Mathlib's unconditional sum, which is $0$ if the family is not summable; for $y>0$ the
series converges (\lean{summable_jCoeffMajorant_mul_exp}), so the junk value never occurs where it matters.
\end{dictionary}
\begin{physics}
$J$ is the torus partition function of the moonshine CFT with the vacuum energy $-c/24=-1$ removed.
\lean{IsJDominated c} says that a ``twisted partition function'' $q^{-1}+\sum c(n)q^n$ never has more states (in
absolute value of its trace) at level $n$ than the untwisted one: this is the bound $|\tr(g|V_n)|\le\dim V_n$ for a
unitary (finite-order) symmetry.  \lean{jTail} is the contribution of all excited states at inverse temperature
$2\pi y$.
\end{physics}
\begin{status}
Definitions; standard axioms (\leanpath{VOA/Moonshine/MoonshineCNGenusZeroAudit.lean} prints the lemmas about them,
and \leanpath{VOA/Moonshine/MoonshineLRM5Audit.lean} prints \lean{isJDominated_cJ}).
\end{status}

\subsection{Conventions for formal series in two variables}\label{ch19:sec:formalconv}

The denominator identities live in rings of formal power series in two variables $p$ and $q$.  The library uses two
models, always with $p$ as the \emph{outer} variable:
\begin{itemize}
\item \lean{PowerSeries (LaurentSeries ℂ)}: power series in $p$ whose coefficients are formal Laurent series in $q$
  (finitely many negative powers allowed).  This is $\C((q))[[p]]$.  It is used for the untwisted identity, where
  $J(q)$ has a pole.  Mathlib's \lean{LaurentSeries ℂ} is \lean{HahnSeries ℤ ℂ}; \lean{HahnSeries.single k a} is
  the monomial $aq^k$, and \lean{HahnSeries.ofPowerSeries ℤ ℂ} embeds $\C[[q]]$ into $\C((q))$.
\item \lean{PowerSeries (PowerSeries ℂ)}: $\C[[q]][[p]]$, used for the twisted identities, which are written in
  ``height coordinates'' where no negative powers occur (\S\ref{ch19:sec:twisted}).
\end{itemize}
\lean{PowerSeries.X} is the outer variable $p$; \lean{PowerSeries.C a} the constant $a$; \lean{PowerSeries.coeff n F}
the coefficient of $p^n$ (so \lean{PowerSeries.coeff N (PowerSeries.coeff A F)} is the coefficient of $p^Aq^N$ in
$F\in\C[[q]][[p]]$); \lean{PowerSeries.derivative R F} the formal derivative $\partial_pF$; \lean{PowerSeries.map φ F}
applies a ring map to every coefficient; \lean{PowerSeries.rescale a F} substitutes $p\mapsto ap$.  For a polynomial
$P\in\C[X]$ and an element $x$ of a $\C$-algebra, \lean{aeval x P} is $P(x)$.  The operator $D=p\,\partial_p$ (the
\emph{Euler derivation}) appears either as \lean{PowerSeries.X * PowerSeries.derivative _ F} or as \lean{eulerDeriv F}
(Definition~\ref{ch19:def:twistedTD}); it multiplies the coefficient of $p^n$ by $n$.

\begin{remark}[Why logarithmic derivatives instead of logarithms]\label{ch19:rem:logfree}
Classically one writes the denominator identity as $\log(\text{product})=-\sum_{k}\frac1k\sum\cdots$.  Mathlib v4.30
has no formal logarithm or exponential on these rings, and in characteristic $0$ none is needed: for a series $A$ with
constant term $1$ (in $p$), the \emph{logarithmic derivative} $DA/A$ determines $A$ (Theorem~\ref{ch19:thm:ode}).
Every identity of the form ``$\log A=L$'' is therefore stated as ``$DA=A\cdot DL$''.  This is a reformulation, not a
weakening: two series with constant term $1$ and the same logarithmic derivative are equal.
\end{remark}

%======================================================================
\section{Faber polynomials}\label{ch19:sec:faber}
%======================================================================

Replication is a statement about polynomials in a $q$-series.  If $f=q^{-1}+\sum_{n\ge1}c(n)q^n$, then $f^m$ has a pole
of order $m$, and by subtracting suitable multiples of lower powers one can remove every non-positive power of $q$
except $q^{-m}$.  The result is the \emph{Faber polynomial} $F_m$ of $f$, the unique monic polynomial of degree $m$ with
$F_m(f)=q^{-m}+O(q)$.  For $f=J$ it computes the Hecke transform $J|T_m$ (Chapter~\ref{ch:16}); for a McKay--Thompson
series it is the polynomial of the replication formula.  The two files \leanpath{VOA/Moonshine/MoonshineLRFaber.lean}
and \leanpath{VOA/Moonshine/MoonshineLRFaberGenerating.lean} (namespace \lean{VOA.Moonshine.ConwayNorton.M8}) build them
for an arbitrary coefficient sequence $c:\N\to\C$; no convergence is involved.

\begin{definition}[Unit part, Laurent series, Faber polynomials]\label{ch19:def:faber}
For $c:\N\to\C$ let $u_c=1+\sum_{n\ge1}c(n)q^{n+1}\in\C[[q]]$ (no $q^1$ term) and $f_c=q^{-1}u_c=q^{-1}+\sum_{n\ge1}c(n)q^n\in
\C((q))$.  The Faber polynomials $F_m=F_{c,m}\in\C[X]$ are defined by strong recursion on $m$:
\begin{equation}\label{ch19:eq:faberrec}
  F_m=X^m-\sum_{j=0}^{m-1}\bigl([q^{-j}]f_c^{\,m}\bigr)\,F_j .
\end{equation}
\end{definition}
\begin{leancode}
noncomputable def qSeriesUnit (c : ℕ → ℂ) : PowerSeries ℂ :=
  PowerSeries.mk fun k => if k = 0 then 1 else if k = 1 then 0 else c (k - 1)
noncomputable def qSeriesLaurent (c : ℕ → ℂ) : LaurentSeries ℂ :=
  HahnSeries.single (-1 : ℤ) (1 : ℂ) * HahnSeries.ofPowerSeries ℤ ℂ (qSeriesUnit c)
noncomputable def faberPoly (c : ℕ → ℂ) (m : ℕ) : ℂ[X] :=
  X ^ m - ∑ j : Fin m, C ((qSeriesLaurent c ^ m).coeff (-((j : ℕ) : ℤ))) * faberPoly c j
termination_by m
decreasing_by exact j.isLt
\end{leancode}
(\leanpath{VOA/Moonshine/MoonshineLRFaber.lean}.)
\begin{dictionary}
\lean{PowerSeries.mk g} is the power series $\sum_kg(k)q^k$; so \lean{qSeriesUnit c} has coefficients $1,0,c(1),c(2),\dots$
in degrees $0,1,2,3,\dots$, i.e.\ it is $u_c$.  \lean{HahnSeries.single (-1) 1} is $q^{-1}$, and the product with the
embedded power series is $f_c$.  In \lean{faberPoly}: \lean{X} is the polynomial variable, \lean{C a} the constant $a$,
\lean{Fin m} the indices $0\le j<m$, and \lean{(qSeriesLaurent c ^ m).coeff (-j)} the coefficient of $q^{-j}$ in $f_c^m$.
\lean{termination_by}/\lean{decreasing_by} only certify that the recursion terminates.  So \lean{faberPoly c m} is
literally \eqref{ch19:eq:faberrec}.  As in \lean{qSeriesFun}, the value \lean{c 0} plays no role: the normalisation
``constant term $0$'' is built in.
\end{dictionary}
\begin{example}[The first Faber polynomials]\label{ch19:ex:faber}
$F_0=1$, $F_1=X$ (since $[q^0]f=0$), $F_2=X^2-2c(1)$ (since $f^2=q^{-2}+2c(1)+\cdots$), and
$F_3=X^3-3c(1)X-3c(2)$.  For $J$: $F_2(J)=J^2-393768=q^{-2}+42987520\,q+\cdots$, and $42987520=2\cdot21493760=2c(2)$, in
accordance with Theorem~\ref{ch19:thm:fabergen}(c) below.
\end{example}
\begin{physics}
For the moonshine CFT, $F_m(J)$ is the partition function of the theory ``wrapped $m$ times'': the sum over all
$m$-sheeted unbranched covers of the torus (equivalently, over sublattices of index $m$ of the period lattice).  Its only
polar term $q^{-m}$ is the vacuum of the $m$-fold cover.  This is the classical reading of the Hecke operator in
symmetric-product orbifolds (Dijkgraaf--Moore--Verlinde--Verlinde).
\end{physics}
\begin{status}
Definitions; printed with standard axioms in \leanpath{VOA/Moonshine/MoonshineLRFaberAudit.lean}.
\end{status}

\begin{theorem}[Faber polynomials: characterisation]\label{ch19:thm:faber}
For every $c$ and $m$:
\begin{enumerate}[label=(\alph*)]
\item $F_m$ is monic of degree $m$;
\item for every $k\le0$, $[q^k]F_m(f_c)=1$ if $k=-m$ and $0$ otherwise; that is, $F_m(f_c)-q^{-m}\in q\,\C[[q]]$;
\item if $P\in\C[X]$ satisfies the conclusion of (b), then $P=F_m$.
\end{enumerate}
\end{theorem}
\begin{leancode}
theorem faberPoly_monic (c : ℕ → ℂ) (m : ℕ) : (faberPoly c m).Monic
theorem faberPoly_natDegree (c : ℕ → ℂ) (m : ℕ) : (faberPoly c m).natDegree = m
theorem coeff_aeval_faberPoly (c : ℕ → ℂ) (m : ℕ) {k : ℤ} (hk : k ≤ 0) :
    (aeval (qSeriesLaurent c) (faberPoly c m)).coeff k = if k = -(m : ℤ) then 1 else 0
theorem eq_faberPoly_of_coeff (c : ℕ → ℂ) (m : ℕ) (P : ℂ[X])
    (h : ∀ k : ℤ, k ≤ 0 → (aeval (qSeriesLaurent c) P).coeff k = if k = -(m : ℤ) then 1 else 0) :
    P = faberPoly c m
\end{leancode}
(\leanpath{VOA/Moonshine/MoonshineLRFaber.lean}.)
\begin{dictionary}
\lean{aeval (qSeriesLaurent c) P} is $P(f_c)\in\C((q))$, and \lean{.coeff k} its coefficient of $q^k$
($k\in\Z$).  \lean{Monic} and \lean{natDegree} are Mathlib's.  The statement (b) is exactly the textbook normalisation
$F_m(f)=q^{-m}+0+O(q)$, with \emph{zero constant term} included ($k=0$, $m\ge1$).  Note that for $m=0$ the
statement says $F_0(f)$ has constant term $1$: $F_0=1$.  Uniqueness (c) is the reason why the library can later use an
\emph{existential} monic polynomial in \lean{CompletelyReplicable} and still be talking about the Faber polynomial.
\end{dictionary}
\paragraph{Proof on paper.}  (a) By induction: $X^m$ is monic of degree $m$ and every subtracted term has degree $<m$.
(b) A polynomial $P$ of degree $d$ evaluated at $f$ has $[q^{-d}]P(f)=$ the leading coefficient of $P$ and no lower
powers (\lean{coeff_aeval_qSeriesLaurent_of_le}), because $f^j=q^{-j}u^j$ with $u$ a unit power series.  By induction,
$F_j(f)=q^{-j}+(\text{positive powers})$ for $j<m$.  In $f^m-\sum_{j<m}[q^{-j}](f^m)F_j(f)$ the coefficient of $q^{-j}$
($0\le j<m$) is $[q^{-j}]f^m-[q^{-j}]f^m=0$, because only $F_j$ contributes to $q^{-j}$ (the $F_{j'}$ with $j'<j$ contribute
only $q^{-j'}$ and positive powers, those with $j'>j$ only $q^{-j'}$, which is $\ne q^{-j}$, and positive powers).  (c) If
$P$ and $F_m$ both satisfy (b), $P-F_m$ has all coefficients of $q^{k}$, $k\le0$, equal to zero; by the first remark a
nonzero polynomial has a nonzero coefficient at $q^{-\deg}$, $\deg\ge0$; so $P=F_m$. \qed
\begin{physics}
The uniqueness says: a quantity built polynomially from the partition function is determined by its polar part
(its ``ground-state'' content).  This is the formal shadow of the analytic statement of
Theorem~\ref{ch19:thm:vanish} below.
\end{physics}
\begin{status}
No hypotheses; standard axioms (\leanpath{VOA/Moonshine/MoonshineLRFaberAudit.lean}).
\end{status}

\begin{theorem}[Faber generating identity and recursion]\label{ch19:thm:fabergen}
Let $G(p)=p\bigl(f(p)-X\bigr)=1-Xp+\sum_{n\ge1}c(n)p^{n+1}\in\C[X][[p]]$ and $H(p)=\sum_{m\ge1}F_m(X)p^m$.  Then
\begin{enumerate}[label=(\alph*)]
\item $p\,G'(p)=-G(p)\,H(p)$, i.e.\ $\log G(p)=-\sum_{m\ge1}F_m(X)\,p^m/m$;
\item $F_{n+2}=X\,F_{n+1}-\sum_{j<n}c(j+1)F_{n-j}-(n+2)\,c(n+1)$ for all $n\ge0$;
\item $[q^1]F_m(f)=m\,c(m)$ for $m\ge1$.
\end{enumerate}
\end{theorem}
\begin{leancode}
noncomputable def faberGen (c : ℕ → ℂ) : PowerSeries ℂ[X] :=
  PowerSeries.mk fun k => if k = 0 then 1 else if k = 1 then -X else C (c (k - 1))
noncomputable def faberSeries (c : ℕ → ℂ) : PowerSeries ℂ[X] :=
  PowerSeries.mk fun m => if m = 0 then 0 else faberPoly c m
theorem faberGen_generating (c : ℕ → ℂ) :
    PowerSeries.X * PowerSeries.derivative ℂ[X] (faberGen c) = -(faberGen c * faberSeries c)
theorem faberPoly_recursion (c : ℕ → ℂ) (n : ℕ) :
    faberPoly c (n + 2) =
      X * faberPoly c (n + 1) - ∑ j ∈ Finset.range n, C (c (j + 1)) * faberPoly c (n - j) -
        C (((n : ℂ) + 2) * c (n + 1))
theorem coeff_one_aeval_faberPoly (c : ℕ → ℂ) {m : ℕ} (hm : 1 ≤ m) :
    (aeval (qSeriesLaurent c) (faberPoly c m)).coeff 1 = m * c m
\end{leancode}
(\leanpath{VOA/Moonshine/MoonshineLRFaberGenerating.lean}.)
\begin{dictionary}
Here the outer power-series variable \lean{PowerSeries.X} is $p$ and the polynomial variable \lean{X} (in \lean{-X})
is the indeterminate of $\C[X]$.  \lean{faberGen c} has $p$-coefficients $1,-X,c(1),c(2),\dots$: it is $G$.
\lean{faberSeries c} is $H$ (the $p^0$ term is $0$).  \lean{PowerSeries.X * PowerSeries.derivative ℂ[X] (faberGen c)}
is $p\,\partial_pG$.  So (a) is literally $pG'=-GH$, the logarithm-free form of $\log G=-\sum F_mp^m/m$
(Remark~\ref{ch19:rem:logfree}).  (b) is the coefficient of $p^{n+2}$ in (a).  (c) is stated with \lean{coeff 1}, the
coefficient of $q^1$.
\end{dictionary}
\paragraph{Proof on paper.}  Define polynomials $Q_m$ by $-pG'/G=\sum_{m\ge1}Q_mp^m$ ($G$ is invertible since its constant
term is $1$).  The coefficient of $X^m$ in $Q_m$ is $1$ (only the term $Xp$ of $G$ can produce $X^mp^m$, through
$-p\partial_p\log(1-Xp)=\sum_{k\ge1}X^kp^k$), and $\deg Q_m\le m$; so $Q_m$ is monic of degree $m$.  Now substitute
$X=f(q)$, a ring map $\C[X][[p]]\to\C((q))[[p]]$.  Here and below $f(p)$ denotes the same $q$-series written in the outer
variable, $f(p)=p^{-1}+\sum_{n\ge1}c(n)p^n$, so that $pf(p)=1+\sum c(n)p^{n+1}$ is a power series in $p$ and
$G(p)|_{X=f(q)}=p(f(p)-f(q))$.  Since $p\,p^{-1}=1$ and $p(p^n-q^n)=-(1-p/q)\,pq\sum_{i+j=n-1}p^iq^j$,
\begin{equation}\label{ch19:eq:factor}
  p\bigl(f(p)-f(q)\bigr)=\Bigl(1-\frac pq\Bigr)\,S(p,q),\qquad S=1-\sum_{a,b\ge1}c(a+b-1)\,p^aq^b\in1+pq\,\C[[p,q]] .
\end{equation}
Taking $-p\partial_p\log$ of both sides: $-p\partial_p\log(1-p/q)=\sum_{k\ge1}p^kq^{-k}$, and $-p\partial_pS/S\in pq\,\C[[p,q]]$.
Comparing coefficients of $p^m$: $Q_m(f)=q^{-m}+(\text{element of }q\C[[q]])$.  By Theorem~\ref{ch19:thm:faber}(c),
$Q_m=F_m$, which is (a).  (b) is the $p^{n+2}$-coefficient of $pG'=-GH$ rearranged.  (c): write $f'$ for the formal
derivative; the residue $[q^{-1}](P(f)f')$ vanishes for every polynomial $P$, because $f^jf'=(f^{j+1})'/(j+1)$ is a
derivative and derivatives have no $q^{-1}$ term (\lean{coeff_neg_one_pow_mul_deriv}).  Apply this to $P=F_m$: since
$F_m(f)=q^{-m}+\sum_{r\ge1}a_rq^r$ and $f'=-q^{-2}+\sum_{n}nc(n)q^{n-1}$, the residue is $-a_1+m\,c(m)$ (the
$q^{-m}\cdot mc(m)q^{m-1}$ term and the $a_1q\cdot(-q^{-2})$ term), so $a_1=mc(m)$. \qed
\begin{physics}
(a) has the form ``free energy $=\log$ of a partition function'': $\log\bigl(p(f(p)-X)\bigr)$ is a sum over connected
$m$-fold covers weighted by $1/m$, the standard relation between a second-quantised (multi-string) partition function and
the single-string one.  (c) is the first excited level of the $m$-fold cover: $m\,c(m)$ states.
\end{physics}
\begin{status}
No hypotheses; standard axioms (\leanpath{VOA/Moonshine/MoonshineLRFaberAudit.lean}, 35 reports).
\end{status}

\paragraph{Supporting declarations.}  \lean{constantCoeff_qSeriesUnit}, \lean{qSeriesLaurent_pow}
($f^j=q^{-j}u^j$), \lean{coeff_qSeriesLaurent_pow}, \lean{coeff_qSeriesLaurent_pow_of_le}, \lean{faberPoly_eq},
\lean{faberPoly_zero}, \lean{faberPoly_one}, \lean{faberPoly_degree_and_monic}, \lean{coeff_aeval_C_mul},
\lean{coeff_algebraMap_mul}, \lean{coeff_aeval_qSeriesLaurent_of_le}, \lean{eq_zero_of_coeff_aeval_qSeriesLaurent}
(file \leanpath{MoonshineLRFaber.lean}); \lean{coeff_mul_eq_zero_of_lt}, \lean{coeff_qSeriesUnit},
\lean{coeff_qSeriesLaurent_natCast}, \lean{coeff_qSeriesLaurent_of_lt_one}, the explicit formal derivative
\lean{qSeriesLaurentDeriv} ($f'=-q^{-2}u+q^{-1}u'$, written out because Mathlib v4.30 has no product rule for derivatives
of Laurent series), \lean{coeff_qSeriesLaurentDeriv_natCast}, \lean{coeff_qSeriesLaurentDeriv_sub_of_neg}, the residue
lemmas \lean{coeff_neg_one_pow_mul_deriv} and \lean{coeff_neg_one_aeval_mul_deriv}, and
\lean{coeff_aeval_faberPoly_sub_of_lt_one} (file \leanpath{MoonshineLRFaberGenerating.lean}).  All are coefficient
bookkeeping in $\C((q))$ and are printed in \leanpath{MoonshineLRFaberAudit.lean}.

%======================================================================
\section{\texorpdfstring{$J$}{J} is replicable}\label{ch19:sec:Jrep}
%======================================================================

The untwisted denominator identity is, in Borcherds' proof, a consequence of the Lie algebra structure.  The library
proves it the other way round: from the modularity of $J$, by showing that $J$ is \emph{replicable}, i.e.\ that the
Hecke transform of $J$ is the Faber polynomial of $J$.  The reason for this direction is recorded in
\S\ref{ch19:sec:denominator} (the generation of $\chnineteenm$ by its simple roots is \emph{equivalent} to the product
identity, so the identity must come from outside the Lie algebra).

\subsection{Hecke sums of a \texorpdfstring{$q$}{q}-series}

\begin{definition}[Hecke blocks, Hecke sums, Hecke coefficients]\label{ch19:def:qhecke}
For $c:\N\to\C$, $f=\mathtt{qSeriesFun}\ c$, $a,d\ge1$ and $M\ge1$:
\begin{equation*}
  B_{a,d}(\tau)=\sum_{0\le b<d}f\Bigl(\frac{a\tau+b}d\Bigr),\qquad R_M(\tau)=\sum_{a\mid M}B_{a,M/a}(\tau),\qquad
  h_c(M,r)=\sum_{a\mid\gcd(M,r)}\frac Ma\,c\Bigl(\frac{Mr}{a^2}\Bigr).
\end{equation*}
\end{definition}
\begin{leancode}
noncomputable def qHeckeBlock (c : ℕ → ℂ) (a d : ℕ) (τ : ℍ) : ℂ :=
  if h : 0 < a ∧ 0 < d then
    ∑ b ∈ Finset.range d, qSeriesFun c (heckePoint a d b h.1 h.2 τ)
  else 0
noncomputable def qHeckeSum (c : ℕ → ℂ) (M : ℕ) (τ : ℍ) : ℂ :=
  ∑ a ∈ M.divisors, qHeckeBlock c a (M / a) τ
noncomputable def qHeckeCoeff (c : ℕ → ℂ) (M r : ℕ) : ℂ :=
  ∑ a ∈ (Nat.gcd M r).divisors, ((M / a : ℕ) : ℂ) * c (M * r / (a * a))
\end{leancode}
(\leanpath{VOA/Moonshine/MoonshineLRM5Hecke.lean}.)
\begin{dictionary}
\lean{heckePoint a d b ha hd τ} is $(a\tau+b)/d\in\chnineteenH$ (Chapter~\ref{ch:16}).  The \lean{if h : 0 < a ∧ 0 < d}
only makes the definition total (the value $0$ is never used for $a,d\ge1$).  \lean{M.divisors} is the finite set of
positive divisors of $M$; \lean{M / a} is exact division for $a\mid M$.  \lean{(Nat.gcd M r).divisors} runs over
$a\mid\gcd(M,r)$, and \lean{M * r / (a * a)} is $Mr/a^2$, an integer for such $a$.  So $R_M$ is the sum over the
$\sigma_1(M)$ upper-triangular matrices $\chnineteenmat ab0d$ with $ad=M$, $0\le b<d$: the unnormalised Hecke operator
$f|T_M$ of Chapter~\ref{ch:16}.
\end{dictionary}

\begin{physics}
Hecke operators act on partition functions as sums over the $p$-fold covers of the torus. Physically, a Hecke sum at
a prime $p$ adds up the partition functions of the theory on all tori that cover the given one $p$ times. In
moonshine, the twisted sectors $g^a$ appear because winding $a$ times around a cycle twists by $g^a$. This is the
origin of the replication formulas.
\end{physics}

\begin{theorem}[The $q$-expansion of a Hecke sum]\label{ch19:thm:qheckesum}
If $c$ is $J$-dominated and $M\ge1$, then for every $\tau\in\chnineteenH$
\begin{equation*}
  R_M(\tau)-q^{-M}=\sum_{r\ge1}h_c(M,r)\,q^r ,
\end{equation*}
the series converging (in Mathlib's \lean{HasSum} sense) at every point.
\end{theorem}
\begin{leancode}
theorem hasSum_qHeckeBlock {c : ℕ → ℂ} (hc : IsJDominated c) (a d : ℕ) (ha : 0 < a) (hd : 0 < d)
    (τ : ℍ) :
    HasSum (fun r : ℕ => (if r ≠ 0 ∧ a ∣ r then (d : ℂ) * c (d * (r / a)) else 0) * qPar τ ^ r)
      (qHeckeBlock c a d τ - if d = 1 then (qPar τ ^ a)⁻¹ else 0)
theorem hasSum_qHeckeSum {c : ℕ → ℂ} (hc : IsJDominated c) {m : ℕ} (hm : 0 < m) (τ : ℍ) :
    HasSum (fun r : ℕ => (if r = 0 then 0 else qHeckeCoeff c m r) * qPar τ ^ r)
      (qHeckeSum c m τ - (qPar τ ^ m)⁻¹)
\end{leancode}
(\leanpath{VOA/Moonshine/MoonshineLRM5Hecke.lean}.)
\begin{dictionary}
\lean{HasSum (fun r => x r) s} says that the (unconditional) series $\sum_rx_r$ converges to $s$.  The block version
says $B_{a,d}(\tau)-[d=1]q^{-a}=\sum_{r}[r\ne0,\ a\mid r]\;d\,c(dr/a)\,q^r$.  The term \lean{if r = 0 then 0 else ...}
makes the $r=0$ coefficient zero: there is no constant term.  \lean{(qPar τ ^ m)⁻¹} is $q^{-M}$.
\end{dictionary}
\paragraph{Proof on paper.}  Write $f((a\tau+b)/d)=\sum_{n\ge-1}c(n)e^{2\pi in(a\tau+b)/d}$ (with $c(-1)=1$ and the $n=0$
term absent).  Summing over $b$, the character sum $\sum_{b<d}e^{2\pi inb/d}=d\,[d\mid n]$ keeps exactly the exponents
$n=ds$, each with weight $d$, and $e^{2\pi i\,ds\,a\tau/d}=q^{as}$.  The pole $n=-1$ survives only if $d=1$, giving
$q^{-a}$.  So $B_{a,d}=[d=1]q^{-a}+\sum_{s\ge1}d\,c(ds)\,q^{as}$, which is the block statement with $r=as$.  Summing over
$a\mid M$, $d=M/a$: the pole is $q^{-M}$ (only $a=M$ has $d=1$), and the coefficient of $q^r$ collects the $a\mid r$, i.e.\
$a\mid\gcd(M,r)$, with value $\frac Ma\,c\bigl(\frac Ma\cdot\frac ra\bigr)$.  Absolute convergence follows from $|c(n)|\le c_J(n)$.
\qed
\begin{physics}
Projecting onto $d$-periodic states: averaging over the $d$ shifts $\tau\mapsto\tau+b/d$ is the orbifold projection onto
states whose spin $L_0-\bar L_0$ is divisible by $d$; this is why only every $d$-th coefficient survives.
\end{physics}
\begin{status}
Hypothesis \lean{IsJDominated c} only; standard axioms (\leanpath{VOA/Moonshine/MoonshineLRM5Audit.lean}).
\end{status}

\begin{theorem}[Polynomials of $q$-series]\label{ch19:thm:polyq}
If $c$ is $J$-dominated, $P\in\C[X]$ has degree $D$ and $\tau\in\chnineteenH$, then
$q^D\,P\bigl(f(\tau)\bigr)=\sum_{n\ge0}[q^{n-D}]P(f_c)\;q^n$, where $P(f_c)\in\C((q))$ is the formal evaluation.
\end{theorem}
\begin{leancode}
theorem hasSum_mul_eval_qSeriesFun {c : ℕ → ℂ} (hc : IsJDominated c) (P : ℂ[X]) (τ : ℍ) :
    HasSum (fun n : ℕ => (aeval (qSeriesLaurent c) P).coeff ((n : ℤ) - P.natDegree) * qPar τ ^ n)
      (qPar τ ^ P.natDegree * P.eval (qSeriesFun c τ))
\end{leancode}
(\leanpath{VOA/Moonshine/MoonshineLRM5Poly.lean}.)
\begin{dictionary}
\lean{P.eval (qSeriesFun c τ)} is the number $P(f(\tau))$; \lean{aeval (qSeriesLaurent c) P} is the formal Laurent series
$P(f_c)$.  The theorem says that ``evaluate, then expand'' equals ``expand formally, then evaluate'': the analytic and
formal worlds agree on polynomials.  The factor $q^D$ clears the pole so that the right side is a power series.
\end{dictionary}
\paragraph{Proof on paper.}  $q f(\tau)=u_c(q)$ is an absolutely convergent power series in $q$ ($|q|<1$, coefficients
dominated by those of $J$), and $q^DP(f)=\sum_{j\le D}p_j\,q^{D-j}u_c^j$ where $P=\sum p_jX^j$; absolutely convergent power
series can be added and multiplied termwise (Cauchy product), and the same expression computes $q^DP(f_c)$ formally.
\qed
\begin{physics}
Products of partition functions (multi-copy theories) have the $q$-expansion obtained by multiplying the state counts:
the analytic and the combinatorial pictures agree as long as everything converges.
\end{physics}
\begin{status}
Hypothesis \lean{IsJDominated c}; standard axioms (\leanpath{MoonshineLRM5Audit.lean}).  Supporting declarations in the same
file: \lean{PowerSeriesAbsConv}, \lean{powerSeriesEval} and their closure under $+$, $\times$, powers, finite sums
(\lean{PowerSeriesAbsConv.add}, \lean{powerSeriesEval_mul}, \dots), \lean{powerSeriesAbsConv_qSeriesUnit},
\lean{powerSeriesEval_qSeriesUnit}, \lean{polyUnitSeries}, \lean{aeval_qSeriesLaurent_eq_single_mul},
\lean{powerSeriesEval_polyUnitSeries}.
\end{status}

\subsection{\texorpdfstring{$\mathrm{SL}_2(\mathbb Z)$}{SL2(Z)}-invariance of the Hecke sum}

\begin{theorem}[Invariance of $R_M$]\label{ch19:thm:heckeinv}
\begin{enumerate}[label=(\alph*)]
\item A function on $\chnineteenH$ invariant under $S=\chnineteenmat0{-1}10$ and $T=\chnineteenmat1101$ is
  $\chnineteenSL$-invariant.
\item For every $c$ and $M$, $R_M$ is invariant under $T$.
\item If $f=\mathtt{qSeriesFun}\ c$ is $\chnineteenSL$-invariant, then so is $R_M$ for every $M$.
\end{enumerate}
\end{theorem}
\begin{leancode}
theorem smul_invariant_of_S_T {R : ℍ → ℂ} (hS : ∀ τ : ℍ, R (ModularGroup.S • τ) = R τ)
    (hT : ∀ τ : ℍ, R (ModularGroup.T • τ) = R τ) (γ : SL(2, ℤ)) (τ : ℍ) : R (γ • τ) = R τ
theorem qHeckeSum_T_smul (c : ℕ → ℂ) (M : ℕ) (τ : ℍ) :
    qHeckeSum c M (ModularGroup.T • τ) = qHeckeSum c M τ
theorem qHeckeSum_smul {c : ℕ → ℂ}
    (hf : ∀ (γ : SL(2, ℤ)) (τ : ℍ), qSeriesFun c (γ • τ) = qSeriesFun c τ) (M : ℕ)
    (γ : SL(2, ℤ)) (τ : ℍ) : qHeckeSum c M (γ • τ) = qHeckeSum c M τ
\end{leancode}
(\leanpath{VOA/Moonshine/MoonshineLRM5Invariance.lean}, \leanpath{VOA/Moonshine/MoonshineLRM5CosetMap.lean}.)
\begin{dictionary}
\lean{ModularGroup.S}, \lean{ModularGroup.T} are Mathlib's generators of \lean{SL(2, ℤ)}, acting on \lean{ℍ} by
\lean{γ • τ}.  The hypothesis of (c) is \emph{invariance of the function}, not of the coefficient sequence; it holds for
$c=\mathtt{cJ}$ (\lean{qSeriesFun_cJ_smul}, from $\mathtt{qSeriesFun}\ \mathtt{cJ}=j-744$, \lean{qSeriesFun_cJ}).
\end{dictionary}
\paragraph{Proof on paper.}  (a) $S$ and $T$ generate $\chnineteenSL$ (Mathlib's \lean{SL2Z_generators}).  (b) In the block
$(a,d)$, $\tau\mapsto\tau+1$ changes $(a\tau+b)/d$ into $(a\tau+(b+a))/d$; since $f$ is $1$-periodic only $b+a$ modulo $d$
matters, and $b\mapsto b+a\pmod d$ is a bijection of $\{0,\dots,d-1\}$.  (c) The classical argument: for $A$ in the coset
set $\Delta(M)=\{\chnineteenmat ab0d:ad=M,0\le b<d\}$ and $\gamma\in\chnineteenSL$ there are unique $\gamma'\in\chnineteenSL$ and
$A'\in\Delta(M)$ with $A\gamma=\gamma'A'$, and $A\mapsto A'$ is a bijection of $\Delta(M)$; then
$R_M(\gamma\tau)=\sum_Af(\gamma'A'\tau)=\sum_{A'}f(A'\tau)=R_M(\tau)$.  By (a) and (b) it suffices to do this for $\gamma=S$, where
the library gives the bijection explicitly: for $A=\chnineteenmat ab0d$ put $g=\gcd(b,d)$, $d=gd_1$, $b=gb_1$, and let $r\in[0,d_1)$
solve $rb_1\equiv-1\pmod{d_1}$; then $A'=\chnineteenmat g{ar}0{ad_1}$, the map $\langle a,b\rangle\mapsto\langle g,ar\rangle$ is an
involution of $\Delta(M)$, and $(a\,S\tau+b)/d=\gamma'\bigl((g\tau+ar)/(ad_1)\bigr)$ with
$\gamma'=\chnineteenmat{b_1}{*}{d_1}{*}\in\chnineteenSL$ (the ``point identity'' \lean{heckePoint_S_smul_eq}).  For example, for
$M=2$ the involution fixes $\langle1,1\rangle$ and exchanges $\langle1,0\rangle$ and $\langle2,0\rangle$: indeed
$J(-1/(2\tau))=J(2\tau)$ and $J\bigl((\tau-1)/(2\tau)\bigr)=J\bigl((\tau+1)/2\bigr)$, the matrix $\chnineteenmat1{-1}2{-1}$ relating
$(\tau+1)/2$ to $(\tau-1)/(2\tau)$. \qed
\begin{physics}
The sum over all $m$-sheeted covers of a torus is again modular invariant: a modular transformation of the base permutes
the covers.  This is the modular invariance of the symmetric-product (second-quantised) partition function.
\end{physics}
\begin{status}
(a), (b) no hypotheses; (c) hypothesis: invariance of $f$.  Standard axioms (\leanpath{MoonshineLRM5Audit.lean}).
Supporting declarations: \lean{qHeckeBlock_vadd_one} (file \leanpath{MoonshineLRM5Invariance.lean}); \lean{cosetMatrix},
\lean{heckePoint_S_smul_eq} (file \leanpath{MoonshineLRM5Cosets.lean}); \lean{negInvMod} and its lemmas
\lean{negInvMod_cast}, \lean{negInvMod_lt}, \lean{dvd_negInvMod_mul_add_one}, \lean{eq_negInvMod},
\lean{coprime_of_dvd_mul_add_one}; \lean{heckeCosets}, \lean{cosetS}, the arithmetic record \lean{CosetData} and its
eight projections, \lean{cosetData_of_mem}, \lean{cosetS_mem}, \lean{cosetS_cosetS} (involution), \lean{heckeVal},
\lean{heckeSumOf}, \lean{qHeckeSum_eq_heckeSumOf}, \lean{heckeVal_S_smul}, \lean{heckeSumOf_S_smul} (file
\leanpath{MoonshineLRM5CosetMap.lean}).
\end{status}

\subsection{A modular invariant without polar part vanishes}

\begin{theorem}[Vanishing of invariant cusp forms of weight $0$]\label{ch19:thm:vanish}
Let $D:\chnineteenH\to\C$ be $\chnineteenSL$-invariant and $D(\tau)=\sum_{r\ge0}\delta_rq^r$ for every $\tau$, with
$\delta_0=0$.  Then $\delta_r=0$ for all $r$ (hence $D=0$).
\end{theorem}
\begin{leancode}
theorem eq_zero_of_invariant_hasSum {δ : ℕ → ℂ} {D : ℍ → ℂ}
    (hD : ∀ (γ : SL(2, ℤ)) (τ : ℍ), D (γ • τ) = D τ)
    (hs : ∀ τ : ℍ, HasSum (fun r => δ r * qPar τ ^ r) (D τ)) (h0 : δ 0 = 0) (r : ℕ) :
    δ r = 0
\end{leancode}
(\leanpath{VOA/Moonshine/MoonshineLRM5Replicable.lean}.)
\begin{dictionary}
No holomorphy is assumed: the hypotheses are invariance and convergence of the $q$-expansion at every point.  (A
convergent power series in $q$ is automatically holomorphic, so this is the classical statement that a holomorphic
modular function of weight $0$ for $\chnineteenSL$ that vanishes at the cusp is identically zero.)
\end{dictionary}
\paragraph{Proof on paper.}  The classical proof uses the maximum principle or Liouville on $X(1)$.  The library gives an
elementary proof.  Every $\tau$ is $\chnineteenSL$-equivalent to a point of height $\ge\frac12$
(\lean{exists_smul_half_le_im}), where $|q|\le\rho=e^{-\pi}<\frac12$.  First, $|D|$ is bounded on $\chnineteenH$ by
$B_0=\sum|\delta_r|\rho^r$ (absolute convergence at $\tau_0=i/2$ and invariance).  Second, a bounded function with
$q$-expansion $\sum A_rq^r$ has $|A_r|\le\sup|D|$ (the coefficient bound, Theorem~\ref{ch19:thm:coeffbound} below).  Third,
if $|\delta_r|\le B$ for all $r$, then on height $\ge\frac12$ we get $|D|\le\sum_{r\ge1}B\rho^r=B\rho/(1-\rho)$, hence
$|D|\le B\kappa$ everywhere with $\kappa=\rho/(1-\rho)<1$, hence $|\delta_r|\le B\kappa$ (\lean{norm_le_mul_of_invariant_hasSum}).
Iterating, $|\delta_r|\le B_0\kappa^n\to0$. \qed
\begin{physics}
``The polar part determines the partition function'': a modular-invariant partition function with no states of negative
energy and no vacuum degeneracy at the cusp must vanish.  This is the principle behind the Farey-tail and Rademacher
expansions of partition functions (Dijkgraaf--Maldacena--Moore--Verlinde; Duncan--Frenkel).
\end{physics}
\begin{status}
No hypotheses beyond those displayed; standard axioms (\leanpath{MoonshineLRM5Audit.lean}).  Supporting:
\lean{exp_neg_pi_lt_half}, \lean{exists_smul_half_le_im}, \lean{norm_le_mul_of_invariant_hasSum}.
\end{status}

\subsection{The replicability of \texorpdfstring{$J$}{J}}

\begin{theorem}[$J$ is replicable]\label{ch19:thm:Jrep}
Let $P_M=F_{\mathtt{cJ},M}$ be the Faber polynomial of $J$.  For every $M\ge1$:
\begin{enumerate}[label=(\alph*)]
\item $P_M(J(\tau))=\sum_{ad=M}\sum_{0\le b<d}J\bigl((a\tau+b)/d\bigr)$ for every $\tau\in\chnineteenH$;
\item $[q^r]P_M(J)=\sum_{a\mid\gcd(M,r)}\frac Ma\,c\bigl(\frac{Mr}{a^2}\bigr)$ for every $r\ge1$ (``(Rep$_1$)'').
\end{enumerate}
\end{theorem}
\begin{leancode}
theorem qSeriesFun_cJ (τ : ℍ) : qSeriesFun cJ τ = jInvariant τ - 744
theorem faberPoly_cJ_eval_eq_qHeckeSum {M : ℕ} (hM : 0 < M) (τ : ℍ) :
    (faberPoly cJ M).eval (qSeriesFun cJ τ) = qHeckeSum cJ M τ
theorem coeff_faberPoly_cJ_eq_qHeckeCoeff {M : ℕ} (hM : 0 < M) {r : ℕ} (hr : r ≠ 0) :
    (aeval (qSeriesLaurent cJ) (faberPoly cJ M)).coeff (r : ℤ) = qHeckeCoeff cJ M r
\end{leancode}
(\leanpath{VOA/Moonshine/MoonshineLRM5Replicable.lean}.)
\begin{dictionary}
\lean{jInvariant} is the analytic $j$-function of Chapter~\ref{ch:L14}; \lean{qSeriesFun_cJ} identifies the library's
$q$-series encoding of $J$ with $j-744$ on all of $\chnineteenH$ (from \lean{hasSum_JShift_laurent}, Chapter~\ref{ch:05}).  So
(a) is the identity of functions $P_M\circ J=J|T_M$, and (b) the identity of coefficients.  (b) is stated for the
\emph{formal} series $P_M(J)\in\C((q))$; (a) for the function.  Both have no hypotheses.
\end{dictionary}
\begin{example}[$M=2$]\label{ch19:ex:J2}
(b) at $M=2$, $r=2$: $[q^2](J^2-2c(1))=2c(4)+c(1)$.  The left side is $2c(3)+c(1)^2$.  Hence $c(4)=c(3)+\frac12(c(1)^2-c(1))$,
Mahler's recursion of Chapter~\ref{ch:16}: $864299970+\frac12(196884^2-196884)=20245856256$.
\end{example}
\paragraph{Proof on paper.}  Let $D=P_M\circ J-R_M$.  By Theorem~\ref{ch19:thm:faber}(b) and Theorem~\ref{ch19:thm:polyq},
$P_M(J(\tau))=q^{-M}+\sum_{r\ge1}[q^r]P_M(J)\,q^r$ (\lean{hasSum_faberPoly_eval_sub}); by Theorem~\ref{ch19:thm:qheckesum},
$R_M=q^{-M}+\sum_{r\ge1}h(M,r)q^r$.  So $D=\sum_{r\ge1}\delta_rq^r$ with $\delta_r=[q^r]P_M(J)-h(M,r)$
(\lean{replicationDefect}, \lean{hasSum_replicationDefect}).  $D$ is $\chnineteenSL$-invariant, since $J$ is and $R_M$ is
(Theorem~\ref{ch19:thm:heckeinv}).  By Theorem~\ref{ch19:thm:vanish}, all $\delta_r=0$: this is (b), and summing, (a). \qed
\begin{physics}
The partition function of the $M$-fold cover is a polynomial in the single-cover partition function.  For a holomorphic
CFT with $c=24$ this is forced by modular invariance: the $M$-fold cover theory has a modular-invariant partition function
whose only polar term is $q^{-M}$, and such a function is a polynomial in $J$.
\end{physics}
\begin{status}
\textbf{Unconditional} (no hypotheses); standard axioms (\leanpath{MoonshineLRM5Audit.lean}).  Supporting declarations:
\lean{isJDominated_cJ}, \lean{qSeriesFun_cJ_smul}, \lean{hasSum_faberPoly_eval_sub}, \lean{replicationDefect},
\lean{hasSum_replicationDefect}.
\end{status}

%======================================================================
\section{The product identity for \texorpdfstring{$J$}{J} (the denominator identity)}\label{ch19:sec:denominator}
%======================================================================

\subsection{The identity, and what it means}

\emph{Motivation only.}  The following two paragraphs explain where the identity comes from in Borcherds' proof; nothing in
them is used below, and no Lean theorem of this chapter refers to $\chnineteenm$ in connection with the identity (the link with
$\chnineteenm$ is made in Chapter~\ref{ch:20}).  Let $\chnineteenm$ be the Monster Lie algebra (Chapter~\ref{ch:18}): graded by $I\!I_{1,1}$, with root spaces
$\chnineteenm_{(m,n)}\cong\Vnat_{1+mn}$ in Borcherds' coordinates, so that $\dim\chnineteenm_{(m,n)}=c(mn)$ for $(m,n)\ne0$.  Its
simple roots are the real root $(1,-1)$ and the imaginary roots $(1,n)$, $n\ge1$, with multiplicity $c(n)$, and its Weyl
vector is $\rho=(-1,0)$ in Borcherds' sign convention (it pairs with every simple root $\alpha$ to $-\frac12(\alpha,\alpha)$; Chapter~\ref{ch:18}
uses the opposite sign).  The Weyl--Kac--Borcherds denominator formula of a generalised Kac--Moody algebra, in which $\epsilon(\mu)$ is
$(-1)^n$ if $\mu$ is a sum of $n$ pairwise orthogonal imaginary simple roots and $0$ otherwise (we do not derive it here),
\begin{equation*}
  e^{\rho}\prod_{\alpha>0}(1-e^{\alpha})^{\operatorname{mult}\alpha}=\sum_{w\in W}\det(w)\,w\Bigl(e^\rho\sum_{\mu}\epsilon(\mu)e^{\mu}\Bigr),
\end{equation*}
becomes for $\chnineteenm$, with $p=e^{(1,0)}$ and $q=e^{(0,1)}$, Borcherds' identity
\begin{equation}\label{ch19:eq:borcherds}
  p^{-1}\prod_{m>0,\;n\in\Z}\bigl(1-p^mq^n\bigr)^{c(mn)}=J(p)-J(q) .
\end{equation}
(The Weyl group has order $2$, generated by the reflection $(m,n)\mapsto(n,m)$ in the real root, which is why the right
side is antisymmetric in $p,q$.)  In \eqref{ch19:eq:borcherds} the factors with $n\le0$ are $(1-pq^{-1})^{c(-1)}=1-p/q$
(from $mn=-1$) and $(1-p^m)^{c(0)}=1$; all others with $n\le0$ have $mn\le-2$ and exponent $0$.  Using the factorisation
\eqref{ch19:eq:factor}, $p(J(p)-J(q))=(1-p/q)S$ with $S=1-\sum_{a,b\ge1}c(a+b-1)p^aq^b$, the identity
\eqref{ch19:eq:borcherds} is therefore equivalent to the \emph{Koike--Norton--Zagier} (KNZ) form
\begin{equation}\label{ch19:eq:knz}
  \prod_{m,n\ge1}\bigl(1-p^mq^n\bigr)^{c(mn)}=1-\sum_{a,b\ge1}c(a+b-1)\,p^aq^b .\tag{KNZ}
\end{equation}
This identity about $J$ (with $\prod$ read as \lean{coneProd}) is what the library proves.  Its route is as follows:
\begin{equation*}
  \text{$j$ is $\chnineteenSL$-invariant}\ \Longrightarrow\ \text{$J$ is replicable (Thm.~\ref{ch19:thm:Jrep})}
  \ \Longrightarrow\ D\log S=D\log\prod\ \Longrightarrow\ \eqref{ch19:eq:knz}.
\end{equation*}

\begin{warning}
\textbf{The product identity cannot be derived from the Lie algebra alone} (a recorded dead end).  Given the no-ghost theorem and the invariant form, the statement ``$\chnineteenm$ is generated by its simple
root spaces'' is \emph{equivalent} to \eqref{ch19:eq:knz}: the Casimir argument (Chapter~\ref{ch:18}) shows that the
subalgebra generated by the known simple roots has character $1/S$, and it is all of $\chnineteenm$ exactly when the product
of the multiplicities $c(mn)$ is $S$.  The Casimir argument cannot exclude further imaginary simple roots $(m,n)$ with
$m\ge2$, because $\rho$ is not known to be a Weyl vector for them.  So any proof of generation \emph{is} a proof of the
product formula, i.e.\ of a modularity property of $J$.  The project therefore proves \eqref{ch19:eq:knz} analytically
(this section) and deduces generation from it (Chapter~\ref{ch:20}, \lean{M8.monsterLieSimpleRootGeneration}, which uses
\lean{knz_product} and \lean{coneProd_exponent_unique}).  Attempting a Lie-theoretic proof of generation is a recorded
dead end.
\end{warning}

\subsection{The logarithmic form}

\begin{definition}[The series $S$ and the logarithmic series]\label{ch19:def:knz}
For $c:\N\to\C$ define in $\C((q))[[p]]$
\begin{equation*}
  S_c=1-\sum_{a,b\ge1}c(a+b-1)\,p^aq^b,\qquad E_c=\sum_{M\ge1}\bigl(F_M(f_c)-q^{-M}\bigr)p^M .
\end{equation*}
\end{definition}
\begin{leancode}
noncomputable def knzCoeff (c : ℕ → ℂ) (a : ℕ) : LaurentSeries ℂ :=
  if a = 0 then 1 else
    -HahnSeries.ofPowerSeries ℤ ℂ (PowerSeries.mk fun b => if b = 0 then 0 else c (a + b - 1))
noncomputable def knzSeries (c : ℕ → ℂ) : PowerSeries (LaurentSeries ℂ) :=
  PowerSeries.mk (knzCoeff c)
noncomputable abbrev qInv : LaurentSeries ℂ := HahnSeries.single (-1 : ℤ) (1 : ℂ)
noncomputable def knzLogCoeff (c : ℕ → ℂ) (M : ℕ) : LaurentSeries ℂ :=
  if M = 0 then 0 else aeval (qSeriesLaurent c) (faberPoly c M) - qInv ^ M
noncomputable def knzLogSeries (c : ℕ → ℂ) : PowerSeries (LaurentSeries ℂ) :=
  PowerSeries.mk (knzLogCoeff c)
noncomputable def evalAtQSeries (c : ℕ → ℂ) : ℂ[X] →+* LaurentSeries ℂ :=
  (aeval (qSeriesLaurent c)).toRingHom
\end{leancode}
(\leanpath{VOA/Moonshine/MoonshineLRM5KNZ.lean}.)
\begin{dictionary}
\lean{knzCoeff c a} is the $p^a$-coefficient of $S_c$: $1$ for $a=0$ and $-\sum_{b\ge1}c(a+b-1)q^b$ for $a\ge1$.
\lean{qInv} is $q^{-1}$.  \lean{knzLogCoeff c M} is $F_M(f_c)-q^{-M}$ for $M\ge1$ and $0$ for $M=0$.  \lean{evalAtQSeries c}
is the ring map $\C[X]\to\C((q))$, $X\mapsto f_c$; mapping \lean{faberGen c} along it substitutes $X=f(q)$ in
$G(p)=p(f(p)-X)$, giving $p(f(p)-f(q))$.
\end{dictionary}
\begin{physics}
$S$ is the ``single-particle'' side: $1-S$ is the generating function of the positive-cone simple roots and their
real-root strings, i.e.\ of the generating space $U$ of Chapter~\ref{ch:18}.
\end{physics}
\begin{status}Definitions; standard axioms (\leanpath{MoonshineLRM5Audit.lean}).\end{status}

\begin{theorem}[Factorisation and logarithmic form]\label{ch19:thm:knzlog}
For every $c$:
\begin{enumerate}[label=(\alph*)]
\item $p\bigl(f_c(p)-f_c(q)\bigr)=(1-q^{-1}p)\,S_c$ in $\C((q))[[p]]$;
\item $p\,\partial_pS_c=-S_c\,E_c$.
\end{enumerate}
For $c=\mathtt{cJ}$:
\begin{enumerate}[label=(\alph*),resume]
\item $p\,\partial_pS=-S\cdot\sum_{M\ge1}\Bigl(\sum_{r\ge1}h(M,r)\,q^r\Bigr)p^M$, with $h(M,r)=\sum_{a\mid(M,r)}\frac Ma\,c(\frac{Mr}{a^2})$;
  equivalently $p\,\partial_p\log S=-\sum_{k,m,n\ge1}m\,c(mn)\,p^{mk}q^{nk}$.
\end{enumerate}
\end{theorem}
\begin{leancode}
theorem faberGen_map_eq_mul (c : ℕ → ℂ) :
    PowerSeries.map (evalAtQSeries c) (faberGen c) =
      (1 - PowerSeries.C qInv * PowerSeries.X) * knzSeries c
theorem knz_log (c : ℕ → ℂ) :
    PowerSeries.X * PowerSeries.derivative (LaurentSeries ℂ) (knzSeries c) =
      -(knzSeries c * knzLogSeries c)
theorem knz_log_J :
    PowerSeries.X * PowerSeries.derivative (LaurentSeries ℂ) (knzSeries cJ) =
      -(knzSeries cJ * PowerSeries.mk fun M => if M = 0 then 0 else
        HahnSeries.ofPowerSeries ℤ ℂ
          (PowerSeries.mk fun r => if r = 0 then 0 else qHeckeCoeff cJ M r))
\end{leancode}
(\leanpath{VOA/Moonshine/MoonshineLRM5KNZ.lean}.)
\begin{dictionary}
\lean{PowerSeries.map (evalAtQSeries c) (faberGen c)} is $1-f_c(q)\,p+\sum_{n\ge1}c(n)p^{n+1}=p(f_c(p)-f_c(q))$: here
$f_c(p)$ means the series $p^{-1}+\sum c(n)p^n$ in the outer variable, and $pf_c(p)$ is a power series in $p$.  In
(a), \lean{1 - PowerSeries.C qInv * PowerSeries.X} is $1-q^{-1}p$.  In (c), the inner \lean{PowerSeries.mk fun r => ...}
is $\sum_{r\ge1}h(M,r)q^r$ embedded into $\C((q))$, and \lean{qHeckeCoeff cJ M r} is $h(M,r)$.  The equivalence with
the form $-\sum_{k,m,n}m\,c(mn)p^{mk}q^{nk}$ stated in the docstring is the identity
$\sum_{a\mid(M,r)}\frac Ma\,c(\frac{Mr}{a^2})=[p^Mq^r]\sum_{k,m,n}m\,c(mn)p^{mk}q^{nk}$ (put $k=a$, $m=M/a$, $n=r/a$); it is
used in Lean as \lean{coneLog_cJ} below.
\end{dictionary}
\paragraph{Proof on paper.}  (a) is the telescoping identity behind \eqref{ch19:eq:factor}:
$p(p^{-1}-q^{-1})=1-p/q$ and $p(p^n-q^n)=-(1-p/q)\sum_{i+j=n-1}p^{i+1}q^{j+1}$.  (b): apply the ring map $X\mapsto f_c$ to the
Faber generating identity $pG'=-GH$ (Theorem~\ref{ch19:thm:fabergen}(a); the derivative commutes with coefficient maps,
\lean{derivative_map}), which gives $p\Phi'=-\Phi H(f)$ for $\Phi=(1-q^{-1}p)S$ (\lean{faberGen_map_generating}).  Since
$p\partial_p(1-q^{-1}p)=-(1-q^{-1}p)\sum_{k\ge1}q^{-k}p^k$ and $1-q^{-1}p$ is invertible with inverse
$\sum_kq^{-k}p^k$ (\lean{qInvGeom}, \lean{mul_qInvGeom}), dividing by $1-q^{-1}p$ yields
$pS'=-S\,(H(f)-\sum_kq^{-k}p^k)=-S\,E$ (\lean{knzLogSeries_eq}).  (c): for $J$, the coefficient $F_M(J)-q^{-M}$ of $E$ is
$\sum_{r\ge1}h(M,r)q^r$ by (Rep$_1$) (Theorem~\ref{ch19:thm:Jrep}(b); \lean{knzLogCoeff_cJ}). \qed
\begin{physics}
(c) is the denominator identity at the level of ``free energies'': the logarithmic derivative of the multi-particle
partition function equals the sum over connected contributions $m\,c(mn)$ of single physical states at momentum
$(m,n)$, summed over their $k$-fold multiples (the Bose--Einstein sum $\sum_kx^k$).
\end{physics}
\begin{status}
(a), (b): no hypotheses; (c): \textbf{unconditional}.  Standard axioms (\leanpath{MoonshineLRM5Audit.lean}).  Supporting:
\lean{derivative_map}, \lean{faberGen_map_generating}, \lean{coeff_ofPowerSeries_mk}, \lean{qSeriesLaurent_eq_qInv_add},
\lean{qInv_mul_ofPowerSeries_mk} (the telescoping step), \lean{qInvGeom}, \lean{mul_qInvGeom}, \lean{knzLogSeries_eq},
\lean{knzLogCoeff_cJ}.
\end{status}

\subsection{Products defined by their logarithmic derivative}

Mathlib v4.30 has no infinite products with complex exponents in $\C((q))[[p]]$.  The library defines the product
$\prod_{m,n\ge1}(1-p^mq^n)^{e(m,n)}$ as the unique series with constant term $1$ whose logarithmic derivative is the
expected one.

\begin{definition}[Solution of $pA'=AL$]\label{ch19:def:ode}
Let $K$ be a field of characteristic $0$ and $L\in K[[p]]$ with $L(0)=0$.  $\mathrm{odeSol}(L)=\sum_aA_ap^a$ is defined by
$A_0=1$ and $(a+1)A_{a+1}=\sum_{i\le a}A_i\,L_{a+1-i}$.
\end{definition}
\begin{leancode}
noncomputable def odeCoeff (L : PowerSeries K) : ℕ → K
  | 0 => 1
  | a + 1 => ((a + 1 : ℕ) : K)⁻¹ *
      ∑ i : Fin (a + 1), odeCoeff L i * PowerSeries.coeff (a + 1 - i) L
termination_by a => a
decreasing_by exact i.isLt
noncomputable def odeSol (L : PowerSeries K) : PowerSeries K := PowerSeries.mk (odeCoeff L)
\end{leancode}
(\leanpath{VOA/Moonshine/MoonshineLRM5Ode.lean}; section variables \lean{{K : Type*} [Field K] [CharZero K]}.)

\begin{physics}
Solving $pA'=AL$ is how one exponentiates a logarithmic derivative: $L$ plays the role of a sum over single-particle
contributions, and $A$ is the multi-particle (second-quantised) partition function built from them. This is the
familiar passage from a single-string spectrum to the free-field product formula of a string partition function.
\end{physics}

\begin{theorem}[Existence, uniqueness, multiplicativity]\label{ch19:thm:ode}
Let $L,L_1,L_2\in K[[p]]$ have constant term $0$.
\begin{enumerate}[label=(\alph*)]
\item $A=\mathrm{odeSol}(L)$ satisfies $pA'=AL$ and $A(0)=1$;
\item two solutions of $pA'=AL$ with the same constant term are equal;
\item $\mathrm{odeSol}(L_1+L_2)=\mathrm{odeSol}(L_1)\,\mathrm{odeSol}(L_2)$.
\end{enumerate}
\end{theorem}
\begin{leancode}
theorem odeSol_spec (L : PowerSeries K) (hL : PowerSeries.constantCoeff L = 0) :
    PowerSeries.X * PowerSeries.derivative K (odeSol L) = odeSol L * L
theorem eq_of_ode {A B L : PowerSeries K} (hL : PowerSeries.constantCoeff L = 0)
    (hA : PowerSeries.X * PowerSeries.derivative K A = A * L)
    (hB : PowerSeries.X * PowerSeries.derivative K B = B * L)
    (h0 : PowerSeries.constantCoeff A = PowerSeries.constantCoeff B) : A = B
theorem odeSol_add (L₁ L₂ : PowerSeries K) (h₁ : PowerSeries.constantCoeff L₁ = 0)
    (h₂ : PowerSeries.constantCoeff L₂ = 0) :
    odeSol (L₁ + L₂) = odeSol L₁ * odeSol L₂
\end{leancode}
(\leanpath{VOA/Moonshine/MoonshineLRM5Ode.lean}.)
\begin{dictionary}
\lean{PowerSeries.constantCoeff} is the constant term.  The theorem is the formal version of
``$A=\exp\int L\,dp/p$'': (a) existence, (b) uniqueness (in particular uniqueness given the constant term; no
invertibility hypothesis on $A$ is needed), (c) the exponential law.  It is applied below with $K=\C((q))$, which is a
field of characteristic $0$ (\lean{charZero_laurentSeries}).
\end{dictionary}
\paragraph{Proof on paper.}  (a) Comparing coefficients of $p^{a+1}$ in $pA'=AL$ gives exactly the recursion, since
$[p^{a+1}](pA')=(a+1)A_{a+1}$ and $L_0=0$.  (b) The difference $\Delta=A-B$ satisfies $p\Delta'=\Delta L$ and $\Delta_0=0$; by
induction $(a+1)\Delta_{a+1}=\sum_{i\le a}\Delta_iL_{a+1-i}=0$, and $a+1\ne0$ in characteristic $0$.  (c) $D=p\partial_p$ is a
derivation, so $A_1A_2$ solves $p(A_1A_2)'=A_1A_2(L_1+L_2)$ with constant term $1$; apply (b). \qed
\begin{physics}
The exponential of a connected generating function: a multi-particle partition function is the exponential of the
single-particle free energy, and the ODE is its infinitesimal form.
\end{physics}
\begin{status}No hypotheses beyond those displayed; standard axioms (\leanpath{MoonshineLRM5Audit.lean}).  Supporting:
\lean{odeCoeff_zero}, \lean{odeCoeff_succ}, \lean{constantCoeff_odeSol}, \lean{coeff_succ_mul_of_constantCoeff_eq_zero},
\lean{coeff_succ_X_mul_derivative}.\end{status}

\begin{definition}[Cone products]\label{ch19:def:coneprod}
For $e:\N\times\N\to\C$ let
\begin{equation*}
  \mathrm{coneLog}(e)=-\sum_{k,m,n\ge1}m\,e(m,n)\,p^{mk}q^{nk}\in\C((q))[[p]],\qquad
  \mathrm{coneProd}(e)=\mathrm{odeSol}\bigl(\mathrm{coneLog}(e)\bigr),
\end{equation*}
so that $\mathrm{coneProd}(e)$ ``is'' $\prod_{m,n\ge1}(1-p^mq^n)^{e(m,n)}$.
\end{definition}
\begin{leancode}
noncomputable def coneLogCoeff (e : ℕ → ℕ → ℂ) (a : ℕ) : LaurentSeries ℂ :=
  -HahnSeries.ofPowerSeries ℤ ℂ (PowerSeries.mk fun b =>
    ∑ m ∈ a.divisors, if a / m ∣ b ∧ b ≠ 0 then (m : ℂ) * e m (b / (a / m)) else 0)
noncomputable def coneLog (e : ℕ → ℕ → ℂ) : PowerSeries (LaurentSeries ℂ) :=
  PowerSeries.mk (coneLogCoeff e)
noncomputable def coneProd (e : ℕ → ℕ → ℂ) : PowerSeries (LaurentSeries ℂ) :=
  odeSol (coneLog e)
\end{leancode}
(\leanpath{VOA/Moonshine/MoonshineLRM5Product.lean}.)
\begin{dictionary}
\lean{coneLogCoeff e a} is the $p^a$-coefficient of $\mathrm{coneLog}(e)$: $-\sum_{m\mid a}\sum_{b\ge1,\;(a/m)\mid b}m\,e\bigl(m,b/(a/m)\bigr)q^b$.
With $k=a/m$ and $b=nk$ this is $-\sum_{mk=a}\sum_nm\,e(m,n)q^{nk}$, as displayed.  Why this is the logarithmic derivative of
the product: for $x=p^mq^n$, $p\partial_px=mx$, so $p\partial_p\log(1-x)^{e}=-e\,m\,x/(1-x)=-e\,m\sum_{k\ge1}x^k$.

\textbf{Is \lean{coneProd} the product?}  The library does \emph{not} prove that $\mathrm{coneProd}(e)$ equals the
$(p,q)$-adic limit of finite products $\prod_{m,n\le N}(1-p^mq^n)^{e(m,n)}$ (with $(1-x)^e$ the binomial series); no such
limit is defined in Lean, and the lemma \lean{coneProd_single} ($=(1-p^mq^n)^k$ for a single exponent) proposed in the
coordination log was not written.  What the library proves is that $\mathrm{coneProd}$ has the defining properties of the
product: constant term $1$, the logarithmic derivative $\mathrm{coneLog}$, multiplicativity in the exponents
(Theorem~\ref{ch19:thm:knzproduct}(b)) and uniqueness of exponents (Theorem~\ref{ch19:thm:knzproduct}(c)).  On paper,
these characterise the product: the classical infinite product has constant term $1$ and logarithmic derivative
$\mathrm{coneLog}(e)$ (termwise, as computed above), hence equals $\mathrm{coneProd}(e)$ by Theorem~\ref{ch19:thm:ode}(b).
For the downstream use (Chapter~\ref{ch:20}: dimensions of root spaces from equality of products) only these properties
are used, so the definition is adequate; but a referee should read ``$\prod$'' in the Lean docstrings as
``the series with this logarithmic derivative''.
\end{dictionary}
\begin{physics}
The multi-particle (Fock space) partition function of bosons with $e(m,n)$ species at momentum $(m,n)$ is
$\prod(1-p^mq^n)^{-e(m,n)}$; here the exponent appears with the opposite sign, as in the Weyl denominator (the ``fermionic''
or BRST side of the character formula).
\end{physics}
\begin{status}Definitions; standard axioms (\leanpath{MoonshineLRM5Audit.lean}).\end{status}

\begin{theorem}[The product identity (KNZ)]\label{ch19:thm:knzproduct}
\begin{enumerate}[label=(\alph*)]
\item $\prod_{m,n\ge1}(1-p^mq^n)^{c(mn)}=1-\sum_{a,b\ge1}c(a+b-1)\,p^aq^b$ in $\C((q))[[p]]$, where $c(n)$ are the
  coefficients of $J$ and the product is $\mathrm{coneProd}$;
\item $\mathrm{coneProd}(e_1+e_2)=\mathrm{coneProd}(e_1)\,\mathrm{coneProd}(e_2)$;
\item if $\mathrm{coneProd}(e_1)=\mathrm{coneProd}(e_2)$ then $e_1(m,n)=e_2(m,n)$ for all $m,n\ge1$.
\end{enumerate}
\end{theorem}
\begin{leancode}
theorem knz_product : coneProd (fun m n => cJ (m * n)) = knzSeries cJ
theorem coneProd_add (e₁ e₂ : ℕ → ℕ → ℂ) : coneProd (e₁ + e₂) = coneProd e₁ * coneProd e₂
theorem coneProd_exponent_unique {e₁ e₂ : ℕ → ℕ → ℂ} (h : coneProd e₁ = coneProd e₂) :
    ∀ m n : ℕ, 0 < m → 0 < n → e₁ m n = e₂ m n
\end{leancode}
(\leanpath{VOA/Moonshine/MoonshineLRM5Product.lean}.)
\begin{dictionary}
\lean{fun m n => cJ (m * n)} is the exponent $e(m,n)=c(mn)$, and \lean{knzSeries cJ} is $S$.  (a) is \eqref{ch19:eq:knz}
with the reading of $\prod$ explained in the dictionary of Definition~\ref{ch19:def:coneprod}.  (c) is stated for
$m,n\ge1$ because exponents at $m=0$ or $n=0$ do not occur in $\mathrm{coneLog}$.

\textbf{The full Borcherds form} \eqref{ch19:eq:borcherds} is not a single Lean theorem.  It is the conjunction of (a)
and Theorem~\ref{ch19:thm:knzlog}(a) at $c=\mathtt{cJ}$:
\begin{equation*}
  p\bigl(J(p)-J(q)\bigr)=(1-p/q)\cdot\mathrm{coneProd}\bigl((m,n)\mapsto c(mn)\bigr),
\end{equation*}
which, after inserting the two trivial factor classes $n\le0$ described above, is \eqref{ch19:eq:borcherds}.  (The
coordination log describes \lean{faberGen_map_eq_mul} alone as ``the KNZ$'$ form''; strictly, it is only the factorisation,
and the KNZ$'$ form needs \lean{knz_product} as well.)

\textbf{Formal versus analytic.}  The identity is formal, in $\C((q))[[p]]$.  Both sides also converge for $|q|<1$,
$|p|<|q|$, and the formal identity implies the analytic one there; the library does not state the analytic version,
and nothing downstream needs it.
\end{dictionary}
\paragraph{Proof on paper.}  (a) Both $\mathrm{coneProd}(c(mn))$ and $S$ have constant term $1$.  The first satisfies
$pA'=A\,\mathrm{coneLog}$ by construction; the second satisfies $pS'=S\cdot(-E)$ by Theorem~\ref{ch19:thm:knzlog}(c), and
$-E=\mathrm{coneLog}(c(mn))$ because the coefficient of $p^Mq^r$ in $\mathrm{coneLog}(c(mn))$ is
$-\sum_{m\mid M,\ (M/m)\mid r}m\,c\bigl(m\cdot\frac r{M/m}\bigr)=-\sum_{a\mid(M,r)}\frac Ma\,c\bigl(\frac{Mr}{a^2}\bigr)$ (put $a=M/m$),
which is $-h(M,r)$ (\lean{coneLog_cJ}).  Uniqueness (Theorem~\ref{ch19:thm:ode}(b)) gives equality.  (b) is
Theorem~\ref{ch19:thm:ode}(c) with $\mathrm{coneLog}(e_1+e_2)=\mathrm{coneLog}(e_1)+\mathrm{coneLog}(e_2)$.  (c) Equal products have
equal logarithmic derivatives (cancel the nonzero series $\mathrm{coneProd}$ in $pA'=A\,\mathrm{coneLog}$), and the coefficient
of $p^mq^n$ in $\mathrm{coneLog}(e)$ is $-\sum_{m'\mid m}[\,(m/m')\mid n\,]\,m'\,e(m',n/(m/m'))$ (\lean{coeff_coneLog}); the term
$m'=m$ is $-m\,e(m,n)$ and all others involve $m'<m$; strong induction on $m$ and division by $m\ne0$ give $e_1=e_2$. \qed
\begin{physics}
\textbf{The denominator identity as a one-loop string partition function.}  In the bosonic string on
$\Vnat\otimes V_{I\!I_{1,1}}$, the physical states of momentum $(m,n)$ form $\chnineteenm_{(m,n)}$, of dimension $c(mn)$ (no-ghost
theorem).  The Weyl--Kac--Borcherds formula says that the ``Fock space'' $U(\chnineteenu^+)$ of these states has character
$\prod(1-p^mq^n)^{-c(mn)}$ (PBW), and that it is the free (tensor) algebra on the space $U$ of simple-root strings, whose
character is $1-S$, so that its character is also $1/S$; hence $\prod(1-p^mq^n)^{c(mn)}=S$.  Harvey and Moore (1996) showed that Borcherds products arise as one-loop
threshold integrals of string theory, $\log|\text{product}|=$ a regularised integral of the partition function over the
fundamental domain; Paquette, Persson and Volpato (2016) realised the Monster Lie algebra as the algebra of BPS states of
a heterotic string compactification, with \eqref{ch19:eq:borcherds} as its BPS index.  The antisymmetry $J(p)-J(q)$
reflects the T-duality exchanging the two light-cone directions of $I\!I_{1,1}$, i.e.\ the Weyl reflection in the real
root.  These are physical interpretations; the proof above uses only modularity of $j$.
\end{physics}
\begin{status}
(a): \textbf{unconditional}; (b), (c): no hypotheses.  No Lean statement here mentions the Monster Lie algebra: (a) is an
identity about the coefficients of $J$.  Standard axioms (\leanpath{MoonshineLRM5Audit.lean}, 102 reports;
build record of the release, Chapter~\ref{ch:15}).  Supporting: \lean{charZero_laurentSeries},
\lean{ofPowerSeries_mk_zero}, \lean{constantCoeff_coneLog}, \lean{coneProd_spec}, \lean{constantCoeff_coneProd},
\lean{coneLog_add}, \lean{coneLog_cJ}, \lean{coneLog_eq_of_coneProd_eq}, \lean{coeff_coneLog}.
\end{status}

%======================================================================
\section{The twisted denominator identities}\label{ch19:sec:twisted}
%======================================================================

\subsection{The classical statement}

Let $g\in\Aut(\Vnat)$ and write $c_g(n)=\tr(g\,|\,\Vnat_{n+1})$, so $T_g=q^{-1}+\sum_{n\ge1}c_g(n)q^n$.  The group $\Aut(\Vnat)$
acts on $\chnineteenm$ by graded Lie automorphisms, and the no-ghost isomorphism $\chnineteenm_{(m,n)}\cong\Vnat_{1+mn}$ is
equivariant (Chapter~\ref{ch:18}), so $\tr(g\,|\,\chnineteenm_{(m,n)})=c_g(mn)$.  Taking $g$-traces instead of dimensions in
Borcherds' argument gives his \emph{twisted denominator identities}:
\begin{equation}\label{ch19:eq:twisted-exp}
  p^{-1}\exp\Bigl(-\sum_{k\ge1}\frac1k\sum_{m>0,\;n\in\Z}c_{g^k}(mn)\,p^{mk}q^{nk}\Bigr)=T_g(p)-T_g(q) .
\end{equation}
If $g$ has finite order $N$, so that $g$ is diagonalisable with eigenvalues $N$-th roots of unity, this can be written as a
product, $p^{-1}\prod_{m>0,n}\prod_{\zeta^N=1}(1-\zeta p^mq^n)^{\mathrm{mult}(\zeta,\,g|\Vnat_{1+mn})}=T_g(p)-T_g(q)$.  For $g=1$
it is \eqref{ch19:eq:borcherds}.  Exactly as in \S\ref{ch19:sec:denominator}, the factors with $n\le0$ give $1-p/q$, and
\eqref{ch19:eq:twisted-exp} is equivalent to
\begin{equation}\label{ch19:eq:twisted-u}
  1-u_g=\exp\Bigl(-\sum_{k\ge1}\frac1k\sum_{m,n\ge1}c_{g^k}(mn)\,p^{mk}q^{nk}\Bigr),\qquad u_g=\sum_{a,b\ge1}c_g(a+b-1)p^aq^b,
\end{equation}
and, taking $D=p\partial_p$ of the logarithm, to the logarithm-free form
\begin{equation}\label{ch19:eq:TD}
  D\,u_g=(1-u_g)\,P_g,\qquad P_g=\sum_{k\ge1}\sum_{m,n\ge1}m\,c_{g^k}(mn)\,p^{mk}q^{nk}.\tag{TD$_g$}
\end{equation}
The library proves \eqref{ch19:eq:TD} (in ``height coordinates'', below), for every symmetry of finite order of an
abstract Borcherds datum, and then for every finite-order $g\in\Aut(\Vnat)$ on $\chnineteenm$.

\paragraph{The idea of the proof.}  Chapter~\ref{ch:18} proved that the positive part $\chnineteenu^+=\bigoplus_{m,n\ge1}\chnineteenm_{(m,n)}$
satisfies $T(U)\cong U(\chnineteenu^+)$, compatibly with every symmetry: the algebra map from the tensor algebra $T(U)$ to the
universal enveloping algebra is bijective.  Equivalently, $\chnineteenu^+$ is the \emph{free Lie algebra on $U$} (the Lie subalgebra of
$T(U)$ generated by $U$ maps isomorphically onto $\chnineteenu^+$; Chapter~\ref{ch:18}, ``$\chnineteenu^+$ is free on $U$''), and this
freeness is the whole Lie-theoretic input of the twisted identities.  Here $U\subset\chnineteenu^+$ is the generating space spanned
by the real-root strings $(\mathrm{ad}\,e)^k\chnineteenm_{(1,j)}$, $0\le k<j$; the string through $(1,j)$ sits in degrees
$(1+k,j-k)$ and is isomorphic to $\chnineteenm_{(1,j)}$ as a representation of the symmetry group (which fixes $e$).  Compare the
$g$-traces of the two sides degree by degree:
\begin{itemize}
\item on $T(U)=\bigoplus_\ell U^{\otimes\ell}$, the trace of $g$ on a tensor product is the product of traces, so the
  graded trace is $\sum_\ell(\mathrm{ch}_gU)^\ell=(1-\mathrm{ch}_gU)^{-1}$, and $\mathrm{ch}_gU=u_g$ (the string through $(1,j)$
  contributes $c_g(j)$ in each degree $(1+k,j-k)$, i.e.\ $c_g(a+b-1)$ at $(a,b)$);
\item on $U(\chnineteenu^+)$, choose a basis of eigenvectors $y_a$ of $g$, homogeneous of degree $\gamma_a$, eigenvalue
  $\lambda_a$; by PBW the sorted monomials form a basis of eigenvectors, so the graded trace is
  $B=\prod_a(1-\lambda_ae^{\gamma_a})^{-1}$, and $DB/B=\sum_{k\ge1}\sum_a(\gamma_a)_1\lambda_a^ke^{k\gamma_a}=P_g$, because
  $\sum_{a:\gamma_a=(m,n)}\lambda_a^k=\tr(g^k|\chnineteenm_{(m,n)})=c_{g^k}(mn)$.
\end{itemize}
Equating, $B=(1-u_g)^{-1}$ and $DB=BP_g$, which is \eqref{ch19:eq:TD}.  The library carries out exactly this argument,
but without infinite products: both sides are replaced by \emph{weighted counts} of words and multisets of basis
letters, and the comparison is made by transporting a linear isomorphism between coordinate spaces.

\subsection{The statement in Lean: height coordinates}

\begin{definition}[Euler derivation; the twisted series; the twisted denominator identity]\label{ch19:def:twistedTD}
In $\C[[q]][[p]]$ let $D=p\,\partial_p$.  For $c:\N\to\N\to\C$ (with $c(k,\cdot)$ the coefficient sequence of the $k$-th
replicate) define
\begin{equation*}
  \tilde u=\sum_{a,b\ge1}c(1,a+b-1)\,p^aq^{a+b},\qquad
  \tilde P=\sum_{k\ge1}\sum_{a,b\ge1}a\,c(k,ab)\,p^{ak}q^{(a+b)k},
\end{equation*}
and say that $c$ satisfies the \emph{twisted denominator identity} if $D\tilde u=(1-\tilde u)\tilde P$.
\end{definition}
\begin{leancode}
def eulerDeriv (F : PowerSeries R) : PowerSeries R :=
  PowerSeries.mk fun n => (n : R) * PowerSeries.coeff n F
def twistedUnitSeries (c : ℕ → ℂ) : PowerSeries (PowerSeries ℂ) :=
  PowerSeries.mk fun a => PowerSeries.mk fun N => if 1 ≤ a ∧ a < N then c (N - 1) else 0
def twistedPowerSum (c : ℕ → ℕ → ℂ) : PowerSeries (PowerSeries ℂ) :=
  PowerSeries.mk fun A => PowerSeries.mk fun N =>
    ∑ k ∈ A.divisors,
      if k ∣ N ∧ A < N then ((A / k : ℕ) : ℂ) * c k (A / k * ((N - A) / k)) else 0
def TwistedDenominatorIdentity (c : ℕ → ℕ → ℂ) : Prop :=
  eulerDeriv (twistedUnitSeries (c 1)) = (1 - twistedUnitSeries (c 1)) * twistedPowerSum c
\end{leancode}
(\leanpath{VOA/Moonshine/MoonshineCNReplicationFormal.lean}, namespace \lean{VOA.Moonshine.ConwayNorton}; \lean{eulerDeriv}
is in a section with \lean{{R : Type*} [CommRing R]}.)
\begin{dictionary}
\textbf{Decoding.}  \lean{eulerDeriv F} multiplies the $p^n$-coefficient by $n$: it is $D$.  In \lean{twistedUnitSeries c},
the coefficient of $p^aq^N$ is $c(N-1)$ if $1\le a<N$ and $0$ otherwise; writing $N=a+b$ with $b\ge1$ this is
$\sum_{a,b\ge1}c(a+b-1)p^aq^{a+b}$.  In \lean{twistedPowerSum c}, the coefficient of $p^Aq^N$ is
$\sum_{k\mid A,\ k\mid N,\ A<N}\frac Ak\,c\bigl(k,\frac Ak\cdot\frac{N-A}k\bigr)$; writing $A=ak$, $N=(a+b)k$ this is
$\sum_k a\,c(k,ab)$, as displayed.  \lean{TwistedDenominatorIdentity c} uses only $c(1,\cdot)$ in $\tilde u$ and all
$c(k,\cdot)$, $k\ge1$, in $\tilde P$; the values $c(0,\cdot)$ never occur.

\textbf{Height coordinates.}  A root $(m,n)$ of the positive cone is recorded by the exponent $p^mq^{m+n}$ (first
coordinate and height).  So $\tilde u(p,q)=u_g(pq,q)$ and $\tilde P(p,q)=P_g(pq,q)$ with $c(k,n)=c_{g^k}(n)$.  The substitution
$p\mapsto pq$ is an injective ring map that commutes with $D$ (it does not change the $p$-degree), so
$D\tilde u=(1-\tilde u)\tilde P$ holds if and only if \eqref{ch19:eq:TD} holds.  \emph{Why height coordinates?}  Not for the
counting: the positive cone $\{m,n\ge1\}$, which contains all degrees of $\chnineteenu^+$ and of $U$ (the strings sit at $(1+k,j-k)$
with $0\le k<j$), already lies in $\N\times\N$ with positive first coordinate and finite fibres in Borcherds' coordinates.  The reason
is the replication step (\S\ref{ch19:sec:replication}): there one substitutes $X\mapsto f(q)$ in $G(p)=p(f(p)-X)$, which gives
$(1-p/q)(1-u_g)$, with a negative power of $q$; after the further substitution $p\mapsto pq$ this becomes $(1-p)(1-\tilde u)$, and the
whole computation stays in $\C[[q]][[p]]$ (\lean{faberEvaluation_faberGenerating}).  The twisted identity is stated in the same
coordinates so that it can be used there directly; the torus characters $2^{A}$, $2^{N}$ of the proof below then record the
height degree $(A,N)$.

\textbf{Same as the textbook identity?}  Yes, for every $g$ of finite order: \eqref{ch19:eq:TD} is equivalent to
\eqref{ch19:eq:twisted-u} by the uniqueness of solutions of $pA'=AL$ (Theorem~\ref{ch19:thm:ode}) and to
\eqref{ch19:eq:twisted-exp} by the factorisation $p(T_g(p)-T_g(q))=(1-p/q)(1-u_g)$ (Theorem~\ref{ch19:thm:knzlog}(a) with
$c=c_g$).  These two equivalences are paper arguments; the library states and uses only the logarithm-free form, which is
all that replication needs.  \textbf{Could the definition be vacuous?}  No: for $c(k,n)=c_J(n)$ for all $k$ (the trivial
symmetry), it is the logarithmic KNZ identity, which holds (Theorem~\ref{ch19:thm:knzlog}(c)); and a family not coming from a
symmetry generally fails it (it would otherwise force replication, which fails for most sequences).
\end{dictionary}
\begin{physics}
$\tilde u$ counts single ``simple-root'' string states with the symmetry $g$ inserted, and $\tilde P$ is the connected
(single-particle) free energy of all physical states with $g$ inserted, summed over $k$-fold multi-wound copies, which carry
the twist $g^k$.  The identity is the $g$-twisted index of the string's Fock space: inserting $g$ in the trace over a
symmetric (bosonic) Fock space produces the plethystic exponential $\exp\sum_k\frac1k\tr(g^k\cdots)$, whose $k$-th term is a
$k$-fold cover twisted by $g^k$, exactly as in the orbifold of a symmetric product.
\end{physics}
\begin{status}Definitions; standard axioms (\leanpath{VOA/Moonshine/MoonshineCNAudit18.lean}).\end{status}

\subsection{Words and multisets: the two character identities}

The file \leanpath{VOA/Moonshine/MoonshineCNTwistedCombinatorics.lean} contains no Lie algebra.  It proves the two
generating-function identities ``$T(U)$ has character $(1-\mathrm{ch}\,U)^{-1}$'' and ``$\mathrm{Sym}$ has character the plethystic
exponential'', for abstract weighted letters.

\begin{definition}[Graded letters; weighted counts]\label{ch19:def:letters}
A set of \emph{graded letters} is a set $I$ with a degree $\deg:I\to\N\times\N$ such that $(\deg i)_1\ge1$ and each fibre
$\deg^{-1}(\gamma)$ is finite.  Given weights $\mathrm{wt}:I\to\C$, define for $\gamma\in\N\times\N$:
\begin{itemize}
\item the word count $W(\gamma)=\sum_{w:\deg w=\gamma}\prod_{i\in w}\mathrm{wt}(i)$ over words (finite sequences) of total degree $\gamma$;
\item the letter count $\ell(\gamma)=\sum_{\deg i=\gamma}\mathrm{wt}(i)$;
\item the multiset count $M(\gamma)=\sum_{m:\deg m=\gamma}\prod_{i\in m}\mathrm{wt}(i)$ over finite multisets;
\item the power-sum count $P(\varepsilon)=\sum_{k\ge1}\sum_{i:\,k\deg i=\varepsilon}(\deg i)_1\,\mathrm{wt}(i)^k$.
\end{itemize}
\end{definition}
\begin{leancode}
structure GradedLetters (ι : Type*) where
  deg : ι → ℕ × ℕ
  deg_pos : ∀ i, 1 ≤ (deg i).1
  finite_fiber : ∀ γ, {i | deg i = γ}.Finite
def wordCount (wt : ι → ℂ) (γ : ℕ × ℕ) : ℂ := ∑ w ∈ L.wordFiber γ, (w.map wt).prod
def letterCount (wt : ι → ℂ) (δ : ℕ × ℕ) : ℂ :=
  ∑ i ∈ L.lowerLetters δ, if L.deg i = δ then wt i else 0
def msetCount (wt : ι → ℂ) (γ : ℕ × ℕ) : ℂ := ∑ m ∈ L.msetFiber γ, (m.map wt).prod
def powerSumCount (wt : ι → ℂ) (ε : ℕ × ℕ) : ℂ :=
  ∑ x ∈ L.lowerLetters ε ×ˢ Icc 1 ε.1,
    if x.2 • L.deg x.1 = ε then ((L.deg x.1).1 : ℂ) * wt x.1 ^ x.2 else 0
\end{leancode}
(\leanpath{VOA/Moonshine/MoonshineCNTwistedCombinatorics.lean}; the counts are in \lean{namespace GradedLetters} with
\lean{L : GradedLetters ι}.)
\begin{dictionary}
\lean{L.wordFiber γ} is the finite set of lists of degree $\gamma$ (finite because every letter has first degree $\ge1$,
so a word of degree $\gamma$ has at most $\gamma_1$ letters, and the fibres are finite); \lean{L.msetFiber γ} likewise
for multisets.  \lean{(w.map wt).prod} is $\prod_{i\in w}\mathrm{wt}(i)$.  \lean{L.lowerLetters δ} is the finite set of letters
of degree $\le\delta$ componentwise; the \lean{if L.deg i = δ} selects degree exactly $\delta$.  In \lean{powerSumCount},
\lean{Icc 1 ε.1} is $1\le k\le\varepsilon_1$ (no larger $k$ can satisfy $k\deg i=\varepsilon$), \lean{x.2 • L.deg x.1} is $k\deg i$,
and \lean{(L.deg x.1).1} is $(\deg i)_1$.  The generating series of a count $f$ is
\lean{bigradedSeries f} $=\sum_{A,N}f(A,N)p^Aq^N$ (\lean{wordSeries}, \lean{letterSeries}, \lean{msetSeries},
\lean{powerSumSeries}).
\end{dictionary}

\begin{physics}
Graded letters and weighted counts are the bookkeeping of a Fock space: a word in graded letters is a multi-particle
state, and the weighted count is its contribution to a character, with weights recording the charges (momenta and
energies) of the constituent particles.
\end{physics}

\begin{theorem}[The two character identities]\label{ch19:thm:characters}
For graded letters with weights:
\begin{enumerate}[label=(\alph*)]
\item $W(\gamma)=[\gamma=0]+\sum_i\mathrm{wt}(i)\,W(\gamma-\deg i)$, hence $W\cdot(1-\ell)=1$ as series;
\item $\gamma_1M(\gamma)=\sum_\varepsilon M(\gamma-\varepsilon)P(\varepsilon)$, hence $DM=M\cdot P$ as series;
\item in $\C[[q]][[p]]$: if $B(1-u)=1$ and $DB=BP$, then $Du=(1-u)P$.
\end{enumerate}
\end{theorem}
\begin{leancode}
theorem wordCount_eq (wt : ι → ℂ) (γ : ℕ × ℕ) :
    L.wordCount wt γ = (if γ = 0 then 1 else 0) +
      ∑ i ∈ L.lowerLetters γ, wt i * L.wordCount wt (γ - L.deg i)
theorem wordSeries_mul_one_sub_letterSeries (wt : ι → ℂ) :
    L.wordSeries wt * (1 - L.letterSeries wt) = 1
theorem firstDegree_mul_msetCount_eq (wt : ι → ℂ) (γ : ℕ × ℕ) :
    (γ.1 : ℂ) * L.msetCount wt γ = bigradedConv (L.msetCount wt) (L.powerSumCount wt) γ
theorem eulerDeriv_msetSeries (wt : ι → ℂ) :
    eulerDeriv (L.msetSeries wt) = L.msetSeries wt * L.powerSumSeries wt
theorem eulerDeriv_eq_of_inverse {B u P : PowerSeries (PowerSeries ℂ)} (h1 : B * (1 - u) = 1)
    (h2 : eulerDeriv B = B * P) : eulerDeriv u = (1 - u) * P
\end{leancode}
(\leanpath{VOA/Moonshine/MoonshineCNTwistedCombinatorics.lean}.)
\begin{dictionary}
\lean{γ - L.deg i} is truncated subtraction in $\N\times\N$, harmless because only letters of degree $\le\gamma$ occur.
\lean{bigradedConv f g γ} is the Cauchy convolution $\sum_{\delta+\delta'=\gamma}f(\delta)g(\delta')$ (\lean{bigradedSeries_mul}
identifies it with multiplication of series).  (b) is the logarithm-free form of
$M=\exp\sum_{k\ge1}\frac1k\sum_i\mathrm{wt}(i)^ke^{k\deg i}$, i.e.\ of $M=\prod_i(1-\mathrm{wt}(i)e^{\deg i})^{-1}$.
\end{dictionary}
\paragraph{Proof on paper.}  (a) Split a nonempty word into its first letter and the rest.  (b) For a multiset $m$ of degree
$\gamma$, $\gamma_1=\sum_{i}(\text{multiplicity of }i\text{ in }m)(\deg i)_1=\sum_i\sum_{k\ge1}[\,\mathrm{mult}_i(m)\ge k\,](\deg i)_1$.  So
$\gamma_1M(\gamma)=\sum_i\sum_{k\ge1}(\deg i)_1\sum_{m\ni i^k}\mathrm{wt}(m)$, and the multisets of degree $\gamma$ containing at least $k$
copies of $i$ are exactly $\{i^k\}\cup m'$ with $\deg m'=\gamma-k\deg i$, of weight $\mathrm{wt}(i)^k\mathrm{wt}(m')$
(\lean{sum_msetFiber_count_ge}).  Hence $\gamma_1M(\gamma)=\sum_\varepsilon P(\varepsilon)M(\gamma-\varepsilon)$.  (c) $B$ is a unit with
inverse $1-u$; from $D(B(1-u))=0$ and the Leibniz rule, $DB\,(1-u)=B\,Du$, so $B\,Du=BP(1-u)$; multiply by $1-u$. \qed
\begin{physics}
(a) is the partition function of a ``free'' (no relations) algebra of creation operators: $T(U)$ is the Fock space of
\emph{distinguishable} insertions.  (b) is the bosonic Fock space: the grand-canonical partition function of bosons
is $\prod(1-x)^{-1}=\exp\sum_k x^k/k$, the $k$-th term being $k$ particles in the same state (a $k$-fold winding).
\end{physics}
\begin{status}No hypotheses; standard axioms (\leanpath{MoonshineCNAudit18.lean}).  Supporting declarations (same file):
\lean{bigradedSeries}, \lean{bigradedConv}, \lean{bigradedSeries_ext}, \lean{bigradedSeries_mul}, \lean{one_eq_bigradedSeries},
\lean{eulerDeriv_bigradedSeries}, \lean{bigradedConv_congr}, \lean{bigradedConv_comm}, \lean{sum_antidiagonal_ite_fst},
\lean{sum_shift_eq_bigradedConv}, \lean{prod_add_tsub_cancel_left/right}, \lean{GradedLetters.finite_le},
\lean{lowerLetters}, \lean{listDeg}, \lean{finite_wordFiber}, \lean{wordFiber}, \lean{msetDeg}, \lean{card_le_msetDeg},
\lean{finite_msetFiber}, \lean{msetFiber}, \lean{msetDeg_le_of_le}, \lean{sum_msetFiber_count_ge}; and the
\lean{eulerDeriv} calculus of \leanpath{MoonshineCNReplicationFormal.lean} (\lean{eulerDeriv_add}, \lean{_sub},
\lean{_neg}, \lean{_one}, \lean{_X}, \lean{_mul} (Leibniz), \lean{_map}, \lean{_rescale}).\end{status}

\subsection{Transporting weighted counts along PBW}

\begin{theorem}[Transfer of weighted counts]\label{ch19:thm:transfer}
Let $\Psi:\C^{(X)}\to\C^{(Y)}$ be a linear isomorphism between free vector spaces on sets $X,Y$, and let $\kappa^X_j:X\to\C$,
$\kappa^Y_j:Y\to\C$ ($j\in J$) be families of labels such that $\Psi$ intertwines the diagonal operators
$e_x\mapsto\kappa^X_j(x)e_x$ and $e_y\mapsto\kappa^Y_j(y)e_y$ for every $j$.  Then for every property $P$ of label vectors and
every $j_0$, $\sum_{x:P(\kappa^X(x))}\kappa^X_{j_0}(x)=\sum_{y:P(\kappa^Y(y))}\kappa^Y_{j_0}(y)$ (for finite such sets).
\end{theorem}
\begin{leancode}
def finsuppDiag {X : Type*} (κ : X → ℂ) : (X →₀ ℂ) →ₗ[ℂ] (X →₀ ℂ) :=
  Finsupp.linearCombination ℂ fun x => Finsupp.single x (κ x)
theorem sum_eq_sum_of_diag_equiv (Ψ : (X →₀ ℂ) ≃ₗ[ℂ] (Y →₀ ℂ))
    (κX : J → X → ℂ) (κY : J → Y → ℂ)
    (h : ∀ j v, Ψ (finsuppDiag (κX j) v) = finsuppDiag (κY j) (Ψ v))
    (P : (J → ℂ) → Prop) (S : Finset X) (T : Finset Y)
    (hS : ∀ x, x ∈ S ↔ P fun j => κX j x) (hT : ∀ y, y ∈ T ↔ P fun j => κY j y) (j₀ : J) :
    ∑ x ∈ S, κX j₀ x = ∑ y ∈ T, κY j₀ y
\end{leancode}
(\leanpath{VOA/Moonshine/MoonshineCNTwistedTransfer.lean}.)
\begin{dictionary}
\lean{X →₀ ℂ} is the space of finitely supported functions $X\to\C$, i.e.\ the free vector space $\C^{(X)}$ with basis
$e_x=$\lean{single x 1}; \lean{≃ₗ[ℂ]} is a linear isomorphism; \lean{finsuppDiag κ} is the diagonal operator.  The
finsets \lean{S}, \lean{T} are the points whose label vector $j\mapsto\kappa_j$ satisfies \lean{P}.
\end{dictionary}
\paragraph{Proof on paper.}  An operator intertwining diagonal operators maps the joint eigenspace of a label vector
$\sigma$ into the joint eigenspace of $\sigma$ (\lean{apply_single_eq_zero_of_diag}); the joint eigenspaces are spanned by the
basis vectors with that label; an isomorphism therefore matches their dimensions, i.e.\ the cardinalities of label fibres
(\lean{card_le_of_injective_support}, \lean{card_fiber_eq_of_diag_equiv}).  A weighted count over a union of fibres is a
sum over fibres of (cardinality)$\times$(common value of $\kappa_{j_0}$). \qed
\begin{physics}Equal traces of commuting diagonal operators on isomorphic Hilbert spaces: if two bases of the same space are
both eigenbases of the same commuting family (charges), each joint eigenvalue has the same multiplicity in both.\end{physics}
\begin{status}No hypotheses; standard axioms (\leanpath{MoonshineCNAudit18.lean}).\end{status}

\begin{theorem}[PBW and word coordinates intertwine diagonal symmetries]\label{ch19:thm:wordpbw}
Let $L$ be a complex Lie algebra with a linearly ordered basis $(Y_a)_{a\in A}$, $S\subseteq L$ a subspace with basis
$(Z_i)_{i\in I}$, and suppose the algebra map $T(S)\to U(L)$ is bijective.  Let $\Psi:\C^{(\mathrm{List}\,I)}\to\C^{(\mathrm{Multiset}\,A)}$
be the composite of the word basis of $T(S)$, the bijection $T(S)\to U(L)$ and the PBW coordinates of $U(L)$.  If a Lie
endomorphism $\varphi$ of $L$ preserves $S$ and is diagonal in both bases ($\varphi Y_a=\kappa_aY_a$, $\varphi Z_i=\nu_iZ_i$), then
$\Psi$ intertwines $e_w\mapsto\prod_{i\in w}\nu_i\,e_w$ with $e_m\mapsto\prod_{a\in m}\kappa_a\,e_m$.
\end{theorem}
\begin{leancode}
def pbwCoordinates : UniversalEnvelopingAlgebra ℂ L ≃ₗ[ℂ] (Multiset A →₀ ℂ) :=
  envelopingInducedEquiv ≪≫ₗ
    relativePBWEquiv (⊥ : LieSubalgebra ℂ L) zeroBottomRep bot_isCompl_top (topBasis Y)
theorem pbwCoordinates_sorted (l : List A) (hl : l.Pairwise (· ≤ ·)) :
    pbwCoordinates Y (l.map fun a => UniversalEnvelopingAlgebra.ι ℂ (Y a)).prod =
      single (l : Multiset A) 1
def wordPBWEquiv : (List I →₀ ℂ) ≃ₗ[ℂ] (Multiset A →₀ ℂ) :=
  Finsupp.domLCongr (R := ℂ) (M := ℂ) FreeMonoid.toList.symm ≪≫ₗ Z.tensorAlgebra.repr.symm ≪≫ₗ
    LinearEquiv.ofBijective (tensorGeneratorToEnveloping S).toLinearMap hbij ≪≫ₗ pbwCoordinates Y
theorem wordPBWEquiv_diag (φ : L →ₗ⁅ℂ⁆ L) (hφ : ∀ u ∈ S, φ u ∈ S) (κ : A → ℂ) (ν : I → ℂ)
    (hY : ∀ a, φ (Y a) = κ a • Y a) (hZ : ∀ i, φ (Z i : L) = ν i • (Z i : L))
    (v : List I →₀ ℂ) :
    wordPBWEquiv S hbij Y Z (finsuppDiag (fun w => (w.map ν).prod) v) =
      finsuppDiag (fun m => (m.map κ).prod) (wordPBWEquiv S hbij Y Z v)
\end{leancode}
(\leanpath{VOA/Moonshine/MoonshineCNTwistedPBW.lean}; section variables \lean{(Y : Basis A ℂ L)} with \lean{[LinearOrder A]},
\lean{(S : Submodule ℂ L)}, \lean{(hbij : Function.Bijective (tensorGeneratorToEnveloping S))},
\lean{(Z : Basis I ℂ S)}.)
\begin{dictionary}
\lean{pbwCoordinates Y} is a linear isomorphism $U(L)\cong\C^{(\mathrm{Multiset}\,A)}$ obtained from the repository's relative PBW
theorem with the zero subalgebra (Chapter~\ref{ch:18}): \lean{pbwCoordinates_sorted} says that a sorted monomial
$Y_{a_1}\cdots Y_{a_r}$ ($a_1\le\dots\le a_r$) has coordinate vector $e_{\{a_1,\dots,a_r\}}$; this is PBW.
\lean{tensorGeneratorToEnveloping S} is the algebra map $T(S)\to U(L)$ of Chapter~\ref{ch:18}.  \lean{Z.tensorAlgebra} is the
basis of $T(S)$ by words (Mathlib, \lean{FreeMonoid I}, converted to lists).  \lean{wordPBWEquiv_diag} is the statement
displayed.
\end{dictionary}
\paragraph{Proof on paper.}  $\varphi$ acts on $U(L)$ by an algebra automorphism $U(\varphi)$; on a sorted monomial of eigenvectors it
multiplies by the product of eigenvalues (\lean{pbwCoordinates_envelopingMap}).  Likewise on words in $T(S)$
(\lean{tensorAlgebra_repr_lift_diag}).  The map $T(S)\to U(L)$ is natural in $\varphi$ (\lean{tensorGeneratorToEnveloping_natural},
Chapter~\ref{ch:18}). \qed
\begin{physics}The free-field (word) description and the oscillator (PBW) description of the same Fock space carry the same
symmetry charges.\end{physics}
\begin{status}Hypotheses as displayed (bijectivity of $T(S)\to U(L)$ is an explicit argument of \lean{wordPBWEquiv}); standard
axioms (\leanpath{MoonshineCNAudit18.lean}).  Supporting (same file): \lean{list_prod_map_smul}, \lean{finsuppDiag_domLCongr},
\lean{topBasis}, \lean{bot_isCompl_top}, \lean{zeroBottomRep}, \lean{balancing_bot_eq_bot}, \lean{envelopingInducedEquiv},
\lean{relativeInducedWord_bot}, \lean{pbwCoordinates_symm_single}, \lean{pbwCoordinates_envelopingMap},
\lean{tensorAlgebra_basis_apply}, \lean{tensorAlgebra_repr_lift_diag}; (file \leanpath{MoonshineCNTwistedTransfer.lean})
\lean{finsuppDiag_single}, \lean{finsuppDiag_apply}, \lean{apply_single_eq_zero_of_diag}, \lean{card_le_of_injective_support},
\lean{card_fiber_eq_of_diag_equiv}.\end{status}

\subsection{Eigenbases, root-space bases and torus elements}

\begin{definition}[A symmetry on one root space]\label{ch19:def:weightAction}
For a Borcherds datum $D$ on $\mathfrak g$ (Chapter~\ref{ch:18}) and a group action $A$ of $G$ by graded, form-preserving Lie
automorphisms fixing $\mathfrak g_0$ and $\mathfrak g_{\pm\alpha_r}$ (a \lean{D.GAction G}), $\sigma\in G$ restricts to an endomorphism of each
root space $\mathfrak g_\gamma$.
\end{definition}
\begin{leancode}
def weightAction (σ : G) (γ : ℤ × ℤ) : Module.End ℂ (D.weightSpace γ) :=
  (A.action σ).toLinearMap.restrict fun x hx => A.graded σ γ x hx
\end{leancode}
(\leanpath{VOA/Moonshine/MoonshineCNTwistedBases.lean}, namespace \lean{BorcherdsData}, section variables
\lean{(D : BorcherdsData g)}, \lean{{G : Type*} [Group G] (A : D.GAction G)}.)
\begin{status}Definition; \lean{weightAction_pow} ($\sigma^k$ restricts to the $k$-th power) is printed in
\leanpath{MoonshineCNAudit18.lean}.\end{status}

\begin{physics}
A symmetry acting on one root space is a global symmetry of the string theory restricted to the physical states of a
fixed spacetime momentum. Its trace there is the corresponding twisted contribution to the one-loop amplitude, the
ingredient of the twisted denominator identity.
\end{physics}

\begin{theorem}[Eigenbases of finite-order operators; traces of powers]\label{ch19:thm:eigen}
If $T$ is an endomorphism of a finite-dimensional complex vector space with $T^N=1$ for some $N\ge1$, then $T$ has a basis of
eigenvectors; and for any eigenbasis with eigenvalues $\mu_i$, $\tr(T^k)=\sum_i\mu_i^k$.
\end{theorem}
\begin{leancode}
theorem exists_eigenbasis_of_pow_eq_one {V : Type u} [AddCommGroup V] [Module ℂ V]
    [FiniteDimensional ℂ V] (T : Module.End ℂ V) {N : ℕ} (hN : 0 < N) (hT : T ^ N = 1) :
    ∃ (ι : Type u) (_ : Fintype ι) (b : Basis ι ℂ V) (μ : ι → ℂ), ∀ i, T (b i) = μ i • b i
theorem trace_pow_eq_sum_of_eigenbasis {V : Type*} [AddCommGroup V] [Module ℂ V]
    {ι : Type*} [Fintype ι] (b : Basis ι ℂ V) (T : Module.End ℂ V) (μ : ι → ℂ)
    (h : ∀ i, T (b i) = μ i • b i) (k : ℕ) :
    LinearMap.trace ℂ V (T ^ k) = ∑ i, μ i ^ k
\end{leancode}
(\leanpath{VOA/Moonshine/MoonshineCNTwistedBases.lean}.)
\paragraph{Proof on paper.}  The minimal polynomial of $T$ divides $X^N-1$, which is separable over $\C$; a separable
minimal polynomial means $T$ is semisimple, hence diagonalisable over the algebraically closed field $\C$.  The trace of
$T^k$ in an eigenbasis is the sum of the diagonal entries $\mu_i^k$. \qed
\begin{physics}A unitary symmetry of finite order has a basis of states with definite charge (eigenvalue a root of unity),
and the twisted partition function counts states weighted by their charge.\end{physics}
\begin{status}No hypotheses beyond those displayed; standard axioms (\leanpath{MoonshineCNAudit18.lean}).\end{status}

The remaining declarations of \leanpath{MoonshineCNTwistedBases.lean} build the bases in which the comparison is made.
Given bases $b_\gamma$ of all root spaces: \lean{positiveBasis} is the basis of $\chnineteenu^+$ collected from the $b_\gamma$,
$\gamma$ in the positive cone (using that $\chnineteenu^+$ is the internal direct sum of its root spaces,
\lean{positiveWeight_isInternal}); \lean{generatorBasis} is the basis of the generating space $U$ formed by the string vectors
$(\mathrm{ad}\,e)^k b_{(1,j),i}$, $0\le k<j$ (using \lean{generatorStrings_isInternal} and the string isomorphisms of
Chapter~\ref{ch:18}), with \lean{generatorBasis_coe}, \lean{generatorBasis_mem} recording that the vector indexed by the string
$(j,k)$ lies in degree $(1+k,j-k)$.  The \emph{torus elements} \lean{weightTorus χ} multiply $\mathfrak g_\gamma$ by $\chi(\gamma)$;
for multiplicative $\chi$ they are Lie endomorphisms (\lean{weightTorus_lie}) and restrict to $\chnineteenu^+$
(\lean{positiveTorus}) and to $U$ (\lean{weightTorus_genSpace}, \lean{positiveTorus_generator_mem}).  Finally, if the $b_\gamma$ are
eigenbases of $\sigma$ with eigenvalues $\mu_\gamma$, then \lean{positiveBasis} and \lean{generatorBasis} are eigenbases of $\sigma$ on
$\chnineteenu^+$ and $U$, the string vector $(\mathrm{ad}\,e)^kb_{(1,j),i}$ having the eigenvalue $\mu_{(1,j),i}$ because $\sigma$ fixes $e$
(\lean{positivePartAction_positiveBasis}, \lean{positivePartAction_generatorBasis}), and eigenbases of every torus element
(\lean{positiveTorus_positiveBasis}, \lean{positiveTorus_generatorBasis}).  Other supporting declarations:
\lean{isInternal_comap_subtype}, \lean{borcherdsStringIndex_decidableEq}, \lean{weightDecompose},
\lean{weightDecompose_of_mem}, \lean{weightTorus_of_mem}, \lean{mem_weightSpace_iSup}, \lean{weightTorus_mem_weightSpan},
\lean{positiveTorus_coe}, \lean{positiveBasis_coe}, \lean{positiveBasis_mem}, \lean{weightAction_pow}.  All are printed in
\leanpath{MoonshineCNAudit18.lean} with standard axioms.

\subsection{The twisted denominator identity for Borcherds data}

\begin{theorem}[Twisted denominator identity, abstract form]\label{ch19:thm:twistedabstract}
Let $D$ be a Borcherds datum with a $G$-action $A$ and $\sigma\in G$.  In (a) and (b) we are moreover given a normalised real pair
$(e,f)$ and bases $b_\gamma$ of the root spaces; (c) has neither (only $D$, $A$, $\sigma$, finite order, the family $c$ and the trace
hypothesis).
\begin{enumerate}[label=(\alph*)]
\item If every $b_\gamma$ is an eigenbasis of $\sigma$ with eigenvalues $\mu_\gamma$, then for every bidegree the word count of
  the letters of $U$ weighted by eigenvalues equals the multiset count of the letters of $\chnineteenu^+$ so weighted (on paper these counts are the graded traces
  $\tr(\sigma|T(U)_\gamma)$ and $\tr(\sigma|U(\chnineteenu^+)_\gamma)$, computed in the word and PBW eigenbases; the Lean statement is about the
  counts).
\item If moreover $\sum_i\mu_{\gamma,i}^k=c(k,\gamma_1\gamma_2)$ for all $k\ge1$ and all $\gamma$ in the positive cone, then $c$
  satisfies the twisted denominator identity.
\item If $\sigma$ has finite order and $\tr(\sigma^k|\mathfrak g_\gamma)=c(k,\gamma_1\gamma_2)$ for all $k\ge1$ and $\gamma$ in the positive
  cone, then $c$ satisfies the twisted denominator identity.
\end{enumerate}
\end{theorem}
\begin{leancode}
def borcherdsHeightDegree (γ : ℤ × ℤ) : ℕ × ℕ := (γ.1.toNat, (γ.1 + γ.2).toNat)
theorem generatorLetters_wordCount_eq (σ : G)
    (hb : ∀ γ i, D.weightAction A σ γ (b γ i) = μ γ i • b γ i) (γ : ℕ × ℕ) :
    (generatorLetters ι).wordCount (fun x => μ _ x.2) γ =
      (positiveLetters ι).msetCount (fun x => μ x.1.1 x.2) γ
theorem twistedDenominatorIdentity_of_eigenbasis (σ : G)
    (hb : ∀ γ i, D.weightAction A σ γ (b γ i) = μ γ i • b γ i) (c : ℕ → ℕ → ℂ)
    (htr : ∀ k : ℕ, 1 ≤ k → ∀ γ : ℤ × ℤ, borcherdsPositiveCone γ →
      ∑ i, μ γ i ^ k = c k (γ.1 * γ.2).toNat) :
    TwistedDenominatorIdentity c
theorem twistedDenominatorIdentity_of_finOrder (σ : G) (hσ : IsOfFinOrder σ) (c : ℕ → ℕ → ℂ)
    (htr : ∀ k : ℕ, 1 ≤ k → ∀ γ : ℤ × ℤ, borcherdsPositiveCone γ →
      LinearMap.trace ℂ (D.weightSpace γ) (D.weightAction A (σ ^ k) γ) =
        c k (γ.1 * γ.2).toNat) :
    TwistedDenominatorIdentity c
\end{leancode}
(\leanpath{VOA/Moonshine/MoonshineCNTwistedCharacters.lean}, \leanpath{VOA/Moonshine/MoonshineCNTwistedDenominator.lean};
namespace \lean{BorcherdsData}; section variables \lean{(e : g) (he : e ∈ D.weightSpace borcherdsRealRoot)},
\lean{(f : g) (hf : f ∈ D.weightSpace (-borcherdsRealRoot)) (hef : D.form e f = 1)},
\lean{(b : ∀ γ, Basis (ι γ) ℂ (D.weightSpace γ)) (μ : ∀ γ, ι γ → ℂ)} with \lean{[∀ γ, Fintype (ι γ)]}, and
\lean{{G : Type*} [Group G] (A : D.GAction G)}.)
\begin{dictionary}
\lean{borcherdsHeightDegree} sends $(m,n)$ to $(m,m+n)$ (with \lean{toNat}, harmless on the positive cone).
\lean{positiveLetters ι} and \lean{generatorLetters ι} are the graded letters given by the bases of $\chnineteenu^+$ and $U$,
with degrees in height coordinates; a letter of \lean{generatorLetters} is a pair (string index $(j,k)$, basis index $i$ of
$\mathfrak g_{(1,j)}$), with weight \lean{μ _ x.2} $=\mu_{(1,j),i}$.  \lean{borcherdsPositiveCone γ} is $\gamma_1,\gamma_2\ge1$, and
\lean{(γ.1 * γ.2).toNat} is $mn$ for $\gamma=(m,n)$.  \lean{LinearMap.trace ℂ V T} is the trace.  In (c) the hypothesis \lean{htr}
is ``the traces of the powers of $\sigma$ on the positive root spaces are given by $c$''; finite order is \lean{IsOfFinOrder σ}.
The real pair $(e,f)$ is obtained inside the proof of (c) from \lean{exists_normalized_real_pair} (Chapter~\ref{ch:18}), so it is
not a hypothesis of (c).
\end{dictionary}
\paragraph{Proof on paper.}  (a) Apply Theorem~\ref{ch19:thm:transfer} to the isomorphism $\Psi$ of Theorem~\ref{ch19:thm:wordpbw} with
$L=\chnineteenu^+$, $S=U$ (bijectivity of $T(U)\to U(\chnineteenu^+)$ is \lean{positiveTensorGeneratorToEnveloping_bijective},
Chapter~\ref{ch:18}) and three diagonal operators: $\sigma$, and the torus elements with $\chi_1(\gamma)=2^{\gamma_1}$ and
$\chi_2(\gamma)=2^{\gamma_1+\gamma_2}$ (\lean{heightCharFirst}, \lean{heightCharTotal}).  On a word or multiset of height degree $(A,N)$
the torus labels are $2^A$ and $2^N$; since $2^a=2^b$ forces $a=b$, the property ``labels $=(2^A,2^N)$'' selects exactly the words,
resp.\ multisets, of degree $(A,N)$.  The transferred weighted count with $j_0=\sigma$ is the identity of (a).  (b) By (a) and
Theorem~\ref{ch19:thm:characters}(a), the word series $W$ satisfies $W(1-\ell)=1$, and the letter series $\ell$ of $U$ is
$\tilde u$: the letters of degree $(1+k,1+j)$ are the $(\mathrm{ad}\,e)^kb_{(1,j),i}$, with total weight $\sum_i\mu_{(1,j),i}=c(1,j)$, and
$(a,a+b)=(1+k,1+j)$ gives $j=a+b-1$.  By Theorem~\ref{ch19:thm:characters}(b), the multiset series $M$ satisfies $DM=MP$, and the
power-sum series of the letters of $\chnineteenu^+$ is $\tilde P$: a letter in $\mathfrak g_{(m,n)}$ has height degree $(m,m+n)$, first
coordinate $m$, and $\sum_i\mu^k_{(m,n),i}=c(k,mn)$.  By (a), $W=M$.  Theorem~\ref{ch19:thm:characters}(c) gives
$D\tilde u=(1-\tilde u)\tilde P$.  (c) By Theorem~\ref{ch19:thm:eigen} each $\sigma|_{\mathfrak g_\gamma}$ has an eigenbasis (its $N$-th power is
the identity), and the power sums of eigenvalues are the traces of powers; apply (b). \qed
\begin{physics}
The free (tensor-algebra) description of the BPS Fock space and its oscillator (PBW) description must have the same
$g$-twisted partition function.  The two torus elements are the ``momentum charges'' $(m,m+n)$: comparing traces at fixed
charges is comparing twisted partition functions coefficient by coefficient.
\end{physics}
\begin{status}No hypotheses beyond those displayed; standard axioms (\leanpath{MoonshineCNAudit18.lean}).  Supporting
(\leanpath{MoonshineCNTwistedCharacters.lean}): \lean{borcherdsHeightDegree_injOn}, \lean{borcherdsStringDegree_cone},
\lean{sigmaGradedLetters}, \lean{sum_sigma_fiber}, \lean{heightCharFirst}, \lean{heightCharTotal} and their additivity and
cone lemmas, \lean{list_prod_two_pow}, \lean{multiset_prod_two_pow}, \lean{two_pow_inj}, \lean{GradedLetters.listDeg_fst}
(and \lean{_snd}, \lean{msetDeg_fst}, \lean{msetDeg_snd}), \lean{positiveLetters}, \lean{generatorLetters};
(\leanpath{MoonshineCNTwistedDenominator.lean}) \lean{weightAction_one}.\end{status}

\subsection{The Monster instance}

\begin{theorem}[Twisted denominator identities of the Monster Lie algebra]\label{ch19:thm:twistedmonster}
Assume \lean{MonsterLieSimpleRootGeneration}.  Then:
\begin{enumerate}[label=(\alph*)]
\item for every $\tau\in\Aut(\Vnat)$ and every $(m,n)$ with $m,n\ge1$, the trace of $\tau$ on the root space $\chnineteenm_{(m,n)}$ (Borcherds
  coordinates) is $c_\tau(mn)=\tr(\tau|\Vnat_{1+mn})$;
\item for every $\sigma\in\Aut(\Vnat)$ of finite order, the family $c(k,n)=c_{\sigma^k}(n)$ satisfies the twisted denominator identity.
\end{enumerate}
\end{theorem}
\begin{leancode}
theorem monsterLie_weightAction_trace (hgen : MonsterLieSimpleRootGeneration) (τ : AutVnat)
    (γ : ℤ × ℤ) (hγ : borcherdsPositiveCone γ) :
    LinearMap.trace ℂ ((monsterLieBorcherdsData hgen).weightSpace γ)
        ((monsterLieBorcherdsData hgen).weightAction (monsterLieBorcherdsAction hgen) τ γ) =
      PowerSeries.coeff (γ.1 * γ.2).toNat (mckayThompson τ)
theorem twistedDenominatorIdentity_of_simpleRootGeneration (hgen : MonsterLieSimpleRootGeneration)
    (σ : AutVnat) (hσ : IsOfFinOrder σ) : TwistedDenominatorIdentity (mckayThompsonFamily σ)
def mckayThompsonFamily (g : AutVnat) : ℕ → ℕ → ℂ :=
  fun k n => PowerSeries.coeff n (mckayThompson (g ^ k))
\end{leancode}
(\leanpath{VOA/Moonshine/MoonshineCNTwistedDenominator.lean}, \leanpath{VOA/Moonshine/MoonshineCNReplicationAnalytic.lean}.)
\begin{dictionary}
\lean{monsterLieBorcherdsData hgen} is the Borcherds datum on $\chnineteenm$ and \lean{monsterLieBorcherdsAction hgen} the action of
$\Aut(\Vnat)$ on it (Chapter~\ref{ch:18}); both need \lean{hgen} only because the structure \lean{BorcherdsData} has the field
\lean{generated}.  \lean{PowerSeries.coeff n (mckayThompson τ)} is $\tr(\tau|\Vnat_{n+1})$ (Chapter~\ref{ch:16}).  So (a) is the
equivariant no-ghost theorem read on traces, and \lean{mckayThompsonFamily σ k n} $=c_{\sigma^k}(n)$.  (b) is
\eqref{ch19:eq:TD} for $g=\sigma$, in height coordinates.

\textbf{Strength.}  (b) is stated for finite-order $\sigma$ only.  Every automorphism of $\Vnat$ has finite order
(\lean{M8.autVnat_isOfFinOrder}, Chapter~\ref{ch:20}), so this is not a restriction in the end, but no theorem of the library
states the twisted identity for $\sigma$ without the finite-order hypothesis or without \lean{hgen}; the hypothesis-free
version is the composition of (b) with \lean{M8.monsterLieSimpleRootGeneration} and \lean{M8.autVnat_isOfFinOrder}, which
the library performs only at the level of complete replicability (\lean{completelyReplicable_holds}, Chapter~\ref{ch:22}).
Neither the exponential form \eqref{ch19:eq:twisted-exp} nor the product form with roots of unity is a Lean statement.
\end{dictionary}
\paragraph{Proof on paper.}  (a) The root-space isomorphism $\chnineteenm_{(m,n)}\cong\Vnat_{1-m(-n)}=\Vnat_{1+mn}$ (repository momentum
$(m,-n)$, Chapter~\ref{ch:18}; the sign convention is recalled in the first warning of \S\ref{ch19:sec:traps}) is $\Aut(\Vnat)$-equivariant; conjugating by it identifies the traces.  (b)
Theorem~\ref{ch19:thm:twistedabstract}(c) with $c(k,\cdot)=c_{\sigma^k}$ and (a) applied to $\tau=\sigma^k$. \qed
\begin{physics}
\textbf{Twisted partition functions and orbifold sectors.}  $T_{\sigma}$ is the torus partition function with the symmetry
$\sigma$ inserted along the time circle.  The twisted denominator identity is Borcherds' observation that the
$\sigma$-twisted index of the space of physical string states has the same product structure as the untwisted one, the
$k$-th term of the plethystic logarithm involving $T_{\sigma^k}$: a string state wound $k$ times around the time circle
picks up $\sigma^k$.  Carnahan's generalized moonshine programme (2010--2017) extends this to the twisted-sector
Lie algebras attached to commuting pairs $(g,h)$; Paquette--Persson--Volpato realise the twisted identities as twisted BPS
indices of a heterotic string.
\end{physics}
\begin{status}
\textbf{Conditional on} \lean{hgen : MonsterLieSimpleRootGeneration} (proved in Chapter~\ref{ch:20}); (b) also needs
\lean{IsOfFinOrder σ}.  Standard axioms according to the build of \leanpath{MoonshineCNAudit18.lean} (2,268 reports over the audit's whole import chain, 0 \lean{sorryAx}; recorded in
Chapter~\ref{ch:15}, not re-run for this chapter).
\end{status}

%======================================================================
\section{Replication}\label{ch19:sec:replication}
%======================================================================

\subsection{Replication from the twisted identity: the idea}

Norton's replication formulas (1984) say that the Faber polynomials of $T_g$ are computed by Hecke sums whose $a$-th
block uses the replicate $T_{g^a}$:
\begin{equation}\label{ch19:eq:rep}
  F_{g,m}\bigl(T_g(\tau)\bigr)=\sum_{ad=m}\ \sum_{0\le b<d}T_{g^a}\Bigl(\frac{a\tau+b}d\Bigr),\qquad m\ge1 .
\end{equation}
Borcherds derived them from the twisted denominator identities; Koike had observed the link between replication and
Hecke operators.  The derivation is the twisted analogue of Theorem~\ref{ch19:thm:knzlog}: substitute $X=T_g(q)$ in the
Faber generating identity, use the factorisation $p(T_g(p)-T_g(q))=(1-p/q)(1-u_g)$, and replace $D\log(1-u_g)$ by $-P_g$
using \eqref{ch19:eq:TD}.  Comparing $p^m$-coefficients gives the \emph{formal} replication formula
\begin{equation}\label{ch19:eq:repformal}
  F_{g,m}(T_g)=q^{-m}+\sum_{r\ge1}\Bigl(\sum_{a\mid\gcd(m,r)}\frac ma\,c_{g^a}\Bigl(\frac{mr}{a^2}\Bigr)\Bigr)q^r ,
\end{equation}
and the $q$-expansion of the Hecke sum (the computation of Theorem~\ref{ch19:thm:qheckesum}, with replicates) shows that the
right side of \eqref{ch19:eq:rep} has the same expansion.  Since both sides converge for finite-order $g$, they agree as
functions.

\subsection{Formal replication}

\begin{definition}[The Faber data in height coordinates]\label{ch19:def:faberheight}
For $c:\N\to\C$: $\Phi_c(p,X)=1-Xp+\sum_{n\ge1}c(n)p^{n+1}\in\C[X][[p]]$; the series $\sum_mQ_m(X)p^m=-D\Phi_c/\Phi_c$; the ring map
$\C[X][[p]]\to\C((q))[[p]]$ sending $\sum_mG_m(X)p^m$ to $\sum_mq^mG_m(f_c)p^m$ (substitute $X\mapsto f_c$ and $p\mapsto pq$); and the
replication coefficient $\rho_c(m,r)=\sum_{a\mid m,\ a\mid r}\frac ma\,c\bigl(a,\frac ma\cdot\frac ra\bigr)$ for $c:\N\to\N\to\C$.
\end{definition}
\begin{leancode}
def faberGenerating (c : ℕ → ℂ) : PowerSeries ℂ[X] :=
  PowerSeries.mk fun n => if n = 0 then 1 else if n = 1 then -X else C (c (n - 1))
def faberLogDerivSeries (c : ℕ → ℂ) : PowerSeries ℂ[X] :=
  -eulerDeriv (faberGenerating c) * PowerSeries.invOfUnit (faberGenerating c) 1
def faberEvaluation (c : ℕ → ℂ) : PowerSeries ℂ[X] →+* PowerSeries (LaurentSeries ℂ) :=
  (PowerSeries.rescale (HahnSeries.single (1 : ℤ) (1 : ℂ))).comp
    (PowerSeries.map (aeval (M8.qSeriesLaurent c)).toRingHom)
def replicationCoeff (c : ℕ → ℕ → ℂ) (m r : ℕ) : ℂ :=
  ∑ a ∈ m.divisors, if a ∣ r then ((m / a : ℕ) : ℂ) * c a (m / a * (r / a)) else 0
\end{leancode}
(\leanpath{VOA/Moonshine/MoonshineCNReplicationFormal.lean}.)
\begin{dictionary}
\lean{faberGenerating} is the same series as \lean{M8.faberGen} of \S\ref{ch19:sec:faber} (a duplicate written on the other
side of the project).  \lean{PowerSeries.invOfUnit F 1} is the inverse of a series with
constant term $1$; so \lean{faberLogDerivSeries} is $-D\Phi/\Phi$.  \lean{faberEvaluation c} is the composite of
\lean{PowerSeries.map} along $X\mapsto f_c$ and \lean{PowerSeries.rescale (single 1 1)}, the substitution $p\mapsto qp$.
\lean{replicationCoeff c m r} sums over $a\mid m$ with $a\mid r$, the term $\frac ma\,c(a,\frac ma\cdot\frac ra)$; since
$\frac ma\cdot\frac ra=\frac{mr}{a^2}$ this is the coefficient of \eqref{ch19:eq:repformal} with $c(a,\cdot)=c_{g^a}$.
\end{dictionary}
\begin{status}Definitions; standard axioms (\leanpath{MoonshineCNAudit18.lean}).\end{status}

\begin{physics}
Faber polynomials express the partition functions of the $n$-fold covers of the torus in terms of the partition
function itself. Height coordinates organise the root lattice of the Monster Lie algebra by the energy level of the
corresponding string states. They are the natural variables for comparing the two sides of the denominator identity
level by level.
\end{physics}

\begin{theorem}[Formal replication]\label{ch19:thm:repformal}
Let $c:\N\to\N\to\C$ satisfy the twisted denominator identity, and let $F_m$ be the Faber polynomial of
$f=q^{-1}+\sum_{n\ge1}c(1,n)q^n$.  For every $m\ge1$, in $\C((q))$ and in $\C[[q]]$ respectively,
\begin{equation*}
  q^m\,F_m(f)=1+[p^m]\tilde P,\qquad\text{i.e.}\qquad F_m(f)=q^{-m}+\sum_{r\ge1}\rho_c(m,r)\,q^r .
\end{equation*}
\end{theorem}
\begin{leancode}
theorem faberEvaluation_faberGenerating (c : ℕ → ℂ) :
    faberEvaluation c (faberGenerating c) =
      laurentEmbedding ((1 - PowerSeries.X) * (1 - twistedUnitSeries c))
theorem faberLogDerivSeries_coeff_eq_faberPoly (hTD : TwistedDenominatorIdentity c) (m : ℕ)
    (hm : 0 < m) :
    PowerSeries.coeff m (faberLogDerivSeries (c 1)) = M8.faberPoly (c 1) m
theorem faberPoly_replication (hTD : TwistedDenominatorIdentity c) (m : ℕ) (hm : 0 < m) :
    HahnSeries.single (m : ℤ) (1 : ℂ) * aeval (M8.qSeriesLaurent (c 1)) (M8.faberPoly (c 1) m) =
      HahnSeries.ofPowerSeries ℤ ℂ (1 + PowerSeries.coeff m (twistedPowerSum c))
def homogenizedEval (c : ℕ → ℂ) (m : ℕ) (P : ℂ[X]) : PowerSeries ℂ :=
  ∑ j ∈ Finset.range (m + 1),
    PowerSeries.C (P.coeff j) * PowerSeries.X ^ (m - j) * M8.qSeriesUnit c ^ j
theorem homogenizedEval_faberPoly (c : ℕ → ℕ → ℂ) (hTD : TwistedDenominatorIdentity c) (m : ℕ)
    (hm : 0 < m) :
    homogenizedEval (c 1) m (M8.faberPoly (c 1) m) =
      1 + PowerSeries.coeff m (twistedPowerSum c)
theorem coeff_coeff_twistedPowerSum_add (c : ℕ → ℕ → ℂ) (m r : ℕ) (hr : 0 < r) :
    PowerSeries.coeff (m + r) (PowerSeries.coeff m (twistedPowerSum c)) =
      replicationCoeff c m r
\end{leancode}
(\leanpath{VOA/Moonshine/MoonshineCNReplicationFormal.lean}; the theorems of the section \lean{Replication} have the section
variable \lean{(c : ℕ → ℕ → ℂ)}.)
\begin{dictionary}
\lean{laurentEmbedding} is the inclusion $\C[[q]][[p]]\subset\C((q))[[p]]$.  \lean{HahnSeries.single (m : ℤ) 1} is $q^m$.
\lean{M8.faberPoly (c 1) m} is the Faber polynomial of $f$ (Definition~\ref{ch19:def:faber}).
\lean{homogenizedEval c m P} $=\sum_{j\le m}P_jq^{m-j}u_c^{\,j}$ is $q^mP(f)$ as an honest power series when $\deg P\le m$
(\lean{ofPowerSeries_homogenizedEval}).  \lean{PowerSeries.coeff m (twistedPowerSum c)} is $[p^m]\tilde P\in\C[[q]]$; its
coefficient of $q^{m+r}$ is $\rho_c(m,r)$ (\lean{coeff_coeff_twistedPowerSum_add}) and its coefficients of $q^N$, $N\le m$, vanish
(\lean{coeff_coeff_twistedPowerSum_of_le}).  So the theorem is exactly \eqref{ch19:eq:repformal}, for an abstract family.

\textbf{Hypotheses.}  Only the twisted denominator identity; in particular the theorem holds for \emph{every} $g\in\Aut(\Vnat)$
whose family satisfies it, regardless of finite order (which matters only for convergence, below).
\end{dictionary}
\paragraph{Proof on paper.}  Under the substitution $X\mapsto f(q)$, $p\mapsto pq$, the series $\Phi$ becomes
$1-qf(q)\,p+\sum_{n\ge1}c(1,n)q^{n+1}p^{n+1}$, which is $(1-p)(1-\tilde u)$ (compare coefficients:
\lean{coeff_one_sub_X_mul_one_sub_twistedUnitSeries}).  The substitution is a ring map commuting with $D$
(\lean{faberEvaluation_eulerDeriv}), so it carries $\Phi\cdot\sum_mQ_mp^m=-D\Phi$ to
$(1-p)(1-\tilde u)\cdot\mathcal Q=-D\bigl((1-p)(1-\tilde u)\bigr)$ with $\mathcal Q=\sum_mq^mQ_m(f)p^m$.  The right side is
$p(1-\tilde u)+(1-p)D\tilde u=(1-p)(1-\tilde u)\bigl(\frac p{1-p}+\tilde P\bigr)$ by the twisted identity
(\lean{twistedFaberRecursion}).  Cancelling the unit $(1-p)(1-\tilde u)$ (\lean{laurentEmbedding_injective}),
$\mathcal Q=\sum_{k\ge1}p^k+\tilde P$ (\lean{faberEvaluation_faberLogDerivSeries}), i.e.\ $q^mQ_m(f)=1+[p^m]\tilde P$
(\lean{faberLogDerivSeries_coeff_eval}).  Since $[p^m]\tilde P$ only has powers $q^N$ with $N>m$, $Q_m(f)=q^{-m}+O(q)$, and
$Q_m$ is monic of degree $m$; by uniqueness of Faber polynomials $Q_m=F_m$ (\lean{faberLogDerivSeries_coeff_eq_faberPoly}).
\qed
\begin{physics}
The formal replication formula is the $g$-twisted version of ``the partition function of $m$-fold covers is a polynomial
in the single partition function'': in the $a$-th block the covering map wraps the time circle $a$ times, so the twist
along it becomes $g^a$.  This is the content of the Hecke operators on twisted sectors (Borcherds; Carnahan,
Generalized moonshine I and II), and of the DMVV formula for symmetric-product orbifolds with a symmetry inserted.
\end{physics}
\begin{status}Hypothesis \lean{TwistedDenominatorIdentity c} only; standard axioms (\leanpath{MoonshineCNAudit18.lean}).
Supporting: \lean{positiveGeometricSeries}, \lean{one_sub_X_mul_positiveGeometricSeries},
\lean{coeff_coeff_twistedPowerSum_of_le}, \lean{constantCoeff_faberGenerating}, \lean{faberGenerating_mul_faberLogDerivSeries},
\lean{coeff_faberEvaluation}, \lean{faberEvaluation_eulerDeriv}, \lean{laurentEmbedding}, \lean{laurentEmbedding_injective},
\lean{coeff_one_sub_X_mul_one_sub_twistedUnitSeries}, \lean{algebraMap_laurentSeries}, \lean{twistedFaberRecursion},
\lean{faberEvaluation_faberLogDerivSeries}, \lean{faberLogDerivSeries_coeff_eval}, \lean{ofPowerSeries_homogenizedEval}.\end{status}

\subsection{Analytic replication}

\begin{theorem}[Analytic replication]\label{ch19:thm:repanalytic}
\begin{enumerate}[label=(\alph*)]
\item If $h\in\Aut(\Vnat)$ has finite order, its McKay--Thompson coefficients are $J$-dominated.
\item (Hecke side) If $c(a,\cdot)$ is $J$-dominated with $c(a,0)=0$ for every $a\ge0$ (including $a=0$, as in the Lean
  quantifier \lean{∀ a}), then for $m\ge1$ and $\tau\in\chnineteenH$,
  $q^m\sum_{a\mid m}\sum_{b<m/a}T_a\bigl(\frac{a\tau+b}{m/a}\bigr)-1=\sum_N[p^mq^N]\tilde P\;q^N$, where
  $T_a=q^{-1}+\sum_{n\ge0}c(a,n)q^n$.
\item (Faber side) If $c$ satisfies the twisted denominator identity and $c(1,\cdot)$ is $J$-dominated with $c(1,0)=0$,
  then $q^mF_m(T_1(\tau))-1=\sum_N[p^mq^N]\tilde P\;q^N$.
\item (Replication) If $g\in\Aut(\Vnat)$ has finite order and the family $k\mapsto c_{g^k}$ satisfies the twisted denominator
  identity, then \eqref{ch19:eq:rep} holds for every $m\ge1$ and $\tau\in\chnineteenH$, with $F_{g,m}$ the Faber polynomial of
  $T_g$.
\end{enumerate}
\end{theorem}
\begin{leancode}
def poleQSeriesFun (c : ℕ → ℂ) (τ : ℍ) : ℂ :=
  (qPar τ)⁻¹ + ∑' n : ℕ, c n * qPar τ ^ n
theorem isJDominated_mckayThompson (h : AutVnat) (hh : IsOfFinOrder h) :
    IsJDominated (fun n => PowerSeries.coeff n (mckayThompson h))
theorem hasSum_heckeSum (c : ℕ → ℕ → ℂ) (hc : ∀ a, IsJDominated (c a)) (hc0 : ∀ a, c a 0 = 0)
    (m : ℕ) (hm : 0 < m) (τ : ℍ) :
    HasSum (fun N => PowerSeries.coeff N (PowerSeries.coeff m (twistedPowerSum c)) * qPar τ ^ N)
      (qPar τ ^ m * (∑ a ∈ m.divisors.attach, ∑ b ∈ Finset.range (m / a),
        poleQSeriesFun (c a) (heckePoint a (m / a) b (Nat.pos_of_mem_divisors a.2)
          (Nat.div_pos (Nat.divisor_le a.2) (Nat.pos_of_mem_divisors a.2)) τ)) - 1)
theorem hasSum_faberPoly_eval (c : ℕ → ℕ → ℂ) (hTD : TwistedDenominatorIdentity c)
    (hc : IsJDominated (c 1)) (hc0 : c 1 0 = 0) (m : ℕ) (hm : 0 < m) (τ : ℍ) :
    HasSum (fun N => PowerSeries.coeff N (PowerSeries.coeff m (twistedPowerSum c)) * qPar τ ^ N)
      (qPar τ ^ m * (M8.faberPoly (c 1) m).eval (poleQSeriesFun (c 1) τ) - 1)
theorem faberPoly_eval_mckayThompsonFun (g : AutVnat) (hg : IsOfFinOrder g)
    (hTD : TwistedDenominatorIdentity (mckayThompsonFamily g)) (m : ℕ) (hm : 0 < m) (τ : ℍ) :
    (M8.faberPoly (mckayThompsonFamily g 1) m).eval (mckayThompsonFun g τ) =
      ∑ a ∈ m.divisors.attach, ∑ b ∈ Finset.range (m / a),
        mckayThompsonFun (g ^ (a : ℕ)) (heckePoint a (m / a) b (Nat.pos_of_mem_divisors a.2)
          (Nat.div_pos (Nat.divisor_le a.2) (Nat.pos_of_mem_divisors a.2)) τ)
\end{leancode}
(\leanpath{VOA/Moonshine/MoonshineCNReplicationAnalytic.lean}.)
\begin{dictionary}
\lean{poleQSeriesFun c τ} $=q^{-1}+\sum_{n\ge0}c(n)q^n$ (it \emph{includes} $c(0)$; that is why $c(a,0)=0$ is a hypothesis);
\lean{mckayThompsonFun_eq_poleQSeriesFun} identifies \lean{mckayThompsonFun g} of Chapter~\ref{ch:16} with it.  The Hecke sum in
(b) and (d) is written with the same \lean{m.divisors.attach}/\lean{Finset.range (m / a)}/\lean{heckePoint} bookkeeping as
\lean{CompletelyReplicable} (Chapter~\ref{ch:16}), so (d) is literally the clause of \lean{CompletelyReplicable} for $g$ and $m$,
with $P$ the Faber polynomial \lean{M8.faberPoly (mckayThompsonFamily g 1) m}.  \lean{g ^ (a : ℕ)} is the replicate.
\end{dictionary}
\paragraph{Proof on paper.}  (a) Theorem of Chapter~\ref{ch:16} (trace bounds): for finite order the eigenvalues are roots of
unity, so $|\tr(h|\Vnat_n)|\le\dim\Vnat_n$.  (b) As in Theorem~\ref{ch19:thm:qheckesum}, block by block, with the character sum
\lean{sum_range_exp_pow} and the pole computation \lean{sum_inv_qPar_heckePoint}; the coefficient of $q^{m+r}$ after multiplying by
$q^m$ is $\rho_c(m,r)$, which is $[p^mq^{m+r}]\tilde P$.  (c) $q^mF_m(T_1)$ is the value at $q$ of the power series
$\mathtt{homogenizedEval}=1+[p^m]\tilde P$ (Theorem~\ref{ch19:thm:repformal}), and evaluation of absolutely convergent power
series is a ring map (\lean{HasAbsEvalAt}, \lean{hasAbsEvalAt_homogenizedEval}).  (d) By (a) all $c_{g^a}$ are dominated, so (b)
and (c) apply; the two sides of \eqref{ch19:eq:rep}, multiplied by $q^m\ne0$, have the same convergent expansion. \qed
\begin{physics}
Convergence is where finite order enters: a symmetry of finite order is unitary on each energy level, so its twisted
partition function is dominated by the untwisted one and converges on the whole upper half-plane.
\end{physics}
\begin{status}Hypotheses as displayed; standard axioms (\leanpath{MoonshineCNAudit18.lean}).  Supporting:
\lean{HasAbsEvalAt} and its closure lemmas (\lean{HasAbsEvalAt.add}, \lean{.mul}, \lean{.const}, \lean{.one}, \lean{.zero},
\lean{.X}, \lean{.pow}, \lean{.sum}), \lean{qPar_ne_zero}, \lean{summable_norm_coeff_mul_qPar}, \lean{hasAbsEvalAt_qSeriesUnit},
\lean{hasAbsEvalAt_homogenizedEval}, \lean{qPar_heckePoint_eq_exp_mul}, \lean{exp_heckeShift_pow}, \lean{sum_range_exp_pow},
\lean{hasSum_heckeCosetSum_regular}, \lean{sum_inv_qPar_heckePoint}, \lean{mckayThompsonFamily_one}.  (The lemmas
\lean{hasSum_heckeSum} and \lean{hasAbsEvalAt_homogenizedEval} duplicate the lemmas \lean{hasSum_qHeckeSum} and
\lean{hasSum_mul_eval_qSeriesFun}; nothing depends on the choice.)\end{status}

\subsection{Complete replicability}

\begin{theorem}[Complete replicability of the McKay--Thompson series]\label{ch19:thm:completerep}
\begin{enumerate}[label=(\alph*)]
\item If every finite-order $g\in\Aut(\Vnat)$ has a family satisfying the twisted denominator identity, then
  \lean{CompletelyReplicableOfFinOrder} holds; if moreover every $g$ has finite order and every family satisfies the identity,
  \lean{CompletelyReplicable} holds.
\item Given \lean{MonsterLieSimpleRootGeneration}: \lean{CompletelyReplicableOfFinOrder}; and, given also that every
  automorphism has finite order, \lean{CompletelyReplicable} and the Conway--Norton conjecture.
\end{enumerate}
\end{theorem}
\begin{leancode}
theorem completelyReplicableOfFinOrder_of_twistedDenominator
    (hTD : ∀ g : AutVnat, IsOfFinOrder g → TwistedDenominatorIdentity (mckayThompsonFamily g)) :
    CompletelyReplicableOfFinOrder
theorem completelyReplicable_of_twistedDenominator (hfin : ∀ g : AutVnat, IsOfFinOrder g)
    (hTD : ∀ g : AutVnat, TwistedDenominatorIdentity (mckayThompsonFamily g)) :
    CompletelyReplicable
theorem completelyReplicableOfFinOrder_of_simpleRootGeneration
    (hgen : MonsterLieSimpleRootGeneration) : CompletelyReplicableOfFinOrder
theorem completelyReplicable_of_simpleRootGeneration (hgen : MonsterLieSimpleRootGeneration)
    (hfin : ∀ g : AutVnat, IsOfFinOrder g) : CompletelyReplicable
theorem conwayNorton_of_simpleRootGeneration (hgen : MonsterLieSimpleRootGeneration)
    (hfin : ∀ g : AutVnat, IsOfFinOrder g) : ConwayNortonConjecture
\end{leancode}
(\leanpath{VOA/Moonshine/MoonshineCNReplicationAnalytic.lean}, \leanpath{VOA/Moonshine/MoonshineCNTwistedDenominator.lean}.)
\begin{dictionary}
\lean{CompletelyReplicable} is the fixed definition of the statement file (Chapter~\ref{ch:16}): for all $g$ and $m\ge1$ a monic
degree-$m$ polynomial $P$ with $P(T_g(\tau))=\sum_{a\mid m}\sum_{b<m/a}T_{g^a}((a\tau+b)/(m/a))$; \lean{CompletelyReplicableOfFinOrder}
(file \leanpath{MoonshineCNGammaZeroBridge.lean}, Chapter~\ref{ch:16}) is the same restricted to finite-order $g$.  The
witness $P$ is the Faber polynomial.  The last theorem composes with the genus-zero theorem through
\lean{M8.conwayNorton_of_finOrder_of_completelyReplicable} (file \leanpath{MoonshineLRConwayNorton.lean}, Chapter~\ref{ch:20}).
\end{dictionary}
\paragraph{Proof on paper.}  (a) Theorem~\ref{ch19:thm:repanalytic}(d) with $P=F_{g,m}$, monic of degree $m$ by
Theorem~\ref{ch19:thm:faber}.  (b) Theorem~\ref{ch19:thm:twistedmonster}(b) gives the twisted identities, then (a); the
Conway--Norton conjecture follows from the genus-zero theorem of this chapter and the bridge of Chapter~\ref{ch:16}. \qed
\begin{physics}Every twisted partition function of the moonshine CFT closes, with its powers, under the Hecke covering
operations.\end{physics}
\begin{status}
(a): hypotheses as displayed.  (b): \textbf{conditional on} \lean{hgen}; the \lean{CompletelyReplicable} and Conway--Norton
statements also on \lean{hfin}.  Both are theorems of Chapter~\ref{ch:20}, and \lean{completelyReplicable_holds} (no hypotheses)
is recorded in Chapter~\ref{ch:22}.  Standard axioms (\leanpath{MoonshineCNAudit18.lean}).
\end{status}

%======================================================================
\section{Genus zero: the theorem and the route}\label{ch19:sec:gz}
%======================================================================

\subsection{What is proved}

The genus-zero theorem is proved for the abstract interface \lean{ReplicableFamily N} of Chapter~\ref{ch:16}: a doubly
indexed sequence $c(k,n)$, periodic of period $N$ in the replicate index $k$, with $|c(k,n)|\le|c(0,n)|$, $c(0,\cdot)$ the
coefficients of $J$, and satisfying complete replication with replicate $ka$ in the $a$-th Hecke block.  Write
$f_k=\mathtt{qSeriesFun}\,(c(k,\cdot))$ and, for any $f:\chnineteenH\to\C$,
\begin{equation*}
  G(f)=\{\gamma\in\chnineteenGLR:\ \det\gamma=1,\ f(\gamma\tau)=f(\tau)\ \forall\tau\in\chnineteenH\}.
\end{equation*}

\begin{quote}\itshape
\textbf{Genus-zero theorem (the Lean statement).}  Let $N\ge1$ and let $F$ be a replicable family of period $N$ that satisfies
the prime Hecke identity for every member and of which no member is injective modulo translations.  For every $k$ there are a
subgroup $\Gamma\le\chnineteenGLR$ and $N'\ge1$ such that $\Gamma\le\chnineteenSLR$, $\Gamma_1(N')\le\Gamma$, $\Gamma$ is commensurable with
$\chnineteenSL$, and $f_k$ is a normalised Hauptmodul of $\Gamma$ (in the sense of \lean{IsNormalizedHauptmodul}).
\end{quote}
This is \lean{ReplicableFamily.exists_hauptmodul_gammaOne} (\S\ref{ch19:sec:main}).  Three remarks keep the reader honest about
what the displayed Lean type says.  (i) The two hypotheses, the prime Hecke identity and nontriviality, are explicit premises
(\lean{hB1}, \lean{hD7}) of that theorem; they are proved for every family by two \emph{separate} theorems
(\S\ref{ch19:sec:lrinputs}), and no Lean declaration packages the composite ``every replicable family of positive period''.
(ii) The proof takes $\Gamma=G(f_k)$, and by the field \lean{injective} this is the only possible choice up to $\pm1$; but the
identification $\Gamma=G(f_k)$ is in the proof and the docstring, not in the type.  (iii) The proof also gets $N\mid N'$, from
\lean{ReplicableFamily.exists_gamma1GL_le_symmetryGroup} (Theorem~\ref{ch19:thm:level}); the main theorem discards this fact.  The McKay--Thompson family of a finite-order $g$ is an instance (Chapter~\ref{ch:16}), so
complete replicability gives the $\Gamma_1$ form of the Conway--Norton conjecture.

\subsection{The route, on one page}\label{ch19:sec:route}

Fix $F$, $k$, $f=f_k$, $G=G(f)$.  The proof has six stages; the sections that follow treat them in order.
\begin{enumerate}[label=(\arabic*)]
\item \emph{Analytic basics} (\S\ref{ch19:sec:analytic}).  $f$ is holomorphic, $f-q^{-1}\to0$ at $i\infty$, and there is a height
  $t$ depending only on $J$ above which $f$ is injective modulo $\Z$ (\emph{top injectivity}); every value of large modulus is
  taken high up.
\item \emph{Modular equations at primes, as a pairing} (\S\ref{ch19:sec:lrinputs}, \S\ref{ch19:sec:pairing}).  For a prime
  $p\equiv1\pmod N$, the Hecke operator at $p$ maps the Faber polynomials of $f$ to polynomials in $f$; hence for every
  polynomial $P$, $\sum_{\beta\in\Delta(p)}P(f(\beta\tau))$ is a polynomial in $f(\tau)$, with $\Delta(p)=\{\tau\mapsto p\tau\}\cup
  \{\tau\mapsto(\tau+j)/p\}$.  So the multiset $\{f(\beta\tau)\}_\beta$ depends only on $f(\tau)$: if $f(z)=f(w)$, every value
  $f(\beta z)$ is some $f(\beta'w)$ (the \emph{Hecke pairing}), and on a connected domain one partner works throughout.
\item \emph{Local symmetries are real M\"obius maps} (\S\ref{ch19:sec:local}).  If $\alpha$ is holomorphic on a disc with
  $f\circ\alpha=f$ and $\alpha'\ne0$, then $\alpha$ is the restriction of an element of $G$.  Pairing at two different primes
  produces real affine maps $C_n\to\mathrm{id}$ with $\alpha\circ C_n=C'_n\circ\alpha$; differentiating gives a \emph{real linear ODE}
  $(\lambda z+\mu)\alpha'=\lambda'\alpha+\mu'$; its solutions are four explicit families, each holomorphic on all of $\chnineteenH$; a
  boundary argument shows the explicit solution maps $\chnineteenH$ into $\chnineteenH$ and intertwines $f$ globally; growth and
  injectivity of $f$ high up leave only real M\"obius maps.
\item \emph{Property P} (\S\ref{ch19:sec:local}).  If $f(z)=f(w)$ then $w=\gamma z$ for some $\gamma\in G$: from (3) when $f'(w)\ne0$,
  and by approximation and closedness of $G$ when $f'(z)=f'(w)=0$.
\item \emph{The symmetry group} (\S\ref{ch19:sec:group}).  Shimizu's lemma: every $\gamma\in G$ with lower-left entry $c\ne0$ has
  $|c|\ge1$; the stabiliser of $\infty$ is $\pm\langle T\rangle$; bounded subsets of $G$ are finite.  Pairing on all of $\chnineteenH$
  gives a \emph{Hecke step} $\chnineteenmat abcd\mapsto\chnineteenmat{pa}{b-ja}c{(d-jc)/p}$ inside $G$.  Nontriviality ($f$ is not injective
  modulo $\Z$) and property P give an element with $c\ne0$; rationality, integrality and Dirichlet's theorem on primes in
  progressions then give $\Gamma_1(N')\le G$.
\item \emph{Genus zero} (\S\ref{ch19:sec:main}).  Cusp separation by Shimizu's lemma for conjugates; boundedness and values
  of $f$ at the other cusps; commensurability with $\chnineteenSL$ by Shimura's finiteness lemma on a fundamental domain;
  surjectivity onto $\C\cup\{$cusp values$\}$ by an open-and-closed argument; and the assembly of the eight fields of
  \lean{IsNormalizedHauptmodul}.
\end{enumerate}
The comparison with the literature is in \S\ref{ch19:sec:compare}.

\paragraph{Notation in Part IV.}  Several letters were used with other meanings in Parts I--III.  From here on they mean:
\begin{center}\small
\begin{tabular}{p{0.12\textwidth}p{0.78\textwidth}}
\toprule
symbol & meaning in \S\S\ref{ch19:sec:gz}--\ref{ch19:sec:main}\\
\midrule
$\mathfrak c$ & a coefficient sequence $\N\to\C$, with $f=\mathtt{qSeriesFun}\ \mathfrak c$ (Lean: \lean{c}); in Parts I--III it was $c$\\
$a,b,c,d$ & the entries of a real $2\times2$ matrix $\chnineteenmat abcd$ (Lean: \lean{g 0 0}, \lean{g 0 1}, \lean{g 1 0}, \lean{g 1 1}); $c$ is
  always the lower-left entry\\
$p,q$ & $p$ an admissible prime, $q$ either $e^{2\pi i\tau}$ or (in the overlap construction and the level theorem) a second prime;
  the formal variable $p$ of Parts I--III does not occur\\
$G$ & the symmetry group $G(f)$ (not the Faber series $G(p)$ of \S\ref{ch19:sec:faber}, nor the group acting on a Borcherds datum)\\
$H_m$ & the replicate Hecke sum (not the Faber series $H(p)$); half-discs are written $\Omega^+$\\
$\Omega$ & a connected open set (domain); $D(z_0,\eta)$ is the open disc of centre $z_0$ and radius $\eta$\\
$\Delta(p)$ & the $p+1$ maps $\tau\mapsto p\tau$, $\tau\mapsto(\tau+j)/p$; an element is written $\beta$ or $\beta_0$\\
$\delta$ & in \S\ref{ch19:sec:main}, a matrix with $\delta\infty$ a cusp; in \S\ref{ch19:sec:local}, a positive real dilation\\
$T$ & the translation $\chnineteenmat1101$\\
\bottomrule
\end{tabular}
\end{center}

\begin{physics}
\textbf{Hauptmoduln: ``the partition function determines the theory''.}  The genus-zero property says that a twisted
partition function $T_g$ is a global coordinate on the moduli space of tori with a $g$-twist (modulo the symmetries that
preserve the twist).  Duncan and Frenkel (2011) showed that the genus-zero property is equivalent to the convergence of the
Rademacher sum of $q^{-1}$ to $T_g$, i.e.\ to $T_g$ being determined by its polar part in the way a semiclassical
three-dimensional gravity partition function would be (Witten 2007: pure gravity in AdS$_3$ is conjecturally dual to the
extremal $c=24$ theory, whose partition function is $J$; its sum over geometries is a Rademacher-type sum).  These are
heuristics that motivate the theorem; the proof here uses none of them.
\end{physics}

\subsection{The shared interface}

The two halves of the genus-zero proof (Fedor Popov's side: stages (1)--(6); Leonardo Rastelli's side: the prime Hecke
identity, the unboundedness of $J$'s coefficients and nontriviality) communicate through one file.

\begin{definition}[Replicate Hecke sums; prime Hecke identity; injectivity modulo translations; large values; symmetry
group]\label{ch19:def:interface}
For a replicable family $F$ of period $N$, $k\in\N$, $m\ge1$:
\begin{enumerate}[label=(\alph*)]
\item $H_m(\tau)=\sum_{a\mid m}\sum_{0\le b<m/a}f_{ka}\bigl(\frac{a\tau+b}{m/a}\bigr)$, the right side of the replication formula;
\item the \emph{prime Hecke identity} at $p$: for all $m\ge1$, $\tau$:
  $\sum_{j<p}H_m\bigl(\frac{\tau+j}p\bigr)+H_m(p\tau)=H_{mp}(\tau)+[p\mid m]\,p\,H_{m/p}(\tau)$; the family satisfies it for member $k$
  if it holds at every prime $p\equiv1\pmod N$;
\item $f$ is \emph{injective modulo translations} if $f(\sigma)=f(\tau)$ implies $\sigma-\tau\in\Z$; the family is
  \emph{nontrivial} if no member is;
\item $f$ \emph{takes its large values from the top} if there are $R$ and $Y>0$ with every $w$, $|w|\ge R$, equal to $f(\sigma)$ for
  some $\sigma$ with $\operatorname{Im}\sigma\ge Y$;
\item $G(f)$ as displayed above.
\end{enumerate}
\end{definition}
\begin{leancode}
noncomputable def heckeOrbitSum {N : ℕ} (F : ReplicableFamily N) (k m : ℕ) (τ : ℍ) : ℂ :=
  ∑ a ∈ m.divisors.attach, ∑ b ∈ Finset.range (m / a),
    qSeriesFun (F.c (k * a)) (heckePoint a (m / a) b (Nat.pos_of_mem_divisors a.2)
      (Nat.div_pos (Nat.divisor_le a.2) (Nat.pos_of_mem_divisors a.2)) τ)
def PrimeHeckeIdentityAt {N : ℕ} (F : ReplicableFamily N) (k p : ℕ) (hp : 0 < p) : Prop :=
  ∀ m : ℕ, 0 < m → ∀ τ : ℍ,
    (∑ j ∈ Finset.range p, heckeOrbitSum F k m (heckePoint 1 p j one_pos hp τ)) +
        heckeOrbitSum F k m (heckePoint p 1 0 hp one_pos τ) =
      heckeOrbitSum F k (m * p) τ +
        if p ∣ m then (p : ℂ) * heckeOrbitSum F k (m / p) τ else 0
def PrimeHeckeIdentity {N : ℕ} (F : ReplicableFamily N) (k : ℕ) : Prop :=
  ∀ p : ℕ, ∀ hp : p.Prime, p ≡ 1 [MOD N] → PrimeHeckeIdentityAt F k p hp.pos
def IsInjectiveModTranslations (f : ℍ → ℂ) : Prop :=
  ∀ σ τ : ℍ, f σ = f τ → ∃ n : ℤ, (σ : ℂ) = τ + n
def ReplicableFamily.NotInjectiveModTranslations {N : ℕ} (F : ReplicableFamily N) : Prop :=
  ∀ k : ℕ, ¬ IsInjectiveModTranslations (qSeriesFun (F.c k))
def LargeValuesFromTop (f : ℍ → ℂ) : Prop :=
  ∃ R Y : ℝ, 0 < Y ∧ ∀ w : ℂ, R ≤ ‖w‖ → ∃ σ : ℍ, Y ≤ σ.im ∧ f σ = w
def symmetryGroup (f : ℍ → ℂ) : Subgroup (GL (Fin 2) ℝ) where
  carrier := {g | g.det = 1 ∧ ∀ τ : ℍ, f (g • τ) = f τ}
  ...
lemma mem_symmetryGroup {f : ℍ → ℂ} {g : GL (Fin 2) ℝ} :
    g ∈ symmetryGroup f ↔ g.det = 1 ∧ ∀ τ : ℍ, f (g • τ) = f τ
instance hasDetOne_symmetryGroup (f : ℍ → ℂ) : (symmetryGroup f).HasDetOne
\end{leancode}
(\leanpath{VOA/Moonshine/MoonshineCNGenusZeroInterface.lean}, namespace \lean{VOA.Moonshine.ConwayNorton}; the elided fields of
\lean{symmetryGroup} are the proofs that the carrier is closed under products, contains $1$ and inverses.)
\begin{dictionary}
\textbf{(a)} The sum is literally the right side of the field \lean{replicable} of \lean{ReplicableFamily}, so
\lean{ReplicableFamily.exists_monic_eval_eq_heckeOrbitSum} (in the same file) is a restatement of that field: some monic degree-$m$
polynomial in $f_k$ equals $H_m$.  For $m=0$ the sum is empty and $H_0=0$.  \textbf{(b)} \lean{heckePoint 1 p j} is
$(\tau+j)/p$ and \lean{heckePoint p 1 0} is $p\tau$; \lean{p ≡ 1 [MOD N]} is $p\equiv1\pmod N$; \lean{hp.pos} supplies $p>0$.  The
identity is the Hecke-algebra relation $T_mT_p=T_{mp}+[p\mid m]\,p\,T_{m/p}$ (for the unnormalised operators) applied to the
replicate Hecke sums; it is meaningful only for $p\equiv1\pmod N$, because the replicate index of the $a$-th block changes from
$ka$ to $kap$, and $f_{kap}=f_{ka}$ needs $p\equiv1$.  \textbf{(c)} \lean{(σ : ℂ) = τ + n} with $n\in\Z$.  \textbf{(d)} as
stated.  \textbf{(e)} \lean{g.det} is a unit of $\chnineteenR$; \lean{g.det = 1} is $\det g=1$.  Because $-1$ acts trivially on
$\chnineteenH$, $-1\in G(f)$ always (\lean{neg_one_mem_symmetryGroup}), so ``$\Gamma_1(N')\le G$'' and ``$\pm\Gamma_1(N')\le G$'' are the
same statement.  The instance \lean{hasDetOne_symmetryGroup} provides Mathlib's \lean{Subgroup.HasDetOne}, used in the statement of
the conjecture.

\textbf{Why $G(f)$ is the right group.}  It is the \emph{full} invariance group of $f$ in $\chnineteenSLR$ (as a subgroup of
$\chnineteenGLR$).  By the dictionary of the normalised-Hauptmodul definition in Chapter~\ref{ch:16}, a normalised Hauptmodul
determines its group up to $\pm1$; so the only possible choice of $\Gamma$ in the conclusion is $G(f)$, and the theorem
proves that this choice works.
\end{dictionary}
\begin{physics}
$G(f)$ is the group of modular transformations under which the twisted partition function is invariant: the stabiliser of
the twisted sector.  The prime Hecke identity says that the covering operations at two primes commute in the expected way:
a $p$-fold cover of an $m$-fold cover is an $mp$-fold cover, plus a correction from the covers that factor through the
multiplication by $p$.
\end{physics}
\begin{status}Definitions; printed in \leanpath{VOA/Moonshine/MoonshineCNGenusZeroAudit.lean}, standard axioms.\end{status}

%======================================================================
\section{Genus zero I: analytic basics}\label{ch19:sec:analytic}
%======================================================================

\begin{theorem}[Holomorphy, normalisation, growth]\label{ch19:thm:analytic}
Let $\mathfrak c$ be $J$-dominated and $f=\mathtt{qSeriesFun}\ \mathfrak c$.  Then
\begin{enumerate}[label=(\alph*)]
\item $|f(\tau)-q^{-1}|\le\mathtt{jTail}(\operatorname{Im}\tau)$, and $\mathtt{jTail}(y)\to0$ as $y\to\infty$;
\item $f(\tau)-q^{-1}\to0$ as $\operatorname{Im}\tau\to\infty$ (uniformly in $\operatorname{Re}\tau$);
\item $f$ is holomorphic on $\chnineteenH$.
\end{enumerate}
Every member of a replicable family is $J$-dominated, and $c(k,\cdot)=c(k',\cdot)$ whenever $k\equiv k'\pmod N$.
\end{theorem}
\begin{leancode}
theorem norm_qSeriesFun_sub_inv_qPar_le {c : ℕ → ℂ} (hc : IsJDominated c) (τ : ℍ) :
    ‖qSeriesFun c τ - (qPar τ)⁻¹‖ ≤ jTail τ.im
lemma tendsto_jTail_atTop : Tendsto jTail atTop (𝓝 0)
theorem tendsto_qSeriesFun_sub_inv_qPar {c : ℕ → ℂ} (hc : IsJDominated c) :
    Tendsto (fun τ : ℍ => qSeriesFun c τ - (qPar τ)⁻¹) atImInfty (𝓝 0)
theorem mdifferentiable_qSeriesFun {c : ℕ → ℂ} (hc : IsJDominated c) :
    MDifferentiable 𝓘(ℂ) 𝓘(ℂ) (qSeriesFun c)
lemma isJDominated (k : ℕ) : IsJDominated (F.c k)
lemma c_eq_of_modEq {k k' : ℕ} (h : k ≡ k' [MOD N]) : F.c k = F.c k'
\end{leancode}
(\leanpath{VOA/Moonshine/MoonshineCNGenusZeroAnalytic.lean}; the last two are \lean{ReplicableFamily.isJDominated},
\lean{ReplicableFamily.c_eq_of_modEq} with \lean{F : ReplicableFamily N} a section variable.)
\begin{dictionary}
\lean{atImInfty} and \lean{MDifferentiable 𝓘(ℂ) 𝓘(ℂ)} are as in Chapter~\ref{ch:16}: (b) and (c) are exactly the fields
\lean{normalized} and \lean{holomorphic} of \lean{IsNormalizedHauptmodul}.  \lean{Nat.ModEq} (\lean{k ≡ k' [MOD N]}) is
congruence of natural numbers.
\end{dictionary}
\paragraph{Proof on paper.}  (a) $|\mathfrak c(n)q^n|\le c_J(n)e^{-2\pi n\operatorname{Im}\tau}$, summed.  $\mathtt{jTail}(y)\to0$ by dominated
convergence for series (each term decreases to $0$, and the sum converges at $y=1$).  (b) follows from (a).  (c) A power series
in $q$ converging on $|q|<1$ is holomorphic there, and $q=e^{2\pi i\tau}$ is holomorphic.  The family statements: $|c(k,n)|\le|c(0,n)|=
c_J(n)$; and periodicity in $k$. \qed
\begin{physics}At low temperature ($\operatorname{Im}\tau\to\infty$) the twisted partition function is dominated by the
vacuum, $q^{-1}$, with corrections bounded by those of the untwisted theory.\end{physics}
\begin{status}Hypothesis \lean{IsJDominated c}; standard axioms (\leanpath{MoonshineCNGenusZeroAudit.lean}).\end{status}

\begin{theorem}[Top injectivity]\label{ch19:thm:topinj}
There is $t>0$ such that for every $J$-dominated $\mathfrak c$ and all $\sigma,\tau\in\chnineteenH$ with $\operatorname{Im}\sigma,
\operatorname{Im}\tau\ge t$: if $f(\sigma)=f(\tau)$ then $\sigma-\tau\in\Z$.
\end{theorem}
\begin{leancode}
theorem exists_height_injective :
    ∃ t : ℝ, 0 < t ∧ ∀ c : ℕ → ℂ, IsJDominated c → ∀ σ τ : ℍ, t ≤ σ.im → t ≤ τ.im →
      qSeriesFun c σ = qSeriesFun c τ → ∃ k : ℤ, (σ : ℂ) = τ + k
\end{leancode}
(\leanpath{VOA/Moonshine/MoonshineCNGenusZeroAnalytic.lean}.)
\begin{dictionary}
The quantifier order is the strong one: one height $t$ serves for \emph{all} $J$-dominated series.  The library names a
choice of it \lean{topHeight} (file \leanpath{MoonshineCNGenusZeroPartner.lean}), and ``the top'' means $\{\operatorname{Im}\ge t\}$.
\end{dictionary}
\paragraph{Proof on paper.}  Write $f=q^{-1}+h(q)$ with $h(q)=\sum_{n\ge1}\mathfrak c(n)q^n$.  On $|q|\le r$ (with $r\le e^{-4\pi}$ small) $h$
is Lipschitz with constant $\Lambda=\sum_n n\,c_J(n)r^{n-1}$, independent of $\mathfrak c$ (\lean{norm_qTailSeries_sub_le}).  Let
$q_\sigma\ne q_\tau$ be the $q$-parameters of $\sigma,\tau$ and suppose $f(\sigma)=f(\tau)$.  Then
$|q_\sigma^{-1}-q_\tau^{-1}|=|q_\tau-q_\sigma|/(|q_\sigma||q_\tau|)$ equals $|h(q_\tau)-h(q_\sigma)|\le\Lambda|q_\tau-q_\sigma|$, so
$|q_\sigma||q_\tau|\ge1/\Lambda$, impossible if both are small; hence $q_\sigma=q_\tau$, i.e.\ $\sigma-\tau\in\Z$. \qed
\begin{physics}High up, the twisted partition function is essentially $q^{-1}$, which is a coordinate on the cylinder
$\chnineteenH/\Z$: distinct low-temperature tori are distinguished by their partition function.\end{physics}
\begin{status}No hypotheses; standard axioms (\leanpath{MoonshineCNGenusZeroAudit.lean}).\end{status}

\begin{theorem}[Large values come from the top]\label{ch19:thm:largevalues}
For $J$-dominated $\mathfrak c$ and every height $Y_0$ there is $R$ such that every $w$ with $|w|\ge R$ equals $f(\sigma)$ for some $\sigma$ with
$\operatorname{Im}\sigma\ge Y_0$.  In particular every member of a replicable family takes its large values from the top.
\end{theorem}
\begin{leancode}
theorem exists_large_values_from_height {c : ℕ → ℂ} (hc : IsJDominated c) (Y₀ : ℝ) :
    ∃ R : ℝ, ∀ w : ℂ, R ≤ ‖w‖ → ∃ σ : ℍ, Y₀ ≤ σ.im ∧ qSeriesFun c σ = w
theorem ReplicableFamily.largeValuesFromTop {N : ℕ} (F : ReplicableFamily N) (k : ℕ) :
    LargeValuesFromTop (qSeriesFun (F.c k))
\end{leancode}
(\leanpath{VOA/Moonshine/MoonshineCNGenusZeroTop.lean}.)
\paragraph{Proof on paper.}  $\phi(q)=q/(1+qh(q))=1/f$ is holomorphic near $q=0$ with $\phi(0)=0$, $\phi'(0)=1$; by the inverse
function theorem it maps a small disc onto a neighbourhood of $0$.  For $|w|$ large, $w^{-1}=\phi(q)$ with small $q\ne0$, i.e.\
$f=w$ at a point whose $q$-parameter is small, i.e.\ high up. \qed
\begin{physics}Large values of the partition function are low-temperature values: only the vacuum pole produces them.\end{physics}
\begin{status}Hypothesis \lean{IsJDominated c}; standard axioms (\leanpath{MoonshineCNGenusZeroAudit.lean}).  Supporting:
\lean{qTailSeries_differentiableOn}, \lean{largeValuesFromTop_qSeriesFun}.\end{status}

\paragraph{Supporting declarations of the analytic basics.}  File \leanpath{MoonshineCNGenusZeroAnalytic.lean}:
\lean{jCoeffMajorant_nonneg}, \lean{summable_norm_mul_pow_of_summable}, \lean{norm_qPar_eq} ($|q|=e^{-2\pi\operatorname{Im}\tau}$),
\lean{real_exp_mul_succ}, \lean{summable_jCoeffMajorant_mul_exp}, \lean{jTail_nonneg}, \lean{exp_neg_mul_le_of_le},
\lean{jTail_le_of_le}, \lean{qTailSeries}, \lean{qSeriesFun_eq_inv_add_qTailSeries}, \lean{norm_qTail_term_le},
\lean{summable_qTail_term}, \lean{norm_qTailSeries_le}, \lean{norm_pow_sub_pow_le_of_norm_le}, \lean{norm_qTailSeries_sub_le},
\lean{summable_lipschitz_majorant}, and for families \lean{mdifferentiable}, \lean{tendsto_sub_inv_qPar}, \lean{c_add_mul},
\lean{c_mul_eq_of_modEq_one}.  All in \leanpath{MoonshineCNGenusZeroAudit.lean}.

%======================================================================
\section{Genus zero II: the prime Hecke identity, nontriviality and property P}\label{ch19:sec:lrinputs}
%======================================================================

These are the inputs proved on Leonardo Rastelli's side (namespace \lean{VOA.Moonshine.ConwayNorton.M8}).

\subsection{The prime Hecke identity}

\begin{theorem}[Prime Hecke identity]\label{ch19:thm:primehecke}
For every replicable family $F$ of period $N$ and every member $k$, the prime Hecke identity holds at every prime
$p\equiv1\pmod N$.
\end{theorem}
\begin{leancode}
noncomputable def heckeBlock {N : ℕ} (F : ReplicableFamily N) (k a d : ℕ) (τ : ℍ) : ℂ :=
  if h : 0 < a ∧ 0 < d then
    ∑ b ∈ Finset.range d, qSeriesFun (F.c (k * a)) (heckePoint a d b h.1 h.2 τ)
  else 0
theorem heckeBlock_at_mul {N : ℕ} (F : ReplicableFamily N) (k a d p : ℕ) (ha : 0 < a)
    (hd : 0 < d) (hp : 0 < p) (hc : F.c (k * (a * p)) = F.c (k * a)) (τ : ℍ) :
    heckeBlock F k a d (heckePoint p 1 0 hp one_pos τ) = heckeBlock F k (a * p) d τ
theorem sum_heckeBlock_at_div_of_not_dvd {N : ℕ} (F : ReplicableFamily N) (k a d p : ℕ)
    (ha : 0 < a) (hd : 0 < d) (hp : p.Prime) (hpa : ¬ p ∣ a) (τ : ℍ) :
    ∑ j ∈ Finset.range p, heckeBlock F k a d (heckePoint 1 p j one_pos hp.pos τ) =
      heckeBlock F k a (d * p) τ
theorem sum_heckeBlock_at_div_of_eq_mul {N : ℕ} (F : ReplicableFamily N) (k a₀ d p : ℕ)
    (ha₀ : 0 < a₀) (hd : 0 < d) (hp : 0 < p) (hc : F.c (k * (p * a₀)) = F.c (k * a₀)) (τ : ℍ) :
    ∑ j ∈ Finset.range p, heckeBlock F k (p * a₀) d (heckePoint 1 p j one_pos hp τ) =
      (p : ℂ) * heckeBlock F k a₀ d τ
theorem ReplicableFamily_primeHeckeIdentity {N : ℕ} (F : ReplicableFamily N) (k : ℕ) :
    PrimeHeckeIdentity F k
\end{leancode}
(\leanpath{VOA/Moonshine/MoonshineLRHecke.lean}.)
\begin{dictionary}
\lean{heckeBlock F k a d} is $B(a,d)(\tau)=\sum_{b<d}f_{ka}((a\tau+b)/d)$, so $H_m=\sum_{a\mid m}B(a,m/a)$
(\lean{heckeOrbitSum_eq_sum_heckeBlock}).  The theorem has \emph{no} hypotheses besides the family: note that $N=0$ is allowed
(then $p\equiv1\pmod0$ means $p=1$, not a prime, and the statement is vacuous).
\end{dictionary}
\paragraph{Proof on paper.}  Let $p\equiv1\pmod N$ be prime.  Three block identities:
\begin{enumerate}[label=(\roman*)]
\item $B(a,d)(p\tau)=B(ap,d)(\tau)$, because $(a\,p\tau+b)/d$ is the point of the block $(ap,d)$ and $f_{kap}=f_{ka}$ ($kap\equiv ka$);
\item if $p\nmid a$: $\sum_{j<p}B(a,d)\bigl(\frac{\tau+j}p\bigr)=\sum_{j,b}f_{ka}\bigl(\frac{a\tau+aj+bp}{dp}\bigr)=B(a,dp)(\tau)$, since
  $(j,b)\mapsto aj+bp$ is a bijection $\Z/p\times\Z/d\to\Z/dp$ and $f$ is $1$-periodic (\lean{sum_comp_eq_sum_range_of_periodic});
\item if $a=pa_0$: $\sum_{j<p}B(pa_0,d)\bigl(\frac{\tau+j}p\bigr)=\sum_j\sum_bf_{kpa_0}\bigl(\frac{a_0\tau+a_0j+b}d\bigr)=p\,B(a_0,d)(\tau)$, since for each
  $j$ the inner sum runs over a full residue system and $f_{kpa_0}=f_{ka_0}$.
\end{enumerate}
Now sum over $a\mid m$ with $d=m/a$.  The divisors $a'$ of $mp$ are of two kinds: $a'=ap$ with $a\mid m$ (block
$B(ap,m/a)$, from (i)), and $a'\mid m$ with $p\nmid a'$ (block $B(a',mp/a')$, from (ii)); the divisors $a=pa_0$ of $m$ give, by
(iii), $p\,H_{m/p}$. \qed
\begin{physics}The composition law of covering operations: the $p$-fold covers of the $m$-fold covers are the $mp$-fold
covers, each counted once, except those that factor through the isogeny ``multiply the lattice by $p$'', counted $p$ times.
The congruence $p\equiv1\pmod N$ ensures that the twist is unchanged along the $p$-fold covering.\end{physics}
\begin{status}\textbf{Unconditional} (for every family); standard axioms (\leanpath{VOA/Moonshine/MoonshineLRHeckeAudit.lean}, 14
reports).  It uses only the periodicity of the family, not replication or domination.  Supporting: \lean{qPar_eq_of_eq_add_int},
\lean{qSeriesFun_eq_of_eq_add_int}, \lean{coe_heckePoint_apply}, \lean{ReplicableFamily_c_add_mul},
\lean{ReplicableFamily_c_eq_of_modEq}, \lean{ReplicableFamily_c_mul_mul_of_modEq_one}, \lean{sum_comp_eq_sum_range_of_periodic},
\lean{heckeBlock_of_pos}, \lean{heckeOrbitSum_eq_sum_heckeBlock}.\end{status}

\subsection{The expansion of the replicate Hecke sum}

\begin{theorem}[$q$-expansion of $H_m$]\label{ch19:thm:hecke-exp}
For every family, $k$ and $m\ge1$: $H_m(\tau)-q^{-m}=\sum_{r\ge1}\Bigl(\sum_{a\mid\gcd(m,r)}\frac ma\,c\bigl(ka,\frac{mr}{a^2}\bigr)\Bigr)q^r$.
\end{theorem}
\begin{leancode}
noncomputable def faberRemainderCoeff (c : ℕ → ℕ → ℂ) (k m r : ℕ) : ℂ :=
  ∑ a ∈ (Nat.gcd m r).divisors, ((m / a : ℕ) : ℂ) * c (k * a) (m * r / (a * a))
theorem hasSum_heckeOrbitSum {N : ℕ} (F : ReplicableFamily N) (k : ℕ) {m : ℕ} (hm : 0 < m)
    (τ : ℍ) :
    HasSum (fun r : ℕ => (if r = 0 then 0 else faberRemainderCoeff F.c k m r) * qPar τ ^ r)
      (heckeOrbitSum F k m τ - (qPar τ ^ m)⁻¹)
\end{leancode}
(\leanpath{VOA/Moonshine/MoonshineLRTower.lean}, \leanpath{VOA/Moonshine/MoonshineLRHeckeExpansion.lean}.)
\begin{dictionary}
The coefficient \lean{faberRemainderCoeff F.c k m r} is the replication value of $[q^r](F_{k,m}(f_k)-q^{-m})$; with $k=1$ it is
$\rho_c(m,r)$ of \S\ref{ch19:sec:replication}.  No domination hypothesis appears, because members of a family are
dominated by definition.
\end{dictionary}
\paragraph{Proof on paper.}  The computation of Theorem~\ref{ch19:thm:qheckesum}, block by block, with replicate $ka$ in block
$a$ (\lean{hasSum_heckeBlock}, using the root-of-unity sums \lean{sum_range_pow_rootOfUnity} and \lean{sum_inv_mul_rootOfUnity}). \qed
\begin{physics}The state content of the $m$-fold covers with twists.\end{physics}
\begin{status}No hypotheses; standard axioms (\leanpath{VOA/Moonshine/MoonshineLRNontrivialAudit.lean}).  Supporting:
\lean{heckePoint_eq_vadd}, \lean{qPar_heckePoint}, \lean{qPar_heckePoint_zero_pow}, \lean{sum_inv_mul_rootOfUnity},
\lean{hasSum_heckeBlock}.\end{status}

\subsection{The coefficients of \texorpdfstring{$J$}{J} are unbounded}

\begin{theorem}[$J$ has unbounded coefficients]\label{ch19:thm:Junbounded}
There is no $C$ with $|[q^n]\mathtt{JShift}|\le C$ for all $n$.
\end{theorem}
\begin{leancode}
theorem JShift_coeff_unbounded : ¬ ∃ C : ℝ, ∀ n, ‖PowerSeries.coeff n JShift‖ ≤ C
\end{leancode}
(\leanpath{VOA/Moonshine/MoonshineLRJUnbounded.lean}.)
\begin{dictionary}
A trivial-looking statement ($c(1)=196884$, $c(2)=21493760,\dots$ grow like $e^{4\pi\sqrt n}$), but the library does not know
any growth of the individual coefficients: it is proved from modularity.  It is the only input about $J$ in the
nontriviality argument.
\end{dictionary}
\paragraph{Proof on paper.}  Suppose $|a_n|\le C$ for all coefficients $a_n$ of $F(z)=q(j(z)-744)=\sum a_nq^n$, $a_0=1$.  At $z=iY$,
$s=e^{-2\pi Y}$: $|F(z)-1|\le Cs/(1-s)$, so $|F(z)|\ge1-2Cs$ for large $Y$.  At $w=S z=i/Y$, $t=e^{-2\pi/Y}$: $|F(w)|\le C/(1-t)$.
Since $j(w)=j(z)$, $|F(z)|=(s/t)|F(w)|\le\frac{s}{t}\cdot\frac C{1-t}$, which is about $\frac{C\,Y}{2\pi}e^{-2\pi Y}\to0$.  For large $Y$
this contradicts $|F(z)|\ge1-2Cs\to1$ (the library makes this explicit: $e^{2\pi Y}\le KY$ with $K=2C+2Ce^{2\pi}$, false at
$Y=K+1$). \qed
\begin{physics}Modular invariance forces exponential growth of the density of states (the Cardy formula): a CFT whose
partition function had bounded level degeneracies could not be modular invariant, because $S$ exchanges low and high
temperature.\end{physics}
\begin{status}No hypotheses; standard axioms (\leanpath{VOA/Moonshine/MoonshineLRJUnboundedAudit.lean}).  Supporting:
\lean{norm_qJ_le_of_bounded}, \lean{norm_qJ_sub_one_le_of_bounded}, \lean{one_le_two_mul_pi_sq}.\end{status}

\subsection{A coefficient bound}

\begin{theorem}[Bounded $q$-series have bounded coefficients]\label{ch19:thm:coeffbound}
If $u(\tau)=\sum_{r\ge0}A_rq^r$ for every $\tau\in\chnineteenH$ and $|u|\le M$ on $\chnineteenH$, then $|A_r|\le M$ for every $r$.
\end{theorem}
\begin{leancode}
theorem norm_coeff_le_of_hasSum_qPar (A : ℕ → ℂ) (u : ℍ → ℂ) (M : ℝ)
    (hu : ∀ τ : ℍ, HasSum (fun r => A r * qPar τ ^ r) (u τ)) (hM : ∀ τ : ℍ, ‖u τ‖ ≤ M)
    (r : ℕ) : ‖A r‖ ≤ M
\end{leancode}
(\leanpath{VOA/Moonshine/MoonshineLRCoeffBound.lean}.)
\paragraph{Proof on paper.}  This is Cauchy's estimate on circles $|q|=\rho\to1$, done discretely: for $L>r$ and
$\zeta=e^{2\pi i/L}$, the average $\frac1L\sum_{t<L}u(\tau+t/L)\zeta^{(L-r)t}$ picks out the coefficients $A_{r'}$ with $r'\equiv r\pmod L$
(\lean{sum_range_pow_rootOfUnity}), so $|A_r||q|^r\le M+\sum_{r'\ge L}|A_{r'}||q|^{r'}$ (\lean{norm_coeff_mul_le_add_tail}); let $L\to\infty$,
then $\operatorname{Im}\tau\to0$. \qed
\begin{physics}A partition function bounded at all temperatures has bounded level degeneracies (in absolute value).\end{physics}
\begin{status}No hypotheses beyond those displayed; standard axioms (\leanpath{VOA/Moonshine/MoonshineLRCoeffBoundAudit.lean}).
Supporting: \lean{sum_range_pow_rootOfUnity}, \lean{norm_exp_two_pi_I_div_natCast}, \lean{qPar_vadd_div},
\lean{norm_coeff_mul_le_add_tail}.\end{status}

\subsection{The tower and nontriviality}

\begin{theorem}[The tower down to $J$]\label{ch19:thm:tower}
\begin{enumerate}[label=(\alph*)]
\item For any $c:\N\to\N\to\C$ and $k$: if for every $m\ge1$ the numbers $\rho_{k}(m,r)=\sum_{a\mid(m,r)}\frac ma\,c(ka,\frac{mr}{a^2})$,
  $r\ge1$, are bounded, then for every $m\ge1$ the sequence $\ell\mapsto c(km,\ell)$, $\ell\ge1$, is bounded.
\item For a replicable family of period $N\ge1$ and every $k$, the numbers $\rho_k(m,r)$ are \emph{not} bounded for every $m$.
\end{enumerate}
\end{theorem}
\begin{leancode}
theorem bounded_replicate_of_bounded_faberRemainderCoeff (c : ℕ → ℕ → ℂ) (k : ℕ)
    (h : ∀ m : ℕ, 1 ≤ m → ∃ M : ℝ, ∀ r : ℕ, 1 ≤ r → ‖faberRemainderCoeff c k m r‖ ≤ M) :
    ∀ m : ℕ, 1 ≤ m → ∃ M : ℝ, ∀ ℓ : ℕ, 1 ≤ ℓ → ‖c (k * m) ℓ‖ ≤ M
theorem ReplicableFamily_c_mul_period {N : ℕ} (F : ReplicableFamily N) (k : ℕ) :
    F.c (k * N) = F.c 0
theorem not_bounded_faberRemainderCoeff {N : ℕ} (hN : 0 < N) (F : ReplicableFamily N) (k : ℕ) :
    ¬ ∀ m : ℕ, 1 ≤ m → ∃ M : ℝ, ∀ r : ℕ, 1 ≤ r → ‖faberRemainderCoeff F.c k m r‖ ≤ M
\end{leancode}
(\leanpath{VOA/Moonshine/MoonshineLRTower.lean}.)
\paragraph{Proof on paper.}  (a) Strong induction on $m$.  At $r=m\ell$ (\lean{faberRemainderCoeff_mul_self}),
$\rho_k(m,m\ell)=c(km,\ell)+\sum_{a\mid m,\,a<m}\frac ma\,c\bigl(ka,\frac{m^2\ell}{a^2}\bigr)$; every term of the sum involves a replicate $ka$ with
$a<m$, bounded by the induction hypothesis; so $c(km,\ell)$ is bounded.  (b) $c(kN,\cdot)=c(0,\cdot)$ by periodicity, the
coefficients of $J$; by (a) with $m=N$ they would be bounded, contradicting Theorem~\ref{ch19:thm:Junbounded}. \qed
\begin{physics}Walking down the tower of covers from the twisted sector $g^k$ to $g^{kN}=1$ reaches the untwisted theory,
whose density of states grows; so the twisted sectors cannot all have bounded remainders.\end{physics}
\begin{status}(a) no hypotheses; (b) hypothesis $0<N$; standard axioms (\leanpath{VOA/Moonshine/MoonshineLRTowerAudit.lean}).\end{status}

\begin{theorem}[Nontriviality]\label{ch19:thm:nontrivial}
For a replicable family of period $N\ge1$, no member $f_k$ is injective modulo translations.
\end{theorem}
\begin{leancode}
theorem norm_lt_of_injective_of_largeValues {f : ℍ → ℂ} (hinj : IsInjectiveModTranslations f)
    {R Y : ℝ} (hR : ∀ w : ℂ, R ≤ ‖w‖ → ∃ σ : ℍ, Y ≤ σ.im ∧ f σ = w) {τ : ℍ} (hτ : τ.im < Y) :
    ‖f τ‖ < R
theorem exists_bound_heckeOrbitSum_sub {N : ℕ} (F : ReplicableFamily N) (k : ℕ) {m : ℕ}
    (hm : 0 < m) {R Y : ℝ} (hY : 0 < Y)
    (hlow : ∀ τ : ℍ, τ.im < Y → ‖qSeriesFun (F.c k) τ‖ < R) :
    ∃ M : ℝ, ∀ τ : ℍ, ‖heckeOrbitSum F k m τ - (qPar τ ^ m)⁻¹‖ ≤ M
theorem ReplicableFamily_notInjectiveModTranslations {N : ℕ} (hN : 0 < N)
    (F : ReplicableFamily N) : F.NotInjectiveModTranslations
\end{leancode}
(\leanpath{VOA/Moonshine/MoonshineLRNontrivial.lean}.)
\begin{dictionary}
The conclusion \lean{F.NotInjectiveModTranslations} is $\forall k,\ \neg\,\mathtt{IsInjectiveModTranslations}(f_k)$.  The
only hypothesis is $0<N$.  The statement is stronger than what the blueprint asked for (it does not need the prime Hecke
identity).
\end{dictionary}
\paragraph{Proof on paper.}  Suppose $f=f_k$ is injective modulo $\Z$.  By Theorem~\ref{ch19:thm:largevalues} there are $R,Y>0$ such
that every value of modulus $\ge R$ is taken at height $\ge Y$; if $\operatorname{Im}\tau<Y$ and $|f(\tau)|\ge R$, then $f(\tau)=f(\sigma)$ with
$\operatorname{Im}\sigma\ge Y$, and injectivity modulo $\Z$ forces $\operatorname{Im}\sigma=\operatorname{Im}\tau$, a contradiction; so $|f|<R$ below
height $Y$.  Fix $m\ge1$ and the replication polynomial $P$ with $P(f)=H_m$.  Then $u=H_m-q^{-m}$ is bounded on $\chnineteenH$: below
height $Y$ by $\max_{|w|\le R}|P(w)|+e^{2\pi mY}$, and above by its absolutely convergent expansion (Theorem~\ref{ch19:thm:hecke-exp}).
By Theorem~\ref{ch19:thm:coeffbound} its coefficients $\rho_k(m,r)$ are bounded, for every $m$, contradicting
Theorem~\ref{ch19:thm:tower}(b). \qed
\begin{physics}A twisted partition function that separated all points of the cylinder $\chnineteenH/\Z$ would have no
nontrivial modular symmetry; the covering relations, together with the growth of the untwisted density of states, forbid
this.  This excludes in particular the ``modular fictions'' $q^{-1}+\zeta q$ of Cummins--Gannon, which are injective modulo
$\Z$ and satisfy replication, but are not attached to $J$ (their period-$N$ families violate the clause \lean{zero_eq_J}, $c(0,\cdot)=J$, of \lean{ReplicableFamily}).\end{physics}
\begin{status}Hypothesis $0<N$; \textbf{unconditional} otherwise.  Standard axioms
(\leanpath{VOA/Moonshine/MoonshineLRNontrivialAudit.lean}).\end{status}

\begin{theorem}[Property P, unconditionally]\label{ch19:thm:propPLR}
For a replicable family of period $N\ge1$, every $k$, and $z,w\in\chnineteenH$: if $f_k(z)=f_k(w)$ then $\gamma z=w$ for some $\gamma\in
G(f_k)$.
\end{theorem}
\begin{leancode}
theorem ReplicableFamily_exists_mem_symmetryGroup_smul_eq {N : ℕ} (hN : 0 < N)
    (F : ReplicableFamily N) (k : ℕ) {z w : ℍ}
    (hzw : qSeriesFun (F.c k) z = qSeriesFun (F.c k) w) :
    ∃ γ ∈ symmetryGroup (qSeriesFun (F.c k)), γ • z = w
\end{leancode}
(\leanpath{VOA/Moonshine/MoonshineLRPropertyP.lean}.)
\paragraph{Proof on paper.}  Theorem~\ref{ch19:thm:propP} below, whose only hypothesis, the prime Hecke identity, is
Theorem~\ref{ch19:thm:primehecke}. \qed
\begin{physics}Equal twisted partition functions at two moduli imply the two tori (with twist) are related by a symmetry.\end{physics}
\begin{status}Hypothesis $0<N$; standard axioms (\leanpath{MoonshineLRNontrivialAudit.lean}, which imports
\leanpath{MoonshineLRPropertyP.lean}).\end{status}

%======================================================================
\section{Genus zero III: the Hecke pairing}\label{ch19:sec:pairing}
%======================================================================

From now on $F$ is a replicable family of period $N$, $k$ a member, $f=f_k$, and $p$ a prime with $p\equiv1\pmod N$ (an
\emph{admissible} prime).  $\Delta(p)$ is the set of the $p+1$ maps $\tau\mapsto p\tau$ and $\tau\mapsto(\tau+j)/p$, $0\le j<p$.

\begin{theorem}[Polynomial Hecke sums; membership; pairing]\label{ch19:thm:pairing}
Assume the prime Hecke identity at $p$.
\begin{enumerate}[label=(\alph*)]
\item For every $P\in\C[X]$ there is $R\in\C[X]$ with $\sum_{\beta\in\Delta(p)}P\bigl(f(\beta\tau)\bigr)=R\bigl(f(\tau)\bigr)$ for all $\tau$.
\item If two finite multisets $s,t$ of complex numbers have $\sum_{x\in s}P(x)=\sum_{y\in t}P(y)$ for every polynomial $P$, then
  every member of $s$ is a member of $t$.
\item (Hecke pairing) If $f(z)=f(w)$, then every value $f(\beta z)$, $\beta\in\Delta(p)$, equals some value $f(\beta'w)$, $\beta'\in\Delta(p)$.
\end{enumerate}
\end{theorem}
\begin{leancode}
noncomputable def heckeValueSum (P : ℂ[X]) (τ : ℍ) : ℂ :=
  (∑ j ∈ Finset.range p, P.eval (qSeriesFun (F.c k) (heckePoint 1 p j one_pos hp τ))) +
    P.eval (qSeriesFun (F.c k) (heckePoint p 1 0 hp one_pos τ))
noncomputable def heckeValues (τ : ℍ) : Multiset ℂ :=
  (Finset.range p).val.map (fun j => qSeriesFun (F.c k) (heckePoint 1 p j one_pos hp τ)) +
    {qSeriesFun (F.c k) (heckePoint p 1 0 hp one_pos τ)}
theorem exists_poly_heckeValueSum (hB1 : PrimeHeckeIdentityAt F k p hp) (P : ℂ[X]) :
    ∃ R : ℂ[X], ∀ τ : ℍ, heckeValueSum F k p hp P τ = R.eval (qSeriesFun (F.c k) τ)
theorem mem_of_forall_sum_eval_eq {s t : Multiset ℂ}
    (h : ∀ P : ℂ[X], (s.map P.eval).sum = (t.map P.eval).sum) {x : ℂ} (hx : x ∈ s) :
    x ∈ t
theorem heckeValues_pairing (hB1 : PrimeHeckeIdentityAt F k p hp) {z w : ℍ}
    (h : qSeriesFun (F.c k) z = qSeriesFun (F.c k) w) {x : ℂ}
    (hx : x ∈ heckeValues F k p hp z) : x ∈ heckeValues F k p hp w
\end{leancode}
(\leanpath{VOA/Moonshine/MoonshineCNGenusZeroPairing.lean}; section variables \lean{{N : ℕ} (F : ReplicableFamily N) (k p : ℕ)
(hp : 0 < p)}.)
\begin{dictionary}
\lean{heckeValueSum F k p hp P τ} is $\sum_{\beta\in\Delta(p)}P(f(\beta\tau))$, \lean{heckeValues} the multiset of the $p+1$ values.
\lean{Multiset ℂ} is a finite multiset; \lean{s.map P.eval} applies $P$ to its members.  Note that (a)--(c) are stated for a
fixed $p$ with the prime Hecke identity at $p$ as hypothesis (\lean{hB1 : PrimeHeckeIdentityAt F k p hp}); primality and
admissibility are not needed here, only the identity.
\end{dictionary}
\paragraph{Proof on paper.}  (a) For $P=P_m$, the monic replication polynomial with $P_m(f)=H_m$, the prime Hecke identity
gives $\sum_\beta P_m(f(\beta\tau))=\sum_\beta H_m(\beta\tau)=H_{mp}(\tau)+[p\mid m]pH_{m/p}(\tau)=P_{mp}(f(\tau))+[p\mid m]pP_{m/p}(f(\tau))$.  The
set of $P$ for which (a) holds is closed under sums and scalar multiples and contains $1$ and a monic polynomial of each
positive degree, so it is all of $\C[X]$ (\lean{forall_poly_of_monic_basis}).  (b) Given $x\in s$, test with
$P=\prod_{y\in t}(X-y)\cdot\prod_{x'\in s,\ x'\ne x}(X-x')$ (distinct $x'$): the $t$-sum is $0$, the $s$-sum is
$(\text{multiplicity of }x)\cdot P(x)$, nonzero unless $x\in t$.  No Newton identities are needed.  (c) By (a) the power sums
$\sum_\beta f(\beta z)^r$ and $\sum_\beta f(\beta w)^r$ agree for every $r$ (both are $R(f(z))=R(f(w))$); apply (b). \qed
\begin{physics}The $p$-fold covering correspondence is a ``modular equation'': the multiset of twisted partition functions of
the $p+1$ covers of a torus is determined by the partition function of the torus itself.  So two tori with equal partition
functions have covers with equal partition functions.\end{physics}
\begin{status}Hypothesis \lean{PrimeHeckeIdentityAt F k p hp}; standard axioms (\leanpath{MoonshineCNGenusZeroAudit.lean}).
Supporting: \lean{forall_poly_of_monic_basis}, \lean{eval_prod_X_sub_C}, \lean{sum_map_eval_heckeValues},
\lean{mem_heckeValues_iff}, \lean{hecke_pairing_shift}, \lean{hecke_pairing_scale} (the two explicit forms of (c)).\end{status}

\begin{theorem}[Pairing on connected domains; partner uniqueness]\label{ch19:thm:partner}
Let $f$ also denote its extension $z\mapsto f(z)$ to $\{z\in\C:\operatorname{Im}z>0\}$ (\lean{qSeriesExt}).
\begin{enumerate}[label=(\alph*)]
\item \emph{Hypothesis: the prime Hecke identity at $p$.}  Let $\Omega\subseteq\C$ be open, connected and nonempty, and
  $u,v:\Omega\to\chnineteenH$ holomorphic with $f\circ u=f\circ v$ on $\Omega$.  For every $\beta_0\in\Delta(p)$ there is one $\beta\in\Delta(p)$
  with $f(\beta_0(u(z)))=f(\beta(v(z)))$ for all $z\in\Omega$.
\item \emph{Same hypothesis.}  Iterating: for every $n$ there are $e\in\Z$, $|e|\le n$, and $b\in\chnineteenR$ with
  $f(p^nu(z))=f(p^ev(z)+b)$ on $\Omega$.
\item (No hypothesis.)  If $p>1$ and $f\circ\beta=f\circ\beta'$ at points accumulating at a point of $\chnineteenH$, then $\beta=\beta'$.
\end{enumerate}
\end{theorem}
\begin{leancode}
noncomputable def affC (a b : ℝ) (z : ℂ) : ℂ := a * z + b
noncomputable def heckeAffine (p : ℕ) (β : Option (Fin p)) : ℂ → ℂ :=
  affC (heckeDil p β) (heckeTr p β)
noncomputable def qSeriesExt (c : ℕ → ℂ) : ℂ → ℂ :=
  qSeriesFun c ∘ ofComplex
theorem exists_heckeAffine_partner (hB1 : PrimeHeckeIdentityAt F k p hp) {D : Set ℂ}
    (hDo : IsOpen D) (hDc : IsPreconnected D) (hDn : D.Nonempty) {u v : ℂ → ℂ}
    (hu : DifferentiableOn ℂ u D) (hv : DifferentiableOn ℂ v D)
    (huU : ∀ z ∈ D, 0 < (u z).im) (hvU : ∀ z ∈ D, 0 < (v z).im)
    (huv : ∀ z ∈ D, qSeriesExt (F.c k) (u z) = qSeriesExt (F.c k) (v z)) (M : Option (Fin p)) :
    ∃ β : Option (Fin p), ∀ z ∈ D,
      qSeriesExt (F.c k) (heckeAffine p M (u z)) = qSeriesExt (F.c k) (heckeAffine p β (v z))
theorem exists_iterate_heckeAffine_partner (hB1 : PrimeHeckeIdentityAt F k p hp) {D : Set ℂ}
    (hDo : IsOpen D) (hDc : IsPreconnected D) (hDn : D.Nonempty) {u v : ℂ → ℂ}
    (hu : DifferentiableOn ℂ u D) (hv : DifferentiableOn ℂ v D)
    (huU : ∀ z ∈ D, 0 < (u z).im) (hvU : ∀ z ∈ D, 0 < (v z).im)
    (huv : ∀ z ∈ D, qSeriesExt (F.c k) (u z) = qSeriesExt (F.c k) (v z)) (n : ℕ) :
    ∃ (e : ℤ) (b : ℝ), |e| ≤ n ∧ ∀ z ∈ D,
      qSeriesExt (F.c k) (affC ((p : ℝ) ^ n) 0 (u z)) =
        qSeriesExt (F.c k) (affC ((p : ℝ) ^ e) b (v z))
theorem heckeAffine_eq_of_frequently_eq {c : ℕ → ℂ} (hc : IsJDominated c) {p : ℕ} (hp1 : 1 < p)
    {β β' : Option (Fin p)} {z₀ : ℂ} (hz₀ : 0 < z₀.im)
    (h : ∃ᶠ z in 𝓝[≠] z₀, qSeriesExt c (heckeAffine p β z) = qSeriesExt c (heckeAffine p β' z)) :
    β = β'
\end{leancode}
(\leanpath{VOA/Moonshine/MoonshineCNGenusZeroPartner.lean}.)
\begin{dictionary}
\lean{affC a b z} $=az+b$ with $a,b$ real.  $\Delta(p)$ is encoded by \lean{Option (Fin p)}: \lean{none} is $z\mapsto pz$ and
\lean{some j} is $z\mapsto(z+j)/p$ (\lean{heckeDil}, \lean{heckeTr} are the dilation and translation parts).
\lean{qSeriesExt c} $=\mathtt{qSeriesFun}\ c\circ\mathtt{ofComplex}$, where Mathlib's \lean{UpperHalfPlane.ofComplex} is a
partial inverse of the inclusion $\chnineteenH\to\C$; only values with $\operatorname{Im}>0$ are ever used, and there
\lean{qSeriesExt c z} $=f(z)$ (\lean{qSeriesExt_of_im_pos}).  \lean{DifferentiableOn ℂ u D} is holomorphy on $D$;
\lean{IsPreconnected} is connectedness (allowing the empty set, excluded by \lean{hDn}).  In (c),
\lean{∃ᶠ z in 𝓝[≠] z₀, ...} says the equation holds at points arbitrarily close to $z_0$, $z\ne z_0$.
\end{dictionary}
\paragraph{Proof on paper.}  (a) At each $z\in\Omega$ the pointwise pairing (Theorem~\ref{ch19:thm:pairing}(c) at $u(z),v(z)$) gives a
partner $\beta(z)$.  Take a small closed disc $K\subset\Omega$; it is uncountable and $\Delta(p)$ is finite, so one $\beta$ serves on an
infinite subset of $K$, which accumulates at a point of $K\subset\Omega$ (compactness).  The identity theorem for the holomorphic
function $f\circ\beta_0\circ u-f\circ\beta\circ v$ on the connected $\Omega$ extends the equation to all of $\Omega$.  (b) Iterate (a) $n$ times
with $\beta_0=(z\mapsto pz)$; each step multiplies the dilation on the $v$-side by $p$ or $p^{-1}$ (the dilations of $\Delta(p)$) and adds a
real translation, so after $n$ steps the dilation is $p^e$ with $|e|\le n$.  (c) The identity theorem extends $f\circ\beta=f\circ\beta'$
to $\chnineteenH$; at $z=iy$ with $y$ large both $\beta z$ and $\beta'z$ are in the top, so by top injectivity they differ by an
integer for all large $y$; comparing imaginary parts gives equal dilations and then equal translations. \qed
\begin{physics}The covering correspondence can be followed continuously in the moduli: along a connected family of pairs
of tori with equal partition functions, the matching of their covers does not jump.\end{physics}
\begin{status}(a), (b): hypothesis \lean{PrimeHeckeIdentityAt F k p hp}; (c): \lean{IsJDominated c}, $p>1$.  Standard axioms
(\leanpath{MoonshineCNGenusZeroAudit2.lean}).  Supporting (file \leanpath{MoonshineCNGenusZeroPartner.lean}): the
\lean{affC} calculus (\lean{affC_apply}, \lean{affC_re}, \lean{affC_im}, \lean{affC_affC}, \lean{affC_one_zero},
\lean{affC_inv_affC}, \lean{affC_affC_inv}, \lean{differentiable_affC}, \lean{continuous_affC}, \lean{affC_im_pos},
\lean{norm_affC_sub_le}), \lean{heckeDil}, \lean{heckeTr}, \lean{heckeDil_pos}, \lean{heckeDil_eq_zpow},
\lean{heckeAffine_im_pos}, \lean{heckeAffine_some_coe}, \lean{heckeAffine_none_coe}, \lean{qSeriesExt_coe},
\lean{qSeriesExt_of_im_pos}, \lean{differentiableOn_qSeriesExt}, \lean{qSeriesExt_add_intCast} ($1$-periodicity),
\lean{topHeight}, \lean{topHeight_pos}, \lean{exists_int_of_qSeriesExt_eq} (top injectivity on $\C$),
\lean{heckeAffine_pairing_coe}, \lean{heckeAffine_pairing}.\end{status}

%======================================================================
\section{Genus zero IV: local symmetries are real M\"obius maps; property P}\label{ch19:sec:local}
%======================================================================

A \emph{local symmetry} of $f$ is a holomorphic map $\alpha$ on a disc in $\chnineteenH$, with values in $\chnineteenH$, such that
$f\circ\alpha=f$.  The goal of this section is: \emph{every local symmetry with $\alpha'\ne0$ is the restriction of an element of
$G(f)$}.  Classically (Cummins--Gannon, \S\S3--4) one continues $\alpha$ analytically along paths to a maximal domain, shows the
maximal domain is invariant under a dense group of affine maps, hence is $\chnineteenH$ or a sector, excludes sectors, and
concludes with $\Aut(\chnineteenH)=\mathrm{PSL}_2(\chnineteenR)$.  None of path continuation, maximal domains, the density lemma or
$\Aut(\chnineteenH)=\mathrm{PSL}_2(\chnineteenR)$ is in Mathlib v4.30.  The library replaces this by a different argument
: \emph{one} intertwining sequence on a fixed disc yields a real linear ODE, whose solutions are
explicit.

\subsection{Local symmetries at non-critical values; relations in the top}

\begin{theorem}[Local symmetries and the top]\label{ch19:thm:localtop}
\begin{enumerate}[label=(\alph*)]
\item If $g$ is holomorphic on $\{\operatorname{Im}>0\}$, $g(z_0)=g(w_0)$ and $g'(w_0)\ne0$, there is a holomorphic $\alpha$ on a disc
  around $z_0$ with $\alpha(z_0)=w_0$, values in $\{\operatorname{Im}>0\}$, and $g\circ\alpha=g$.
\item If $f\circ u=f\circ v$ on a connected set $S$, with $u,v$ continuous and taking values in the top, then $u=v+n$ on $S$ for one
  integer $n$.
\item If $f(\delta z+\tau)=f(\mu\alpha(z)+\nu)$ on a connected open $D$ ($\delta,\mu>0$ real, $\alpha$ holomorphic), the right side lies in the
  top everywhere and the left side strictly above the top height at one point, then $\alpha(z)=az+b$ on $D$ with $a>0$, $a,b$ real.
\end{enumerate}
\end{theorem}
\begin{leancode}
theorem exists_localSymmetry_of_deriv_ne_zero {g : ℂ → ℂ}
    (hg : DifferentiableOn ℂ g upperHalfPlaneSet) {z₀ w₀ : ℂ} (hz₀ : 0 < z₀.im)
    (hw₀ : 0 < w₀.im) (heq : g z₀ = g w₀) (hd : _root_.deriv g w₀ ≠ 0) :
    ∃ η > 0, ∃ α : ℂ → ℂ, α z₀ = w₀ ∧ DifferentiableOn ℂ α (Metric.ball z₀ η) ∧
      ∀ z ∈ Metric.ball z₀ η, 0 < (α z).im ∧ g (α z) = g z
theorem exists_int_eq_add_of_top {c : ℕ → ℂ} (hc : IsJDominated c) {S : Set ℂ}
    (hS : IsPreconnected S) {u v : ℂ → ℂ} (hu : ContinuousOn u S) (hv : ContinuousOn v S)
    (htu : ∀ z ∈ S, topHeight ≤ (u z).im) (htv : ∀ z ∈ S, topHeight ≤ (v z).im)
    (huv : ∀ z ∈ S, qSeriesExt c (u z) = qSeriesExt c (v z)) :
    ∃ n : ℤ, ∀ z ∈ S, u z = v z + n
theorem exists_affC_eqOn_of_top {c : ℕ → ℂ} (hc : IsJDominated c) {D : Set ℂ} (hDo : IsOpen D)
    (hDc : IsPreconnected D) {α : ℂ → ℂ} (hα : DifferentiableOn ℂ α D) {δ τ μ ν : ℝ}
    (hδ : 0 < δ) (hμ : 0 < μ)
    (hrel : ∀ z ∈ D, qSeriesExt c (affC δ τ z) = qSeriesExt c (affC μ ν (α z)))
    (htop : ∀ z ∈ D, topHeight ≤ (affC μ ν (α z)).im) {z₁ : ℂ} (hz₁ : z₁ ∈ D)
    (hz₁top : topHeight < (affC δ τ z₁).im) :
    ∃ a b : ℝ, 0 < a ∧ ∀ z ∈ D, α z = affC a b z
\end{leancode}
(\leanpath{VOA/Moonshine/MoonshineCNGenusZeroLocal.lean}.)
\begin{dictionary}
\lean{upperHalfPlaneSet} is $\{z\in\C:\operatorname{Im}z>0\}$; \lean{_root_.deriv g w₀} the complex derivative;
\lean{Metric.ball z₀ η} the open disc.  In (c), \lean{affC δ τ z} $=\delta z+\tau$.
\end{dictionary}
\paragraph{Proof on paper.}  (a) $\alpha=g_{w_0}^{-1}\circ g$ with $g_{w_0}^{-1}$ the holomorphic local inverse of $g$ near $w_0$
(Mathlib's analytic inverse function theorem).  (b) By top injectivity, $u(z)-v(z)\in\Z$ for every $z$; a continuous integer-valued
function on a connected set is constant (\lean{eq_of_continuousOn_intValued}).  (c) The set where the left side is in the top is
open and nonempty; on a small disc inside it (a connected set) both sides are in the top, so by (b) $\delta z+\tau=\mu\alpha(z)+\nu+n$, i.e.\ $\alpha$ is real affine on an open set,
hence on $D$ by the identity theorem. \qed
\begin{physics}Near the low-temperature limit, a symmetry of the partition function can only be a shift of the modulus by an
integer.\end{physics}
\begin{status}(a) no hypotheses beyond those displayed; (b), (c) \lean{IsJDominated c}.  Standard axioms
(\leanpath{MoonshineCNGenusZeroAudit2.lean}).  Supporting: \lean{not_lt_of_intValued_aux}, \lean{eq_of_continuousOn_intValued}.\end{status}

\subsection{The overlap construction}

\begin{theorem}[Intertwining sequences]\label{ch19:thm:overlap}
Let $\alpha$ be a local symmetry of $f$ on the open disc $B'=D(z_0,\eta)$ with $\eta<\operatorname{Im}z_0$ (so that even the closed disc lies in $\chnineteenH$), with $\operatorname{Im}\alpha\ge m>0$, and assume the
prime Hecke identity for $k$.  If $\alpha$ is not a real affine map on $B'$, there are real affine maps $C_n(z)=a_nz+b_n$ with $a_n>0$,
$a_n\ne1$, $(a_n,b_n)\to(1,0)$, and real affine $C'_n$ of positive dilation, such that $C_n$ maps $D(z_0,\eta/2)$ into $B'$ and
$\alpha\circ C_n=C'_n\circ\alpha$ on $D(z_0,\eta/2)$.
\end{theorem}
\begin{leancode}
theorem exists_intertwining_of_not_realAffine {N : ℕ} (hN : 0 < N) (F : ReplicableFamily N)
    (k : ℕ) (hB1 : PrimeHeckeIdentity F k) {z₀ : ℂ} {η : ℝ} (hη : 0 < η) (hηz : η < z₀.im)
    {α : ℂ → ℂ} (hα : DifferentiableOn ℂ α (Metric.ball z₀ η)) {m : ℝ} (hm : 0 < m)
    (hαm : ∀ z ∈ Metric.ball z₀ η, m ≤ (α z).im)
    (hαf : ∀ z ∈ Metric.ball z₀ η, qSeriesExt (F.c k) (α z) = qSeriesExt (F.c k) z)
    (hna : ¬ ∃ a b : ℝ, 0 < a ∧ ∀ z ∈ Metric.ball z₀ η, α z = affC a b z) :
    ∃ a b a' b' : ℕ → ℝ, (∀ n, 0 < a n) ∧ (∀ n, a n ≠ 1) ∧ (∀ n, 0 < a' n) ∧
      Tendsto a atTop (𝓝 1) ∧ Tendsto b atTop (𝓝 0) ∧
      ∀ n, ∀ z ∈ Metric.ball z₀ (η / 2), affC (a n) (b n) z ∈ Metric.ball z₀ η ∧
        α (affC (a n) (b n) z) = affC (a' n) (b' n) (α z)
\end{leancode}
(\leanpath{VOA/Moonshine/MoonshineCNGenusZeroOverlap.lean}.)
\begin{dictionary}
\lean{hna} is ``$\alpha$ is not of the form $z\mapsto az+b$ with $a>0$ real on $B'$''.  The conclusion is the intertwining sequence:
\lean{affC (a n) (b n)} is $C_n$, \lean{affC (a' n) (b' n)} is $C'_n$, \lean{Tendsto a atTop (𝓝 1)} is $a_n\to1$.  The abstract core,
with the Hecke input replaced by a sequence of relations, is \lean{exists_intertwining_of_relations} (same file).
\end{dictionary}
\paragraph{Proof on paper.}  Write $L_{\delta,\tau}(z)=\delta z+\tau$ for real $\delta>0$, $\tau$.  Choose two distinct admissible primes
$p,q$ (there are infinitely many primes $\equiv1\pmod N$, \lean{infinite_admissible_primes_gt}), and $w$ so large that $p^w\,m$
exceeds the top height.

\emph{Step 1: a family of relations.}  For $i\ge1$ let $\ell=\ell(i)$ be the integer with $p^\ell\le q^i<p^{\ell+1}$.  Apply
Theorem~\ref{ch19:thm:partner}(b) at the prime $p$, on the connected disc $B'$, with $u=\alpha$, $v=\mathrm{id}$ (allowed since
$f\circ\alpha=f$ on $B'$) and $n=w+\ell+1$:
\begin{equation*}
  f\bigl(p^{w+\ell+1}\alpha(z)\bigr)=f\bigl(p^{e}z+b\bigr)\qquad(z\in B'),\qquad|e|\le w+\ell+1 .
\end{equation*}
This is again an equation $f\circ u'=f\circ v'$ on $B'$, now with $u'(z)=p^ez+b$ and $v'(z)=p^{w+\ell+1}\alpha(z)$ (note the roles: the
side carrying $\alpha$ has become $v'$).  Apply Theorem~\ref{ch19:thm:partner}(b) once more, at the prime $q$ with $n=i$:
\begin{equation}\label{ch19:eq:overlaprel}
  f\bigl(\delta_iz+\tau_i\bigr)=f\bigl(\mu_i\alpha(z)+\nu_i\bigr)\quad(z\in B'),\qquad \delta_i=q^ip^{e},\ \tau_i=q^ib,\
  \mu_i=q^{e'_i}p^{w+\ell+1},
\end{equation}
where $|e'_i|\le i$ comes from (b) (each of the $i$ steps multiplies the dilation by $q^{\pm1}$).

\emph{Step 2: the bounds.}  The right side of \eqref{ch19:eq:overlaprel} lies in the top: $\mu_i\ge q^{-i}p^{\ell+1}p^w>p^w$ because
$p^{\ell+1}>q^i$, so $\operatorname{Im}(\mu_i\alpha(z))\ge p^wm$.  (This is why $\ell$ was chosen as it is.)  The dilations $\delta_i$ are
pairwise distinct, since they contain distinct powers $q^i$ of the prime $q\ne p$ (\lean{zpow_ne_pow_of_prime_ne}).  They are bounded below,
$\delta_i\ge q^ip^{-(w+\ell+1)}\ge p^{-(w+1)}$ since $p^\ell\le q^i$.  They are bounded above: if $\delta_i$ were very large, the left side
$\delta_iz+\tau_i$ would lie above the top height at some $z\in B'$, and Theorem~\ref{ch19:thm:localtop}(c) would make $\alpha$ real affine on
$B'$, which is excluded.  Using $1$-periodicity of $f$ we may replace $\tau_i$ by its fractional part.  By compactness of
$[\mathrm{lo},\mathrm{hi}]\times[0,1]$ the pairs $(\delta_i,\tau_i)$ have a convergent subsequence.

\emph{Step 3: unwinding.}  Take $i<j$ in the subsequence, so close that $C=L_{\delta_j,\tau_j}^{-1}\circ L_{\delta_i,\tau_i}$ maps
$D(z_0,\eta/2)$ into $B'$.  For $z\in D(z_0,\eta/2)$ put $z'=Cz\in B'$, so that $\delta_iz+\tau_i=\delta_jz'+\tau_j$.  By
\eqref{ch19:eq:overlaprel} for $i$ at $z$ and for $j$ at $z'$,
\begin{equation*}
  f\bigl(\mu_i\alpha(z)+\nu_i\bigr)=f(\delta_iz+\tau_i)=f(\delta_jz'+\tau_j)=f\bigl(\mu_j\alpha(Cz)+\nu_j\bigr).
\end{equation*}
Both arguments of $f$ at the ends lie in the top, and $z$ ranges over the connected disc $D(z_0,\eta/2)$, so by
Theorem~\ref{ch19:thm:localtop}(b) there is one integer $n$ with $\mu_i\alpha(z)+\nu_i=\mu_j\alpha(Cz)+\nu_j+n$ for all such $z$.  Solving,
$\alpha(Cz)=C'(\alpha(z))$ with $C'(w)=(\mu_iw+\nu_i-\nu_j-n)/\mu_j$, real affine with positive dilation.  $C$ is real affine with dilation
$\delta_i/\delta_j\ne1$, so $C\ne\mathrm{id}$, and $C\to\mathrm{id}$ as $i,j\to\infty$ along the subsequence.  Choosing such pairs $(i_n,j_n)$
gives the sequences $C_n$, $C'_n$. \qed
\begin{physics}Composing the covering correspondences at two different primes produces exact symmetry relations between
slightly different rescalings of the moduli; comparing two of them gives transformations arbitrarily close to the identity, whose
limit is an infinitesimal symmetry (the next theorem).\end{physics}
\begin{status}Hypotheses: $0<N$, \lean{PrimeHeckeIdentity F k}, and those displayed.  Standard axioms
(\leanpath{MoonshineCNGenusZeroAudit2.lean}).  Supporting: \lean{im_pos_of_mem_ball}, \lean{zpow_ne_pow_of_prime_ne},
\lean{exists_intertwining_of_relations}.\end{status}

\subsection{The real linear ODE and its explicit solutions}

\begin{theorem}[The ODE]\label{ch19:thm:ode-alpha}
Let $\alpha$ be holomorphic on an open set $S$, non-constant on $S$, and $\alpha\circ C_n=C'_n\circ\alpha$ on $S$ for real affine $C_n\to\mathrm{id}$,
$C_n\ne\mathrm{id}$, and real affine $C'_n$.  Then there are real $\lambda,\mu,\lambda',\mu'$ with $(\lambda,\mu)\ne(0,0)$ and
$(\lambda z+\mu)\,\alpha'(z)=\lambda'\alpha(z)+\mu'$ on $S$.
\end{theorem}
\begin{leancode}
theorem exists_real_linear_ode {α : ℂ → ℂ} {S : Set ℂ} (hSo : IsOpen S)
    (hα : DifferentiableOn ℂ α S) {a b a' b' : ℕ → ℝ} (ha : Tendsto a atTop (𝓝 1))
    (hb : Tendsto b atTop (𝓝 0)) (hne : ∀ n, a n ≠ 1 ∨ b n ≠ 0)
    (hrel : ∀ n, ∀ z ∈ S, α (affC (a n) (b n) z) = affC (a' n) (b' n) (α z))
    {z₁ z₂ : ℂ} (hz₁ : z₁ ∈ S) (hz₂ : z₂ ∈ S) (h12 : α z₁ ≠ α z₂) :
    ∃ l m l' m' : ℝ, (l ≠ 0 ∨ m ≠ 0) ∧
      ∀ z ∈ S, ((l : ℂ) * z + m) * _root_.deriv α z = l' * α z + m'
\end{leancode}
(\leanpath{VOA/Moonshine/MoonshineCNGenusZeroODE.lean}.)
\paragraph{Proof on paper.}  Put $\varepsilon_n=|a_n-1|+|b_n|>0$; along a subsequence $((a_n-1)/\varepsilon_n,b_n/\varepsilon_n)\to(\lambda,\mu)$ with
$|\lambda|+|\mu|=1$.  Then $(\alpha(C_nz)-\alpha(z))/\varepsilon_n\to(\lambda z+\mu)\alpha'(z)$.  The same quantity equals
$\frac{a'_n-1}{\varepsilon_n}\alpha(z)+\frac{b'_n}{\varepsilon_n}$; evaluating at two points $z_1,z_2$ with $\alpha(z_1)\ne\alpha(z_2)$ solves for the two
coefficients, which therefore converge; their limits $\lambda',\mu'$ are real because the $a'_n,b'_n$ are. \qed
\begin{physics}An infinitesimal symmetry: the limit of exact symmetries accumulating at the identity is a real vector field
$(\lambda z+\mu)\partial_z$ on the moduli, which the local symmetry must intertwine with another real vector field.\end{physics}
\begin{status}No hypotheses beyond those displayed; standard axioms (\leanpath{MoonshineCNGenusZeroAudit2.lean}).
Supporting: \lean{exists_real_of_tendsto_ofReal}.\end{status}

\begin{theorem}[Explicit solutions]\label{ch19:thm:families}
On a connected open $S\subseteq\{\operatorname{Im}>0\}$, a holomorphic solution of $(lz+m)\alpha'=l'\alpha+m'$ ($l,m,l',m'$ real,
$(l,m)\ne0$) is one of: $\alpha=K(z-c)^A+M_0$ ($A\ne0$; power family), $\alpha=B\log(z-c)+K$ (logarithmic), $\alpha=Ke^{Az}+M_0$ ($A\ne0$;
exponential), $\alpha=Bz+K$ (affine), with $A,B,c,M_0$ real and $K$ complex.
\end{theorem}
\begin{leancode}
theorem eqOn_explicit_of_real_linear_ode {α : ℂ → ℂ} {S : Set ℂ} (hSo : IsOpen S)
    (hSc : IsPreconnected S) (hSU : ∀ z ∈ S, 0 < z.im) (hα : DifferentiableOn ℂ α S)
    {l m l' m' : ℝ} (hlm : l ≠ 0 ∨ m ≠ 0)
    (hode : ∀ z ∈ S, ((l : ℂ) * z + m) * _root_.deriv α z = l' * α z + m') :
    (∃ (c A M₀ : ℝ) (K : ℂ), A ≠ 0 ∧ ∀ z ∈ S, α z = K * (z - c) ^ (A : ℂ) + M₀) ∨
    (∃ (c B : ℝ) (K : ℂ), ∀ z ∈ S, α z = B * Complex.log (z - c) + K) ∨
    (∃ (A M₀ : ℝ) (K : ℂ), A ≠ 0 ∧ ∀ z ∈ S, α z = K * Complex.exp (A * z) + M₀) ∨
    (∃ (B : ℝ) (K : ℂ), ∀ z ∈ S, α z = B * z + K)
\end{leancode}
(\leanpath{VOA/Moonshine/MoonshineCNGenusZeroFamilies.lean}.)
\begin{dictionary}
\lean{(z - c) ^ (A : ℂ)} is Mathlib's complex power with the principal branch; for $\operatorname{Im}z>0$ and $c$ real, $z-c$ lies in the
slit plane (\lean{sub_ofReal_mem_slitPlane}), so all four formulas are holomorphic on the whole upper half-plane.  The four
disjuncts are the four cases $l\ne0,l'\ne0$; $l\ne0,l'=0$; $l=0,l'\ne0$; $l=l'=0$, with $c=-m/l$, $A=l'/l$ or $l'/m$, $M_0=-m'/l'$,
$B=m'/l$ or $m'/m$.
\end{dictionary}
\paragraph{Proof on paper.}  In each case multiply by an explicit integrating factor and show the derivative of the result
vanishes on the connected $S$ (\lean{exists_const_of_hasDerivAt_zero}). \qed
\begin{physics}The real one-parameter subgroups of $\mathrm{SL}_2(\chnineteenR)$ and their ``non-group'' cousins: the power family
with $A=\pm1$ contains the M\"obius maps, the others will be excluded.\end{physics}
\begin{status}No hypotheses beyond those displayed; standard axioms (\leanpath{MoonshineCNGenusZeroAudit2.lean}).
Supporting: \lean{sub_ofReal_mem_slitPlane}, \lean{sub_ofReal_ne_zero}, \lean{exists_const_of_hasDerivAt_zero},
\lean{hasDerivAt_of_differentiableOn}, and the four case theorems \lean{eqOn_power_of_ode}, \lean{eqOn_log_of_ode},
\lean{eqOn_exp_of_ode}, \lean{eqOn_affine_of_ode}.\end{status}

\subsection{Boundary-real exclusion and global extension}

\begin{theorem}[Boundary-real exclusion]\label{ch19:thm:boundary}
Assume the prime Hecke identity for $k$, $N\ge1$.  There is no injective holomorphic $v$ from a half-disc $\Omega^+=D(x_1,r)\cap\chnineteenH$
at a real point $x_1$ into a closed disc $\overline D(z_1,\rho)$ with $\rho<\operatorname{Im}z_1$, such that $f\circ v=f$ on $\Omega^+$.
\end{theorem}
\begin{leancode}
theorem eq_of_comp_eq_of_mapsTo_compact {G : ℂ → ℂ} {U H K : Set ℂ} (hUo : IsOpen U)
    (hUc : IsPreconnected U) (hG : DifferentiableOn ℂ G U) (hHU : H ⊆ U) (hKc : IsCompact K)
    (hKH : K ⊆ H) {τ : ℂ → ℂ} (hτK : ∀ w ∈ H, τ w ∈ K) (hτinj : InjOn τ H)
    (hGτ : ∀ w ∈ H, G (τ w) = G w) {w₀ : ℂ} (hw₀ : w₀ ∈ H) (hw₀K : w₀ ∉ K) :
    ∀ z ∈ U, G z = G w₀
theorem false_of_halfDisc_symmetry {N : ℕ} (hN : 0 < N) (F : ReplicableFamily N) (k : ℕ)
    (hB1 : PrimeHeckeIdentity F k) {x₁ r : ℝ} (hr : 0 < r) {v : ℂ → ℂ}
    (hv : DifferentiableOn ℂ v (Metric.ball (x₁ : ℂ) r ∩ upperHalfPlaneSet))
    (hinj : InjOn v (Metric.ball (x₁ : ℂ) r ∩ upperHalfPlaneSet)) {z₁ : ℂ} {ρ : ℝ}
    (hρz : ρ < z₁.im)
    (hvK : ∀ w ∈ Metric.ball (x₁ : ℂ) r ∩ upperHalfPlaneSet, v w ∈ Metric.closedBall z₁ ρ)
    (hvf : ∀ w ∈ Metric.ball (x₁ : ℂ) r ∩ upperHalfPlaneSet,
      qSeriesExt (F.c k) (v w) = qSeriesExt (F.c k) w) : False
\end{leancode}
(\leanpath{VOA/Moonshine/MoonshineCNGenusZeroBoundary.lean}.)
\begin{dictionary}
\lean{Metric.ball (x₁ : ℂ) r ∩ upperHalfPlaneSet} is the half-disc $\Omega^+$; \lean{Metric.closedBall z₁ ρ} the closed disc, which lies in
$\chnineteenH$ because $\rho<\operatorname{Im}z_1$; \lean{InjOn v H} is injectivity on $\Omega^+$ (Lean's \lean{H}).  The first quoted lemma is the dynamical tool
used in the proof.
\end{dictionary}
\paragraph{Proof on paper.}  Take a large admissible prime $p$ and pair on the connected $\Omega^+$ with $u=\mathrm{id}$ scaled by $p$:
$f(pw)=f(\beta(v(w)))$ on $\Omega^+$ for one $\beta\in\Delta(p)$.  If $\beta=(z\mapsto pz)$: then $f(pw)=f(pv(w))$; $pv(w)$ lies in the top for all $w$
(as $\operatorname{Im}v\ge\operatorname{Im}z_1-\rho>0$ and $p$ is large) and $pw$ lies above the top somewhere in $\Omega^+$; by
Theorem~\ref{ch19:thm:localtop}(c) $v$ is real affine, so it maps points of $\Omega^+$ near $x_1$ to points near $\chnineteenR$, contradicting
$\operatorname{Im}v\ge\operatorname{Im}z_1-\rho>0$.  If $\beta=(z\mapsto(z+j)/p)$: put $\phi(z)=f(pz)$ and, with a suitable integer $m$,
$\tau(w)=((v(w)+j)/p+m)/p$.  By periodicity $\phi\circ\tau=\phi$ on $\Omega^+$; $\tau$ is injective and maps $\Omega^+$ into a closed disc $K$ of radius
$\rho/p^2$, which for $p$ large lies inside $\Omega^+$ near $x_1$.  Take $w_0\in\Omega^+\setminus K$; its iterates under $\tau$ are pairwise distinct
(injectivity and $w_0\notin\tau(\Omega^+)$), lie in the compact $K\subset\chnineteenH$, hence accumulate there, and $\phi$ takes the value $\phi(w_0)$
on all of them; by the identity theorem $\phi$ is constant, so $f$ is constant, which it is not
(\lean{eq_of_comp_eq_of_mapsTo_compact}, \lean{qSeriesExt_exists_ne}). \qed
\begin{physics}A symmetry of the twisted partition function cannot push the moduli of a degenerating torus (a point of the real
axis) to a non-degenerate torus: the degeneration limits are intrinsic.\end{physics}
\begin{status}Hypotheses: $0<N$, \lean{PrimeHeckeIdentity F k}, and those displayed.  Standard axioms
(\leanpath{MoonshineCNGenusZeroAudit2.lean}).  Supporting: \lean{norm_affC_sub_affC}, \lean{injOn_affC_comp},
\lean{qSeriesExt_exists_ne}.\end{status}

\begin{theorem}[Global extension]\label{ch19:thm:extension}
Assume the prime Hecke identity for $k$, $N\ge1$.  Let $\beta$ be holomorphic on $\{\operatorname{Im}>0\}$ with $\beta'$ nowhere zero there,
$\operatorname{Im}z_0>0$, $\operatorname{Im}\beta(z_0)>0$, and $f\circ\beta=f$ near $z_0$.  Then $\operatorname{Im}\beta(z)>0$ and $f(\beta(z))=f(z)$ for
every $z$ with $\operatorname{Im}z>0$.
\end{theorem}
\begin{leancode}
theorem forall_im_pos_and_eq_of_local {N : ℕ} (hN : 0 < N) (F : ReplicableFamily N) (k : ℕ)
    (hB1 : PrimeHeckeIdentity F k) {β : ℂ → ℂ} (hβ : DifferentiableOn ℂ β upperHalfPlaneSet)
    (hβ' : ∀ z : ℂ, 0 < z.im → _root_.deriv β z ≠ 0) {z₀ : ℂ} (hz₀ : 0 < z₀.im)
    (hβz₀ : 0 < (β z₀).im)
    (hloc : ∀ᶠ z in 𝓝 z₀, qSeriesExt (F.c k) (β z) = qSeriesExt (F.c k) z) :
    ∀ z : ℂ, 0 < z.im → 0 < (β z).im ∧ qSeriesExt (F.c k) (β z) = qSeriesExt (F.c k) z
\end{leancode}
(\leanpath{VOA/Moonshine/MoonshineCNGenusZeroExtension.lean}.)
\begin{dictionary}\lean{∀ᶠ z in 𝓝 z₀, ...} is ``on a neighbourhood of $z_0$''.\end{dictionary}
\paragraph{Proof on paper.}  Let $W$ be the connected component through $z_0$ of $\{z\in\chnineteenH:\beta(z)\in\chnineteenH\}$; by the identity
theorem $f\circ\beta=f$ on $W$.  $W$ is open; it is also relatively closed in $\chnineteenH$: at $z_1\in\chnineteenH\cap\overline W$, either
$\beta(z_1)\in\chnineteenH$ (then a disc around $z_1$ is in $W$) or $\beta(z_1)=x_1\in\chnineteenR$.  In the latter case the analytic local inverse $v$
of $\beta$ at $z_1$ maps a half-disc at $x_1$ injectively into a small closed disc around $z_1$ inside $\chnineteenH$, and its image (connected,
meeting $W$) lies in $W$, so $f\circ v=f$ on the half-disc, which Theorem~\ref{ch19:thm:boundary} forbids.  $\chnineteenH$ is connected, so
$W=\chnineteenH$. \qed
\begin{physics}A local symmetry of the twisted partition function, if it is given by a globally defined formula, is a global
symmetry: there is no wall in the moduli space across which it could fail.\end{physics}
\begin{status}Hypotheses as displayed; standard axioms (\leanpath{MoonshineCNGenusZeroAudit2.lean}).\end{status}

\subsection{Classification of the global solutions}

\begin{theorem}[Which explicit solutions are global symmetries]\label{ch19:thm:classification}
Let $\mathfrak c$ be $J$-dominated and $f=\mathtt{qSeriesFun}\ \mathfrak c$.
\begin{enumerate}[label=(\alph*)]
\item No $Ke^{Az}+M_0$ ($A,K\ne0$) maps $\chnineteenH$ into $\chnineteenH$.
\item No $B\log(z-c_0)+K$ maps $\chnineteenH$ into $\chnineteenH$ and satisfies $f\circ\beta=f$ on $\chnineteenH$.
\item If $Bz+K$ does, then $B=1$ and $K\in\Z$.
\item If $K(z-c_0)^A+M_0$ ($A,K\ne0$) does, then either $A=1$ and $K>0$ real, or $A=-1$ and $K<0$ real.
\end{enumerate}
\end{theorem}
\begin{leancode}
theorem false_of_exp_family {A M₀ : ℝ} {K : ℂ} (hA : A ≠ 0) (hK : K ≠ 0)
    (hinto : ∀ z : ℂ, 0 < z.im → 0 < (K * Complex.exp (A * z) + M₀).im) : False
theorem false_of_log_family {c : ℕ → ℂ} (hc : IsJDominated c) {B c₀ : ℝ} {K : ℂ}
    (hglob : ∀ z : ℂ, 0 < z.im → 0 < ((B : ℂ) * Complex.log (z - c₀) + K).im ∧
      qSeriesExt c ((B : ℂ) * Complex.log (z - c₀) + K) = qSeriesExt c z) : False
theorem exists_int_of_affine_family {c : ℕ → ℂ} (hc : IsJDominated c) {B : ℝ} {K : ℂ}
    (hglob : ∀ z : ℂ, 0 < z.im → 0 < ((B : ℂ) * z + K).im ∧
      qSeriesExt c ((B : ℂ) * z + K) = qSeriesExt c z) :
    B = 1 ∧ ∃ n : ℤ, K = n
theorem power_family_cases {c : ℕ → ℂ} (hc : IsJDominated c) {c₀ A M₀ : ℝ} {K : ℂ}
    (hA : A ≠ 0) (hK : K ≠ 0)
    (hglob : ∀ z : ℂ, 0 < z.im → 0 < (K * (z - c₀) ^ (A : ℂ) + M₀).im ∧
      qSeriesExt c (K * (z - c₀) ^ (A : ℂ) + M₀) = qSeriesExt c z) :
    (A = 1 ∧ ∃ a : ℝ, 0 < a ∧ K = a) ∨ (A = -1 ∧ ∃ a : ℝ, a < 0 ∧ K = a)
\end{leancode}
(\leanpath{VOA/Moonshine/MoonshineCNGenusZeroClassification.lean}.)
\paragraph{Proof on paper.}  (a) $z\mapsto Az$ maps $\chnineteenH$ onto a half-plane, which $\exp$ maps onto $\C\setminus\{0\}$; so the function takes
the value $M_0-i\notin\chnineteenH$.  (b) On the vertical line $c_0+iy$, $\beta=B(\log y+i\pi/2)+K$ has constant imaginary part, while
$f(c_0+iy)\to\infty$ as $y\to\infty$; but $f$ is bounded on every horizontal line ($|f|\le e^{2\pi\operatorname{Im}}+\mathtt{jTail}$), so
$f(\beta(c_0+iy))=f(c_0+iy)$ is impossible.  (c) $B>0$ (into $\chnineteenH$); high up, both $z$ and $Bz+K$ are in the top, so they differ by an
integer, forcing $B=1$, $K\in\Z$.  (d) On rays $c_0+re^{i\theta}$, $\operatorname{Im}\beta=r^A|K|\sin(\arg K+A\theta)+\operatorname{Im}M_0$ must be positive
for all $r$; letting $r\to\infty$ or $r\to0$ forces $\sin(\arg K+A\theta)>0$ on $(0,\pi)$, hence $|A|\le1$ (\lean{abs_le_one_of_sin_pos}).
$0<A<1$ contradicts top injectivity on the vertical line through $c_0$ ($\beta$ grows like $y^A$ there while $f(\beta)=f$ forces $\beta$ to differ
from $c_0+iy$ by an integer).  $-1<A<0$: three observations.  (1) On the vertical line $c_0+iy$, $\beta(c_0+iy)=K(iy)^A+M_0$ has imaginary part
$y^A|K|\sin(\arg K+A\pi/2)$ ($M_0$ is real); the sine is positive (take $\theta=\pi/2$ above), and $y^A\to+\infty$ as $y\to0^+$ because
$A<0$.  Since $|f(z)|\ge e^{2\pi\operatorname{Im}z}-\mathtt{jTail}(\operatorname{Im}z)$, the values $f(c_0+iy)=f(\beta(c_0+iy))$ are unbounded
as $y\to0^+$.  (2) On the lines $c_0\pm1+is$, as $s\to0^+$, $\beta$ tends to $P_\pm=K(\pm1)^A+M_0$ (principal branch: $(-1)^A=e^{iA\pi}$),
with $\operatorname{Im}P_+=|K|\sin(\arg K)$ and $\operatorname{Im}P_-=|K|\sin(\arg K+A\pi)$.  Both are $\ge0$ by continuity of the sine
positivity at the endpoints $\theta=0,\pi$ (\lean{sin_nonneg_endpoints}); if both vanished, $\arg K$ and $\arg K+A\pi$ would be multiples of
$\pi$, forcing $A\in\Z$, impossible for $-1<A<0$.  So one of them, call it $P$, lies in $\chnineteenH$ (\lean{exists_tendsto_power_family_of_neg}).
(3) On that line $f(c_0\pm1+is)=f(c_0+is)$ by $1$-periodicity, which is unbounded as $s\to0^+$ by (1); but
$f(c_0\pm1+is)=f(\beta(c_0\pm1+is))\to f(P)$ by continuity, a finite value.  Contradiction.  $A=\pm1$: the into-$\chnineteenH$ condition forces $K$ real of the stated sign. \qed
\begin{physics}Of all the analytic ``flows'' compatible with the infinitesimal symmetry, only modular transformations (real
M\"obius maps of determinant one) preserve the partition function.\end{physics}
\begin{status}Hypotheses as displayed (no Hecke input); standard axioms (\leanpath{MoonshineCNGenusZeroAudit2.lean}).
Supporting: \lean{norm_qSeriesExt_le}, \lean{le_norm_qSeriesExt}, \lean{lt_norm_qSeriesExt_of_le_im}, \lean{ofReal_mul_exp_cpow},
\lean{ofReal_mul_I_cpow}, \lean{im_mul_exp_ofReal_mul_I}, \lean{im_mul_exp_pos_of_power_into}, \lean{abs_le_one_of_sin_pos},
\lean{exists_pos_real_of_mul_into}, \lean{exists_neg_real_of_mul_inv_into}, \lean{im_power_family_vertical},
\lean{false_of_power_family_of_pos_of_lt_one}, \lean{sin_nonneg_endpoints}, \lean{exists_tendsto_power_family_of_neg},
\lean{false_of_power_family_of_neg_of_gt_neg_one}.\end{status}

\subsection{Local symmetries are M\"obius; property P}

\begin{theorem}[Local symmetries are real M\"obius maps]\label{ch19:thm:moebius}
Assume the prime Hecke identity for $k$, $N\ge1$.  Let $\alpha$ be a local symmetry on $D(z_0,\eta)$ ($\eta<\operatorname{Im}z_0$) with
$\operatorname{Im}\alpha\ge m>0$ and $\alpha'(z_0)\ne0$.  Then on $D(z_0,\eta/2)$, $\alpha(z)=(az+b)/(cz+d)$ for real $a,b,c,d$ with $ad-bc=1$ and
$\chnineteenmat abcd\in G(f)$ (in Lean the entries are named \lean{p q r s}).
\end{theorem}
\begin{leancode}
def IsMoebiusSymmetryOn (f : ℍ → ℂ) (α : ℂ → ℂ) (S : Set ℂ) : Prop :=
  ∃ (p q r s : ℝ) (h : p * s - q * r = 1), glOfDetOne p q r s h ∈ symmetryGroup f ∧
    ∀ z ∈ S, α z = (p * z + q) / (r * z + s)
theorem isMoebiusSymmetryOn_of_localSymmetry {N : ℕ} (hN : 0 < N) (F : ReplicableFamily N)
    (k : ℕ) (hB1 : PrimeHeckeIdentity F k) {z₀ : ℂ} {η : ℝ} (hη : 0 < η) (hηz : η < z₀.im)
    {α : ℂ → ℂ} (hα : DifferentiableOn ℂ α (Metric.ball z₀ η)) {m : ℝ} (hm : 0 < m)
    (hαm : ∀ z ∈ Metric.ball z₀ η, m ≤ (α z).im)
    (hαf : ∀ z ∈ Metric.ball z₀ η, qSeriesExt (F.c k) (α z) = qSeriesExt (F.c k) z)
    (hd : _root_.deriv α z₀ ≠ 0) :
    IsMoebiusSymmetryOn (qSeriesFun (F.c k)) α (Metric.ball z₀ (η / 2))
\end{leancode}
(\leanpath{VOA/Moonshine/MoonshineCNGenusZeroMoebius.lean}.)
\begin{dictionary}
\lean{glOfDetOne p q r s h} is the matrix $\chnineteenmat abcd$ (with $a,b,c,d=$ \lean{p}, \lean{q}, \lean{r}, \lean{s}) in $\chnineteenGLR$ (Chapter~\ref{ch:16}).
\lean{IsMoebiusSymmetryOn f α S}: $\alpha$ agrees on $S$ with the M\"obius map of a determinant-one real matrix that lies in $G(f)$.
\end{dictionary}
\paragraph{Proof on paper.}  If $\alpha$ is real affine on $D(z_0,\eta)$, the relation $f(az+b)=f(z)$ extends to $\chnineteenH$ (identity
theorem) and Theorem~\ref{ch19:thm:classification}(c) makes $\alpha$ an integer translation.  Otherwise Theorem~\ref{ch19:thm:overlap} gives an
intertwining sequence, Theorem~\ref{ch19:thm:ode-alpha} an ODE ($\alpha$ is non-constant since $\alpha'(z_0)\ne0$), and
Theorem~\ref{ch19:thm:families} an explicit $\beta$ equal to $\alpha$ on $D(z_0,\eta/2)$, holomorphic on $\chnineteenH$ with nowhere vanishing
derivative.  Theorem~\ref{ch19:thm:extension} makes $\beta$ a global symmetry mapping $\chnineteenH$ into itself, and
Theorem~\ref{ch19:thm:classification} sorts the four families: exponential and logarithmic are impossible; affine gives $z\mapsto z+n$;
power gives either $A=1$, $K>0$, which is again an affine map $z\mapsto Kz+(M_0-Kc_0)$ and hence, by part (c), a translation by an integer,
or $A=-1$, $K<0$.  (Constant solutions, $K=0$ or $B=0$, do not occur because $\beta=\alpha$ near $z_0$ and $\alpha'(z_0)\ne0$.)  The last case
is $z\mapsto M_0+K/(z-c_0)$, the M\"obius map of $\frac1{\sqrt{-K}}\chnineteenmat{M_0}{K-M_0c_0}1{-c_0}$, of determinant one (\lean{isMoebiusSymmetryOn_of_affine},
\lean{isMoebiusSymmetryOn_of_inv}). \qed
\begin{physics}Every local symmetry of a twisted partition function is a modular transformation.\end{physics}
\begin{status}Hypotheses: $0<N$, \lean{PrimeHeckeIdentity F k}, and those displayed.  Standard axioms
(\leanpath{MoonshineCNGenusZeroAudit2.lean}).  Supporting: \lean{coe_glOfDetOne_smul}, \lean{glOfDetOne_mem_symmetryGroup_of},
\lean{isMoebiusSymmetryOn_of_affine}, \lean{isMoebiusSymmetryOn_of_inv}, \lean{deriv_eq_of_eqOn_nhds}.\end{status}

\begin{theorem}[Property P]\label{ch19:thm:propP}
Assume the prime Hecke identity for $k$, $N\ge1$.  If $f(z)=f(w)$, there is $\gamma\in G(f)$ with $\gamma z=w$.
\end{theorem}
\begin{leancode}
theorem exists_mem_symmetryGroup_smul_eq_of_not_both_critical {N : ℕ} (hN : 0 < N)
    (F : ReplicableFamily N) (k : ℕ) (hB1 : PrimeHeckeIdentity F k) {z w : ℍ}
    (hzw : qSeriesFun (F.c k) z = qSeriesFun (F.c k) w)
    (hd : _root_.deriv (qSeriesExt (F.c k)) (z : ℂ) ≠ 0 ∨
      _root_.deriv (qSeriesExt (F.c k)) (w : ℂ) ≠ 0) :
    ∃ γ ∈ symmetryGroup (qSeriesFun (F.c k)), γ • z = w
lemma exists_bound_entries_of_smul_eq {Y R : ℝ} (hY : 0 < Y) :
    ∃ B : ℝ, ∀ γ : GL (Fin 2) ℝ, γ.det.val = 1 → ∀ τ σ : ℍ, γ • τ = σ → Y ≤ τ.im → Y ≤ σ.im →
      ‖(τ : ℂ)‖ ≤ R → ‖(σ : ℂ)‖ ≤ R →
        |γ 0 0| ≤ B ∧ |γ 0 1| ≤ B ∧ |γ 1 0| ≤ B ∧ |γ 1 1| ≤ B
theorem exists_mem_symmetryGroup_smul_eq {N : ℕ} (hN : 0 < N) (F : ReplicableFamily N) (k : ℕ)
    (hB1 : PrimeHeckeIdentity F k) {z w : ℍ}
    (hzw : qSeriesFun (F.c k) z = qSeriesFun (F.c k) w) :
    ∃ γ ∈ symmetryGroup (qSeriesFun (F.c k)), γ • z = w
\end{leancode}
(\leanpath{VOA/Moonshine/MoonshineCNGenusZeroPropertyP.lean}, \leanpath{VOA/Moonshine/MoonshineCNGenusZeroCritical.lean}.)
\begin{dictionary}
\lean{γ • z} is the M\"obius action of $\chnineteenGLR$ on $\chnineteenH$.  The first theorem is the case where $f'$ does not vanish at both
points; the second is a uniform bound used for the doubly critical case; the third is the full property P.
\end{dictionary}
\paragraph{Proof on paper.}  \emph{Non-critical case} ($f'(w)\ne0$): Theorem~\ref{ch19:thm:localtop}(a) gives a local symmetry $\alpha$ near $z$
with $\alpha(z)=w$; $\alpha$ is not constant, so $\alpha'(z')\ne0$ at some $z'$ near $z$; Theorem~\ref{ch19:thm:moebius} near $z'$ gives $\gamma\in G$
with $\alpha=\gamma$ near $z'$, and the identity theorem extends $\alpha=\gamma$ to the disc around $z$, so $\gamma z=w$.  If instead $f'(z)\ne0$, exchange
$z,w$ and use $\gamma^{-1}$.  \emph{Doubly critical case}: critical points of $f$ are isolated (\lean{eventually_deriv_qSeriesExt_ne_zero}),
and by the open mapping theorem there are $w_n\to w$ non-critical and $z_n\to z$ with $f(z_n)=f(w_n)$; the non-critical case gives
$\gamma_n\in G$ with $\gamma_nz_n=w_n$.  The entries of $\gamma_n$ are bounded (from $\operatorname{Im}(\gamma\tau)=\operatorname{Im}\tau/|c\tau+d|^2$ and
the real and imaginary parts of $\sigma(c\tau+d)=a\tau+b$), so a subsequence converges to some $\gamma$ with $\gamma z=w$, and $\gamma\in G$ because $G$
is closed (the invariance $f(\gamma\tau)=f(\tau)$ passes to limits). \qed
\begin{physics}Property P: the twisted partition function separates the points of the moduli space of twisted tori, i.e.\ of
$G\backslash\chnineteenH$.  It is the injectivity half of ``$f$ is a coordinate on the modular curve''.\end{physics}
\begin{status}Hypotheses: $0<N$, \lean{PrimeHeckeIdentity F k}.  Standard axioms (\leanpath{MoonshineCNGenusZeroAudit2.lean}).
The hypothesis-free form is Theorem~\ref{ch19:thm:propPLR}.  Supporting: \lean{moebius_denom_ne_zero},
\lean{exists_mem_symmetryGroup_smul_eq_of_deriv_ne_zero}, \lean{not_eventually_eq_qSeriesExt},
\lean{eventually_deriv_qSeriesExt_ne_zero}, \lean{abs_le_sq_add_one}.\end{status}

\begin{warning}
\textbf{Critical pairs in Cummins--Gannon.}  Cummins--Gannon's
Prop.~2.3 asserts a local symmetry at every coincidence $f(z_1)=f(z_2)$, but when both points are critical one needs equal
multiplicities, which their Lemma~2.4 does not give.  The library never constructs a symmetry at a critical target: it uses
non-critical targets only and recovers doubly critical pairs by approximation and \emph{closedness} of $G(f)$ (an earlier plan used
discreteness; closedness suffices).
\end{warning}

\begin{warning}
\textbf{Dead ends of the local-to-global step.}  (i) Canonical cusp charts do not give
M\"obius-ness: Hecke transports only relate non-top components of $\{|f|>R\}$ to each other, never to the top, so one global input
(here: one intertwining sequence on a fixed disc) is unavoidable.  (ii) One cannot hope for two independent ODEs that would pin $\alpha$ down algebraically: already for $J$ and a M\"obius
symmetry $\gamma$, every real affine $C$ with $\gamma C\gamma^{-1}$ also affine fixes the real point $\gamma^{-1}\infty$, so all intertwining
relations share one fixed point and give the same ODE (details in the note).  The case analysis of
Theorem~\ref{ch19:thm:classification} is therefore necessary.  (iii) Maximal domains of analytic continuation need ``maximal
domain of a transport $=$ transport of the maximal domain'', which Cummins--Gannon use silently; the route above avoids maximal
domains altogether.  (iv) No holomorphic inverse germ is needed for the second transport of the overlap construction, and the
classification does not need bijectivity: $\beta(\chnineteenH)\subseteq\chnineteenH$ and $f\circ\beta=f$ suffice.
\end{warning}

%======================================================================
\section{Genus zero V: the symmetry group}\label{ch19:sec:group}
%======================================================================

\subsection{Shimizu's lemma and the stabiliser of \texorpdfstring{$\infty$}{infinity}}

\begin{theorem}[Shimizu's lemma for symmetry groups]\label{ch19:thm:shimizu}
\begin{enumerate}[label=(\alph*)]
\item If $\psi$ is holomorphic on $\chnineteenH$, invariant under $T=\chnineteenmat1101$ and under some $g$ of determinant $1$ with lower-left
  entry $0<|c|<1$, then $\psi$ is constant.
\item If $f$ is holomorphic and non-constant and $T\in G(f)$, every $g\in G(f)$ with $c\ne0$ has $|c|\ge1$.
\item In particular this holds for every $J$-dominated $q$-series.
\end{enumerate}
\end{theorem}
\begin{leancode}
theorem eq_of_invariant_of_lowerLeft_abs_lt_one {ψ : ℍ → ℂ}
    (hψ : MDifferentiable 𝓘(ℂ) 𝓘(ℂ) ψ) {g : GL (Fin 2) ℝ} (hdet : g.det.val = 1)
    (hc0 : g 1 0 ≠ 0) (hc1 : |g 1 0| < 1) (hT : ∀ τ, ψ (glUpperUnipotent 1 • τ) = ψ τ)
    (hg : ∀ τ, ψ (g • τ) = ψ τ) (τ σ : ℍ) : ψ τ = ψ σ
theorem one_le_abs_lowerLeft_of_mem_symmetryGroup {f : ℍ → ℂ}
    (hf : MDifferentiable 𝓘(ℂ) 𝓘(ℂ) f) (hnc : ∃ τ σ : ℍ, f τ ≠ f σ)
    (hT : glUpperUnipotent 1 ∈ symmetryGroup f) {g : GL (Fin 2) ℝ}
    (hg : g ∈ symmetryGroup f) (hc0 : g 1 0 ≠ 0) : 1 ≤ |g 1 0|
theorem one_le_abs_lowerLeft_qSeriesFun {c : ℕ → ℂ} (hc : IsJDominated c) {g : GL (Fin 2) ℝ}
    (hg : g ∈ symmetryGroup (qSeriesFun c)) (hc0 : g 1 0 ≠ 0) : 1 ≤ |g 1 0|
\end{leancode}
(\leanpath{VOA/Moonshine/MoonshineCNGenusZeroShimizu.lean}.)
\begin{dictionary}
\lean{g 1 0} is the lower-left entry $c$ (rows and columns indexed by \lean{Fin 2} $=\{0,1\}$); \lean{g.det.val} the determinant as a
real number; \lean{glUpperUnipotent 1} is $T$ (Chapter~\ref{ch:16}).  The statement is the classical Shimizu lemma (for a discrete
group containing $T$, every element with $c\ne0$ has $|c|\ge1$), in the form ``invariance under $T$ and an element with $0<|c|<1$
forces constancy'', which needs no discreteness hypothesis.
\end{dictionary}
\paragraph{Proof on paper.}  Shimizu's sequence: $B_0=g$, $B_{n+1}=B_nTB_n^{-1}$.  For $\det M=1$ with first column $(a,c)$,
$MTM^{-1}=\chnineteenmat{1-ac}{a^2}{-c^2}{1+ac}$ (\lean{conj_glUpperUnipotent}); so the lower-left entries satisfy $c_{n+1}=-c_n^2$, hence
$|c_n|\le|c|^{n+1}\to0$, and the upper-left entries $a_{n+1}=1-a_nc_n$ stay bounded and tend to $1$.  The elements $C_n=B_{n+1}T^{-1}$
lie in the group generated by $T$ and $g$, converge to the identity, and act non-trivially (as $c_n\ne0$).  A holomorphic function
invariant under a sequence of non-trivial M\"obius maps converging to the identity is constant (the constancy lemma
\lean{eq_of_invariant_of_entries_tendsto_one} of Chapter~\ref{ch:16}: the invariance relations produce a non-zero real vector field
annihilating $\psi$).  (b) apply (a) to $f$ if $0<|c|<1$.  (c) $T\in G(f)$ by $1$-periodicity, and $f$ is non-constant (it has a pole at
$i\infty$, \lean{qSeriesFun_exists_ne}). \qed
\begin{physics}A modular-invariant quantity cannot be invariant under ``too small'' a twist: a symmetry group containing $T$ and an
element with $|c|<1$ is not discrete, and a holomorphic function with a continuous family of approximate symmetries is constant.
This is why the symmetry group of a non-trivial twisted partition function is discrete.\end{physics}
\begin{status}No hypotheses beyond those displayed; standard axioms (\leanpath{MoonshineCNGenusZeroAudit.lean}).  Supporting:
\lean{shimizuPairFrom}, \lean{shimizuMatrixFrom} and its entry, determinant, recursion and membership lemmas,
\lean{shimizuWitnessFrom} with \lean{_mem}, \lean{_val}, \lean{_det}, the estimates \lean{shimizuPairFrom_snd_ne_zero},
\lean{_snd_abs_le}, \lean{_snd_abs_le_self}, \lean{_fst_abs_le}, \lean{_snd_tendsto}, \lean{_fst_tendsto},
\lean{shimizuWitnessFrom_entries_tendsto}, \lean{shimizuWitnessFrom_moves}, \lean{det_val_eq_one_of_mem_symmetryGroup},
\lean{glUpperUnipotent_one_mem_symmetryGroup_qSeriesFun}, \lean{qSeriesFun_exists_ne}.  The file generalises the Shimizu sequence
of \leanpath{MoonshineCNGammaZeroShimizu.lean} (Chapter~\ref{ch:16}), which starts at $\chnineteenmat10y1$, to an arbitrary start.\end{status}

\begin{theorem}[The stabiliser of $\infty$]\label{ch19:thm:stab}
For $J$-dominated $\mathfrak c$: every $\gamma\in G(f)$ with lower-left entry $0$ acts on $\chnineteenH$ as $\tau\mapsto\tau+n$ for an integer $n$.
\end{theorem}
\begin{leancode}
theorem exists_int_smul_eq_add_of_lowerLeft_eq_zero {c : ℕ → ℂ} (hc : IsJDominated c)
    {γ : GL (Fin 2) ℝ} (hγ : γ ∈ symmetryGroup (qSeriesFun c)) (h10 : γ 1 0 = 0) :
    ∃ n : ℤ, ∀ τ : ℍ, ((γ • τ : ℍ) : ℂ) = τ + n
\end{leancode}
(\leanpath{VOA/Moonshine/MoonshineCNGenusZeroGroup.lean}.)
\paragraph{Proof on paper.}  $\gamma=\chnineteenmat ab0d$ acts as $\tau\mapsto(a/d)\tau+b/d$, a real affine global symmetry with $a/d>0$;
Theorem~\ref{ch19:thm:classification}(c). \qed
\begin{physics}The only symmetries fixing the low-temperature limit are the integer shifts of the modulus: the cusp $\infty$ has width
$1$.\end{physics}
\begin{status}Hypothesis \lean{IsJDominated c}; no Hecke input; standard axioms (\leanpath{MoonshineCNGenusZeroAudit2.lean}).\end{status}

\subsection{Discreteness and finiteness}

\begin{theorem}[Shimura's finiteness lemma]\label{ch19:thm:discrete}
For $J$-dominated $\mathfrak c$, $f=\mathtt{qSeriesFun}\ \mathfrak c$ and $G=G(f)$:
\begin{enumerate}[label=(\alph*)]
\item an element of $G$ with $|a-1|<1$, $|b|<1$, $|c|<1$ is the identity;
\item for every $B$, the set of elements of $G$ with all entries bounded by $B$ is finite;
\item for every $c_0\ne0$, the pairs $(\mathrm{fract}(a/c_0),\mathrm{fract}(d/c_0))$ of elements $\chnineteenmat ab{c_0}d\in G$ form a finite set.
\end{enumerate}
\end{theorem}
\begin{leancode}
theorem eq_one_of_mem_symmetryGroup_of_near {c : ℕ → ℂ} (hc : IsJDominated c)
    {g : GL (Fin 2) ℝ} (hg : g ∈ symmetryGroup (qSeriesFun c)) (h00 : |g 0 0 - 1| < 1)
    (h01 : |g 0 1| < 1) (h10 : |g 1 0| < 1) : g = 1
theorem finite_bounded_mem_symmetryGroup {c : ℕ → ℂ} (hc : IsJDominated c) (B : ℝ) :
    {g : GL (Fin 2) ℝ | g ∈ symmetryGroup (qSeriesFun c) ∧ ∀ i j, |g i j| ≤ B}.Finite
theorem finite_fract_diag_div_lowerLeft {c : ℕ → ℂ} (hc : IsJDominated c) {c₀ : ℝ}
    (hc₀ : c₀ ≠ 0) :
    ((fun g : GL (Fin 2) ℝ => (Int.fract (g 0 0 / c₀), Int.fract (g 1 1 / c₀))) ''
      {g | g ∈ symmetryGroup (qSeriesFun c) ∧ g 1 0 = c₀}).Finite
\end{leancode}
(\leanpath{VOA/Moonshine/MoonshineCNGenusZeroDiscrete.lean}.)
\begin{dictionary}
\lean{Int.fract x} $=x-\lfloor x\rfloor\in[0,1)$.  \lean{Set.Finite} is finiteness of a set; \lean{f '' S} is the image.  (a) is
discreteness of $G$ at the identity (in the coordinates $a,b,c$; $d$ is then determined by $\det=1$).
\end{dictionary}
\paragraph{Proof on paper.}  (a) $|c|<1$ forces $c=0$ by Shimizu; then $\gamma$ is a translation by an integer $n$ (Theorem~\ref{ch19:thm:stab}),
so $a=d=\pm1$ and $b=\pm n$; $|a-1|<1$ gives $a=1$ and $|b|<1$ gives $n=0$.  (b) If infinitely many elements had bounded entries, two of
them $g\ne h$ would be entrywise within $1/(4B+4)$ (Bolzano--Weierstrass in $\chnineteenR^4$), and then $gh^{-1}$ (computed with the adjugate,
\lean{glFinTwo_inv_eq}) would satisfy (a), so $g=h$ (\lean{eq_of_mem_symmetryGroup_of_close}).  (c) Multiplying
$\chnineteenmat ab{c_0}d$ by integer translations on both sides changes $a$ and $d$ by multiples of $c_0$ and does not change the residues;
normalise $a,d$ into $[0,|c_0|)$; then $b=(ad-1)/c_0$ is bounded, and (b) applies. \qed
\begin{physics}The modular symmetries of a twisted partition function form a discrete group, with only finitely many elements of
bounded size: the moduli space of twisted tori is an orbifold, not a smeared-out quotient.\end{physics}
\begin{status}Hypothesis \lean{IsJDominated c}; standard axioms (\leanpath{MoonshineCNGenusZeroAudit3.lean}).  Supporting:
\lean{glFinTwo_mul_apply}, \lean{glFinTwo_det_eq}, \lean{det_entries_of_mem_symmetryGroup}, \lean{glFinTwo_inv_eq},
\lean{conj_glUpperUnipotent}, \lean{glUpperUnipotent_intCast_mem_symmetryGroup}, \lean{neg_one_mem_symmetryGroup},
\lean{abs_mul_sub_mul_le}, \lean{eq_of_mem_symmetryGroup_of_close}, \lean{finite_fract_upperLeft_div_lowerLeft},
\lean{finite_fract_lowerRight_div_lowerLeft}.\end{status}

\subsection{The Hecke step}

\begin{theorem}[Hecke step]\label{ch19:thm:heckestep}
Assume the prime Hecke identity at a prime $p$.  If $g=\chnineteenmat abcd\in G(f)$ with $0<|c|<p$, then for some $j<p$ the matrix
$\chnineteenmat{pa}{b-ja}c{(d-jc)/p}$ lies in $G(f)$.
\end{theorem}
\begin{leancode}
theorem exists_heckeStep_mem_symmetryGroup {N : ℕ} (F : ReplicableFamily N) (k : ℕ) {p : ℕ}
    (hp : p.Prime) (hB1 : PrimeHeckeIdentityAt F k p hp.pos) {g : GL (Fin 2) ℝ}
    (hg : g ∈ symmetryGroup (qSeriesFun (F.c k))) (hc0 : g 1 0 ≠ 0) (hcp : |g 1 0| < p) :
    ∃ j : ℕ, j < p ∧ ∃ h ∈ symmetryGroup (qSeriesFun (F.c k)), h 0 0 = p * g 0 0 ∧
      h 0 1 = g 0 1 - j * g 0 0 ∧ h 1 0 = g 1 0 ∧ h 1 1 = (g 1 1 - j * g 1 0) / p
\end{leancode}
(\leanpath{VOA/Moonshine/MoonshineCNGenusZeroHeckeStep.lean}.)
\paragraph{Proof on paper.}  Pairing on the connected domain $\chnineteenH$ (Theorem~\ref{ch19:thm:partner}(a)) with $u=g$, $v=\mathrm{id}$ and
$\beta_0=(z\mapsto pz)$ gives $\beta\in\Delta(p)$ with $f(p\,gz)=f(\beta z)$ for all $z$, i.e.\ $\chnineteenmat p001g\beta^{-1}$ lies in $G$
(with $\beta=\chnineteenmat p001$ or $\chnineteenmat1j0p$, both of determinant $p$, the product has determinant $p\cdot1\cdot p^{-1}=1$).  If $\beta=\chnineteenmat p001$ this element is $\chnineteenmat a{pb}{c/p}d$, whose lower-left entry has $0<|c/p|<1$, excluded by
Shimizu.  So $\beta=\chnineteenmat1j0p$, and $\chnineteenmat p001\chnineteenmat abcd\chnineteenmat1j0p^{-1}=\chnineteenmat{pa}{b-ja}c{(d-jc)/p}$. \qed
\begin{physics}Hecke closure (Cummins 2002): the symmetry group of a replicable twisted partition function is stable under the
covering correspondence at admissible primes.\end{physics}
\begin{status}Hypotheses as displayed (\lean{PrimeHeckeIdentityAt F k p hp.pos}); standard axioms
(\leanpath{MoonshineCNGenusZeroAudit3.lean}).\end{status}

\subsection{The symmetry group contains a \texorpdfstring{$\Gamma_1(N')$}{Gamma1(N')}}

\begin{theorem}[Level]\label{ch19:thm:level}
Let $F$ have period $N\ge1$, satisfy the prime Hecke identity for $k$, and let $f=f_k$ not be injective modulo translations.  Then:
\begin{enumerate}[label=(\alph*)]
\item (rationality) for $\chnineteenmat abcd\in G$ with $c\ne0$: $a/c\in\Q$ and $c^2\in\Q$;
\item (an integral element) $G$ contains an integral matrix $\chnineteenmat ABCD$ with $C\ne0$, $N\mid C$, $D\equiv1\pmod N$;
\item (integrality of the Hecke step) if $\chnineteenmat ABCD\in G$ is integral, $C\ne0$, $p>|C|$ is a prime with the prime Hecke identity and
  $p\mid D$, then $\chnineteenmat{pA}BC{D/p}\in G$;
\item (a lower unipotent) $\chnineteenmat10L1\in G$ for some integer $L\ne0$ with $N\mid L$;
\item (all first columns) if moreover $L\mid c_0$, $a_0\equiv1\pmod L$ and $\gcd(a_0,c_0)=1$, some integral element of $G$ has first column
  $(a_0,c_0)$;
\item $\Gamma_1(|L|)\le G$; hence $\Gamma_1(N')\le G$ for some $N'\ge1$ with $N\mid N'$.
\end{enumerate}
\end{theorem}
\begin{leancode}
theorem exists_rat_upperLeft_div_lowerLeft (hN : 0 < N) (hB1 : PrimeHeckeIdentity F k)
    {g : GL (Fin 2) ℝ} (hg : g ∈ symmetryGroup (qSeriesFun (F.c k))) (hc0 : g 1 0 ≠ 0) :
    ∃ q : ℚ, g 0 0 / g 1 0 = q
theorem exists_rat_lowerLeft_sq (hN : 0 < N) (hB1 : PrimeHeckeIdentity F k)
    {g : GL (Fin 2) ℝ} (hg : g ∈ symmetryGroup (qSeriesFun (F.c k))) (hc0 : g 1 0 ≠ 0) :
    ∃ q : ℚ, g 1 0 ^ 2 = q
theorem exists_integral_mem_symmetryGroup (hN : 0 < N) (hB1 : PrimeHeckeIdentity F k)
    (hD7 : ¬ IsInjectiveModTranslations (qSeriesFun (F.c k))) :
    ∃ m ∈ symmetryGroup (qSeriesFun (F.c k)), ∃ A B C D : ℤ, m 0 0 = A ∧ m 0 1 = B ∧
      m 1 0 = C ∧ m 1 1 = D ∧ C ≠ 0 ∧ (N : ℤ) ∣ C ∧ D ≡ 1 [ZMOD N]
theorem exists_mem_symmetryGroup_heckeStep_of_integral {p : ℕ} (hp : p.Prime)
    (hB1p : PrimeHeckeIdentityAt F k p hp.pos) {m : GL (Fin 2) ℝ}
    (hm : m ∈ symmetryGroup (qSeriesFun (F.c k))) {A B C D : ℤ} (hA : m 0 0 = A)
    (hB : m 0 1 = B) (hC : m 1 0 = C) (hD : m 1 1 = D) (hC0 : C ≠ 0) (hCp : |C| < p)
    (hpD : (p : ℤ) ∣ D) :
    ∃ m' ∈ symmetryGroup (qSeriesFun (F.c k)), m' 0 0 = p * A ∧ m' 0 1 = B ∧ m' 1 0 = C ∧
      m' 1 1 = D / p
theorem exists_glLowerUnipotent_mem_symmetryGroup (hN : 0 < N) (hB1 : PrimeHeckeIdentity F k)
    (hD7 : ¬ IsInjectiveModTranslations (qSeriesFun (F.c k))) :
    ∃ L : ℤ, L ≠ 0 ∧ (N : ℤ) ∣ L ∧ glLowerUnipotent L ∈ symmetryGroup (qSeriesFun (F.c k))
theorem exists_mem_symmetryGroup_firstColumn (hB1 : PrimeHeckeIdentity F k)
    {L : ℤ} (hNL : (N : ℤ) ∣ L)
    (hLmem : glLowerUnipotent L ∈ symmetryGroup (qSeriesFun (F.c k))) {a₀ c₀ : ℤ}
    (hc₀ : L ∣ c₀) (ha₀ : a₀ ≡ 1 [ZMOD L]) (hcop : IsCoprime a₀ c₀) :
    ∃ m ∈ symmetryGroup (qSeriesFun (F.c k)), ∃ B D : ℤ, m 0 0 = a₀ ∧ m 0 1 = B ∧
      m 1 0 = c₀ ∧ m 1 1 = D
theorem gamma1GL_le_symmetryGroup_of_lowerUnipotent (hB1 : PrimeHeckeIdentity F k)
    {L : ℤ} (hNL : (N : ℤ) ∣ L)
    (hLmem : glLowerUnipotent L ∈ symmetryGroup (qSeriesFun (F.c k))) :
    gamma1GL L.natAbs ≤ symmetryGroup (qSeriesFun (F.c k))
theorem ReplicableFamily.exists_gamma1GL_le_symmetryGroup {N : ℕ} (hN : 0 < N)
    (F : ReplicableFamily N) (k : ℕ) (hB1 : PrimeHeckeIdentity F k)
    (hD7 : ¬ IsInjectiveModTranslations (qSeriesFun (F.c k))) :
    ∃ N' : ℕ, 0 < N' ∧ N ∣ N' ∧ gamma1GL N' ≤ symmetryGroup (qSeriesFun (F.c k))
\end{leancode}
(\leanpath{VOA/Moonshine/MoonshineCNGenusZeroLevel.lean}; section variables \lean{{N : ℕ} (F : ReplicableFamily N) (k : ℕ)}.)
\begin{dictionary}
\lean{Int.ModEq} (\lean{D ≡ 1 [ZMOD N]}) is congruence of integers; \lean{IsCoprime a₀ c₀} is Bezout coprimality, equivalent to
$\gcd=1$ over $\Z$.  \lean{gamma1GL L.natAbs} is $\Gamma_1(|L|)$ inside $\chnineteenGLR$ (Chapter~\ref{ch:16}).  \lean{hD7} is the
nontriviality of the single member $k$.
\end{dictionary}
\paragraph{Proof on paper.}  (a) By the Hecke step at admissible primes $p>|c|$ (infinitely many, by Dirichlet), $G$ contains elements with
lower-left entry $c$ and upper-left entry $pa$.  If $a/c$ were irrational, the residues $\mathrm{fract}(pa/c)$ would be pairwise distinct (their
differences $(p-p')a/c$ are not integers), infinitely many, against Theorem~\ref{ch19:thm:discrete}(c).  For $c^2$: apply this to
$gTg^{-1}=\chnineteenmat{1-ac}{a^2}{-c^2}{1+ac}$: $(1-ac)/(-c^2)=a/c-1/c^2\in\Q$.  (b) Nontriviality gives $\sigma,\tau$ with $f(\sigma)=f(\tau)$ and
$\sigma-\tau\notin\Z$; property P gives $g\in G$ with $g\sigma=\tau$, and $c(g)\ne0$ (otherwise $g$ is an integer translation).  By (a), $ac$, $a^2$, $c^2$
are rational; with $H$ a multiple of $N$ clearing their denominators, $gT^Hg^{-1}=\chnineteenmat{1-Hac}{Ha^2}{-Hc^2}{1+Hac}$ is integral,
$C=-Hc^2\ne0$, $N\mid C$ and $D=1+Hac\equiv1\pmod N$.  (c) The Hecke step gives $\chnineteenmat{pA}{B-jA}C{(D-jC)/p}$ for some $j<p$.  If
$j\ne0$, then $p\nmid jC$, so the new lower-right entry has exact denominator $p$; repeating the step keeps the lower-left entry $C$ and
raises the exact $p$-power of the denominator of the lower-right entry by one each time.  The residues $\mathrm{fract}(d_n/C)$ are then pairwise
distinct, against Theorem~\ref{ch19:thm:discrete}(c).  So $j=0$.  (d) Take the integral element of (b).  By Dirichlet choose a prime $p\equiv
D\pmod{|C|}$, $p>|C|$; then $p\equiv D\equiv1\pmod N$, so $p$ is admissible.  Choose $t$ with $Ct+D=p$; by (c) applied to $\chnineteenmat ABCD T^t$,
$\chnineteenmat{pA}{B'}C1\in G$, and a left translation turns it into $\chnineteenmat10C1$.  (e) $\chnineteenmat10{c_0}1$ is a power of
$\chnineteenmat10L1$.  By Dirichlet choose a prime $q\equiv a_0\pmod{|c_0|}$, $q>|c_0|$ (admissible since $a_0\equiv1\pmod L$ and $N\mid L$), and $t$
with $q\mid1+tc_0$; (c) applied to $\chnineteenmat10{c_0}1T^t=\chnineteenmat1t{c_0}{1+tc_0}$ gives $\chnineteenmat qt{c_0}{(1+tc_0)/q}\in G$, and a left
translation makes the upper-left entry $a_0$.  (f) Every $\gamma\in\Gamma_1(|L|)$ has first column $(a_0,c_0)$ as in (e); if $m\in G$ has the same
first column, $m^{-1}\gamma$ is an integer upper unipotent, which lies in $G$. \qed
\begin{physics}The symmetry group of the twisted partition function contains the full congruence group $\Gamma_1(N')$: the twisted
sector is mapped to itself, with the same twist, by every modular transformation that is the identity modulo the level.  Physically
one expects $\Gamma_1(o(g)h)$-invariance from orbifold modular invariance; here it is derived from the covering relations alone.\end{physics}
\begin{status}Hypotheses: $0<N$, \lean{PrimeHeckeIdentity F k}, \lean{¬ IsInjectiveModTranslations (qSeriesFun (F.c k))}.  Standard
axioms (\leanpath{MoonshineCNGenusZeroAudit3.lean}).  Supporting: \lean{infinite_admissible_primes_gt} (Dirichlet, via Mathlib's
\lean{Nat.forall_exists_prime_gt_and_eq_mod}), \lean{rat_mul_den_eq_num_real}, \lean{glLowerUnipotent_mul}, \lean{glLowerUnipotent_zero},
\lean{glLowerUnipotent_inv}, \lean{glFinTwo_mul_glUpperUnipotent_apply}, \lean{glUpperUnipotent_mul_glFinTwo_apply},
\lean{glLowerUnipotent_intCast_mul_mem}.\end{status}

\begin{warning}
\textbf{Two published steps fail as printed} (a recorded dead end).  (i) Cummins--Gannon's
Prop.~5.1 is not usable for $\Gamma_1$: its first step clears every prime $p\mid c'$, $p\nmid K$, using $n=p^\eta\equiv1\pmod K$, and such $n$
need not be admissible (its prime factor $p$ need not be $\equiv1\pmod N$).  The $\Gamma_1$ version is Cummins (2002), Prop.~2.1.  (ii) The
printed last step of Cummins (2002), Prop.~2.1, ``post-multiplying $\chnineteenmat10{bc}1$ by a translation gives $\chnineteenmat1*{bc}q$ with
$q\equiv a\pmod{bc}$'', cannot hold: right multiplication by translations keeps the lower-right entry $\equiv1\pmod{bc}$.  The corrected step
is (e) above: choose the prime $q\equiv a_0\pmod{|c_0|}$, then $t$ with $q\mid1+tc_0$, apply the integral Hecke step, and translate on the left.
Further simplifications recorded there: Cummins' condition ``$c$ coprime to every admissible $n$'' is not needed for $\Gamma_1(N')\le G$;
Cummins--Gannon's Lemmas 5.4, 5.6, 5.7 are not needed in their printed form (rationality of $a/c$ and $c^2$ suffices, and integrality of
the Hecke step is the direct $p$-adic count of (c)); the prime-clearing loop of Prop.~5.1 is avoided by the conjugation $gT^Hg^{-1}$ of (b).
\end{warning}

%======================================================================
\section[Genus zero VI: cusps, commensurability and the main theorem]{Genus zero VI: cusps, commensurability, surjectivity, and the main theorem}\label{ch19:sec:main}
%======================================================================

\subsection{Behaviour at the other cusps}

Throughout this subsection $\mathfrak c$ is $J$-dominated, $f=\mathtt{qSeriesFun}\ \mathfrak c$ has property P, $G=G(f)$, and $\delta$ is a determinant-one matrix
with $\delta T^h\delta^{-1}\in G$ for some $h>0$ (a parabolic element of $G$ fixing the cusp $\delta\infty$) such that $\delta\infty$ is not
$G$-equivalent to $\infty$, i.e.\ $(\gamma\delta)_{10}\ne0$ for all $\gamma\in G$.

\begin{theorem}[Cusps]\label{ch19:thm:cusps}
\begin{enumerate}[label=(\alph*)]
\item (Shimizu for conjugates; cusp separation) if $\delta X\delta^{-1}\in G$ with $\det X=1$, $X_{10}\ne0$, then $h|X_{10}|\ge1$; and
  $h\,(\gamma\delta)_{10}^2\ge1$ for all $\gamma\in G$.
\item (boundedness) $f\circ\delta$ is bounded on $\{\operatorname{Im}\tau\ge Y\}$ for some $Y$.
\item (cusp value) $f\circ\delta$ is $h$-periodic; its cusp function (the function $Q(e^{2\pi i\tau/h})=f(\delta\tau)$ on the punctured disc) is
  analytic at $0$, and $f(\delta\tau)\to Q(0)$ as $\operatorname{Im}\tau\to\infty$.
\item (not attained) if $f(\delta\tau)\to w$ then $f(\tau_0)\ne w$ for all $\tau_0\in\chnineteenH$.
\item (separation of cusp values) if also $\delta'$ satisfies the hypotheses and $f\circ\delta$, $f\circ\delta'$ tend to the same $w$, then
  $(\delta'^{-1}\gamma\delta)_{10}=0$ for some $\gamma\in G$, i.e.\ $\gamma(\delta\infty)=\delta'\infty$.
\item $|f|\to\infty$ at $i\infty$; and a limit of $f\circ\delta$ at $i\infty$ gives \lean{HasCuspValue} at $\delta\infty$.
\end{enumerate}
\end{theorem}
\begin{leancode}
theorem one_le_mul_abs_lowerLeft_of_conj (hc : IsJDominated c) {δ : GL (Fin 2) ℝ}
    (hδ : 0 < δ.det.val) {h : ℝ} (hh : 0 < h)
    (hpar : δ * glUpperUnipotent h * δ⁻¹ ∈ symmetryGroup (qSeriesFun c)) {X : GL (Fin 2) ℝ}
    (hX : X.det.val = 1) (hXG : δ * X * δ⁻¹ ∈ symmetryGroup (qSeriesFun c))
    (hX0 : X 1 0 ≠ 0) : 1 ≤ h * |X 1 0|
theorem exists_bound_comp_smul_of_other (hc : IsJDominated c)
    (hP : ∀ z w : ℍ, qSeriesFun c z = qSeriesFun c w →
      ∃ γ ∈ symmetryGroup (qSeriesFun c), γ • z = w)
    {δ : GL (Fin 2) ℝ} (hδ : δ.det.val = 1) {h : ℝ} (hh : 0 < h)
    (hpar : δ * glUpperUnipotent h * δ⁻¹ ∈ symmetryGroup (qSeriesFun c))
    (hother : ∀ γ ∈ symmetryGroup (qSeriesFun c), (γ * δ) 1 0 ≠ 0) :
    ∃ Y R : ℝ, ∀ τ : ℍ, Y ≤ τ.im → ‖qSeriesFun c (δ • τ)‖ < R
theorem cuspData_of_other (hc : IsJDominated c)
    (hP : ∀ z w : ℍ, qSeriesFun c z = qSeriesFun c w →
      ∃ γ ∈ symmetryGroup (qSeriesFun c), γ • z = w)
    {δ : GL (Fin 2) ℝ} (hδ : δ.det.val = 1) {h : ℝ} (hh : 0 < h)
    (hpar : δ * glUpperUnipotent h * δ⁻¹ ∈ symmetryGroup (qSeriesFun c))
    (hother : ∀ γ ∈ symmetryGroup (qSeriesFun c), (γ * δ) 1 0 ≠ 0) :
    AnalyticAt ℂ (cuspFunction h (fun τ : ℍ => qSeriesFun c (δ • τ))) 0 ∧
      Tendsto (fun τ : ℍ => qSeriesFun c (δ • τ)) atImInfty
        (𝓝 (cuspFunction h (fun τ : ℍ => qSeriesFun c (δ • τ)) 0)) ∧
      ∀ q : ℂ, q ≠ 0 → ‖q‖ < 1 → ∃ τ : ℍ, τ.im = -h / (2 * π) * Real.log ‖q‖ ∧
        cuspFunction h (fun τ : ℍ => qSeriesFun c (δ • τ)) q = qSeriesFun c (δ • τ)
theorem qSeriesFun_ne_of_tendsto_comp_smul (hc : IsJDominated c)
    (hP : ∀ z w : ℍ, qSeriesFun c z = qSeriesFun c w →
      ∃ γ ∈ symmetryGroup (qSeriesFun c), γ • z = w)
    {δ : GL (Fin 2) ℝ} (hδ : δ.det.val = 1) {h : ℝ} (hh : 0 < h)
    (hpar : δ * glUpperUnipotent h * δ⁻¹ ∈ symmetryGroup (qSeriesFun c))
    (hother : ∀ γ ∈ symmetryGroup (qSeriesFun c), (γ * δ) 1 0 ≠ 0) {w : ℂ}
    (hw : Tendsto (fun τ : ℍ => qSeriesFun c (δ • τ)) atImInfty (𝓝 w)) (τ₀ : ℍ) :
    qSeriesFun c τ₀ ≠ w
theorem exists_mem_lowerLeft_eq_zero_of_tendsto_eq (hc : IsJDominated c)
    (hP : ∀ z w : ℍ, qSeriesFun c z = qSeriesFun c w →
      ∃ γ ∈ symmetryGroup (qSeriesFun c), γ • z = w)
    {δ δ' : GL (Fin 2) ℝ} (hδ : δ.det.val = 1) (hδ' : δ'.det.val = 1) {h h' : ℝ}
    (hh : 0 < h) (hh' : 0 < h')
    (hpar : δ * glUpperUnipotent h * δ⁻¹ ∈ symmetryGroup (qSeriesFun c))
    (hpar' : δ' * glUpperUnipotent h' * δ'⁻¹ ∈ symmetryGroup (qSeriesFun c))
    (hother : ∀ γ ∈ symmetryGroup (qSeriesFun c), (γ * δ) 1 0 ≠ 0)
    (hother' : ∀ γ ∈ symmetryGroup (qSeriesFun c), (γ * δ') 1 0 ≠ 0) {w : ℂ}
    (hw : Tendsto (fun τ : ℍ => qSeriesFun c (δ • τ)) atImInfty (𝓝 w))
    (hw' : Tendsto (fun τ : ℍ => qSeriesFun c (δ' • τ)) atImInfty (𝓝 w)) :
    ∃ γ ∈ symmetryGroup (qSeriesFun c), (δ'⁻¹ * γ * δ) 1 0 = 0
theorem hasCuspValue_of_tendsto_comp_smul {f : ℍ → ℂ} {δ : GL (Fin 2) ℝ} {w : ℂ}
    (hw : Tendsto (fun τ : ℍ => f (δ • τ)) atImInfty (𝓝 w)) :
    HasCuspValue f (δ • cuspInfty) w
\end{leancode}
(\leanpath{VOA/Moonshine/MoonshineCNGenusZeroCusps.lean}; section variable \lean{{c : ℕ → ℂ}}.)
\begin{dictionary}
\lean{cuspFunction h F} is Mathlib's cusp function of an $h$-periodic function $F$ on $\chnineteenH$: the function $Q$ on the disc with
$Q(e^{2\pi i\tau/h})=F(\tau)$, extended at $0$ by the limit (Mathlib, \lean{Function.Periodic.cuspFunction}).  \lean{AnalyticAt ℂ Q 0}
is complex analyticity at $0$.  \lean{HasCuspValue} and \lean{cuspInfty} are the statement-file notions of Chapter~\ref{ch:16};
\lean{δ • cuspInfty} is $\delta\infty$.  The hypothesis \lean{hP} is property P, stated as an explicit hypothesis so that these results
apply to any $J$-dominated series that has it.
\end{dictionary}
\paragraph{Proof on paper.}  (a) $\delta T^{h}\delta^{-1}\in G$ conjugated by $\delta X\delta^{-1}$: Shimizu's lemma for the group $\delta^{-1}G\delta$,
whose translations have width $h$, after conjugating by $\mathrm{diag}(h,1)$ to make the width $1$.  For cusp separation: $(\gamma\delta)T^h(\gamma\delta)^{-1}
=\gamma(\delta T^h\delta^{-1})\gamma^{-1}\in G$ has lower-left entry $-h(\gamma\delta)_{10}^2\ne0$, so by Shimizu $h(\gamma\delta)_{10}^2\ge1$.  Consequently
(\lean{im_le_of_eq_comp_smul}): if $f(z)=f(\delta\tau)$, property P gives $z=\gamma\delta\tau$ and $\operatorname{Im}z\le1/((\gamma\delta)_{10}^2\operatorname{Im}\tau)
\le h/\operatorname{Im}\tau$: \emph{points with the value of $f\circ\delta$ at a high point lie low}.  (b) A large value of $f(\delta\tau)$ would be taken at a
high point (Theorem~\ref{ch19:thm:largevalues}), but such points lie low for $\operatorname{Im}\tau$ large.  (c) $h$-periodicity is the invariance under
$\delta T^h\delta^{-1}$; a bounded holomorphic function on a punctured disc extends holomorphically (Riemann), which is Mathlib's analyticity of the
cusp function of a bounded periodic holomorphic function.  (d) If $f(\tau_0)=w$, by the open mapping theorem the values $f(\delta\tau_n)\to w$ are taken
at points $z_n\to\tau_0$; but such points lie low, $\operatorname{Im}z_n\le h/\operatorname{Im}\tau_n\to0$, a contradiction.  (e) The common value $w$ is the value at $q'=0$ of the cusp function $Q'$ of $f\circ\delta'$, and $Q'$ is non-constant
and analytic at $0$; by the open mapping theorem every value near $w$, in particular every $f(\delta\tau)$ with $\tau$ high, equals $Q'(q')$
for some small $q'$.  By (d), $w$ itself is not a value of $f$, so these values $f(\delta\tau)\ne w$ need $q'\ne0$, i.e.\ they are values
$f(\delta'\tau')$ with $\tau'\in\chnineteenH$ high; property P gives
$X=\delta'^{-1}\gamma\delta$ with $X\tau=\tau'$; if $X_{10}\ne0$, (a) for $\delta'$ and $\operatorname{Im}X\tau\le1/(X_{10}^2\operatorname{Im}\tau)$ contradict
both $\tau,\tau'$ high.  (f) $|f|\ge e^{2\pi\operatorname{Im}}-\mathtt{jTail}$; and two matrices sending $\infty$ to $\delta\infty$ differ by an upper
triangular matrix, which preserves $\operatorname{Im}\to\infty$. \qed
\begin{physics}At every degeneration of the twisted torus other than the vacuum one, the twisted partition function tends to a finite
value (heuristically: in the dual channel no state of negative energy, which would produce a pole, survives), this value is never attained at a non-degenerate torus, and distinct
degenerations have distinct values.\end{physics}
\begin{status}Hypotheses as displayed; standard axioms (\leanpath{MoonshineCNGenusZeroAudit3.lean}).  Supporting:
\lean{im_smul_le_of_det_eq_one}, \lean{det_val_mul_eq_one}, \lean{det_val_inv_eq_one}, \lean{one_le_mul_lowerLeft_sq_of_mem},
\lean{im_le_of_eq_comp_smul}, \lean{periodic_comp_smul_ofComplex}, \lean{tendsto_norm_qSeriesFun_atImInfty}.\end{status}

\subsection{Commensurability and surjectivity}

\begin{theorem}[Commensurability]\label{ch19:thm:index}
If $\mathfrak c$ is $J$-dominated, $f=\mathtt{qSeriesFun}\ \mathfrak c$ has property P and $\Gamma_1(M)\le G$ ($M\ge1$), then $G$ is commensurable with $\chnineteenSL$.
\end{theorem}
\begin{leancode}
theorem exists_smul_eq_smul_fd (G : Subgroup (GL (Fin 2) ℝ)) (τ : ℍ) :
    ∃ r : SL(2, ℤ) ⧸ G.comap (mapGL ℝ), ∃ g : SL(2, ℤ), mapGL ℝ g ∈ G ∧
      ∃ σ ∈ 𝒟, mapGL ℝ g • τ = mapGL ℝ (Quotient.out r)⁻¹ • σ
theorem exists_im_bound_norm_sub_ge (hc : IsJDominated c)
    (hP : ∀ z w : ℍ, qSeriesFun c z = qSeriesFun c w →
      ∃ γ ∈ symmetryGroup (qSeriesFun c), γ • z = w)
    {δ : GL (Fin 2) ℝ} (hδ : δ.det.val = 1) {h : ℝ} (hh : 0 < h)
    (hpar : δ * glUpperUnipotent h * δ⁻¹ ∈ symmetryGroup (qSeriesFun c)) {v : ℂ}
    (hv : (∀ γ ∈ symmetryGroup (qSeriesFun c), (γ * δ) 1 0 ≠ 0) →
      ∀ w, Tendsto (fun τ : ℍ => qSeriesFun c (δ • τ)) atImInfty (𝓝 w) → v ≠ w) :
    ∃ Y ε : ℝ, 0 < ε ∧ ∀ σ : ℍ, Y ≤ σ.im → ε ≤ ‖qSeriesFun c (δ • σ) - v‖
theorem commensurable_symmetryGroup_of_gamma1GL_le (hc : IsJDominated c)
    (hP : ∀ z w : ℍ, qSeriesFun c z = qSeriesFun c w →
      ∃ γ ∈ symmetryGroup (qSeriesFun c), γ • z = w)
    {M : ℕ} (hM : 0 < M) (hΓ1 : gamma1GL M ≤ symmetryGroup (qSeriesFun c)) :
    Subgroup.Commensurable (symmetryGroup (qSeriesFun c)) 𝒮ℒ
\end{leancode}
(\leanpath{VOA/Moonshine/MoonshineCNGenusZeroIndex.lean}.)
\begin{dictionary}
\lean{𝒟} is Mathlib's standard fundamental domain \lean{ModularGroup.fd} of $\chnineteenSL$; \lean{SL(2, ℤ) ⧸ G.comap (mapGL ℝ)} is the coset space
of $G\cap\chnineteenSL$ in $\chnineteenSL$ and \lean{Quotient.out r} a chosen representative.  \lean{Subgroup.Commensurable G 𝒮ℒ}: both relative indices
finite (Chapter~\ref{ch:16}).
\end{dictionary}
\paragraph{Proof on paper.}  $[\chnineteenSL:G\cap\chnineteenSL]<\infty$ because $\Gamma_1(M)\le G\cap\chnineteenSL$ has finite index in $\chnineteenSL$.  For
$[G:G\cap\chnineteenSL]<\infty$: given $\gamma\in G$, move $\gamma i$ into $\mathcal D$ by an element of $\chnineteenSL$ written as $\rho^{-1}h$ with $\rho$ one of the
finitely many coset representatives and $h\in G\cap\chnineteenSL$; the element $\kappa=h\gamma\in G$ maps $i$ to $\rho\sigma$ with $\sigma\in\mathcal D$ and
$f(\rho\sigma)=f(i)$.  Near the cusp $\rho\infty$, $f\circ\rho$ stays away from $f(i)$ (it tends to $\infty$ if $\rho\infty\sim\infty$, and to a cusp value not
attained in $\chnineteenH$ otherwise, Theorem~\ref{ch19:thm:cusps}(d)); so $\sigma$ lies in a truncated, compact part of $\mathcal D$, the entries of $\kappa$ are
bounded, and Shimura's finiteness lemma leaves finitely many $\kappa$.  Hence $G$ is a finite union of cosets $\kappa(G\cap\chnineteenSL)$. \qed
\begin{physics}The group of modular symmetries of the twisted partition function is arithmetic: up to finite index it is the full modular
group.\end{physics}
\begin{status}Hypotheses as displayed; standard axioms (\leanpath{MoonshineCNGenusZeroAudit3.lean}).  Supporting:
\lean{upperSL_natCast_mem_Gamma}, \lean{Gamma_le_Gamma1_of}, \lean{mapGL_upperSL_natCast}, \lean{conj_glUpperUnipotent_mem_of_gamma1GL_le}
(every rational cusp $\rho\infty$ carries the parabolic $\rho T^M\rho^{-1}\in G$), \lean{det_val_mapGL}, \lean{finiteIndex_comap_of_gamma1GL_le},
\lean{exists_im_ge_norm_le_of_isCompact}.\end{status}

\begin{theorem}[Surjectivity onto $\chnineteenP(\C)$]\label{ch19:thm:surj}
If $\mathfrak c$ is $J$-dominated, $f=\mathtt{qSeriesFun}\ \mathfrak c$ has property P, $\Gamma_1(M)\le G$ and $G$ is commensurable with $\chnineteenSL$, then every $w\in\C$ is a value of $f$ on
$\chnineteenH$ or the value of $f$ at an other cusp of $G$.
\end{theorem}
\begin{leancode}
theorem exists_SL2Z_of_isOtherCusp {G : Subgroup (GL (Fin 2) ℝ)}
    (hcomm : Subgroup.Commensurable G 𝒮ℒ) {cc : OnePoint ℝ} (hcc : IsOtherCusp G cc) :
    ∃ ρ : SL(2, ℤ), cc = mapGL ℝ ρ • cuspInfty ∧ ∀ γ ∈ G, (γ * mapGL ℝ ρ) 1 0 ≠ 0
theorem surjective_of_gamma1GL_le (hc : IsJDominated c)
    (hP : ∀ z w : ℍ, qSeriesFun c z = qSeriesFun c w →
      ∃ γ ∈ symmetryGroup (qSeriesFun c), γ • z = w)
    {M : ℕ} (hM : 0 < M) (hΓ1 : gamma1GL M ≤ symmetryGroup (qSeriesFun c))
    (hcomm : Subgroup.Commensurable (symmetryGroup (qSeriesFun c)) 𝒮ℒ) (w : ℂ) :
    (∃ τ : ℍ, qSeriesFun c τ = w) ∨
      ∃ cc, IsOtherCusp (symmetryGroup (qSeriesFun c)) cc ∧ HasCuspValue (qSeriesFun c) cc w
\end{leancode}
(\leanpath{VOA/Moonshine/MoonshineCNGenusZeroSurjective.lean}.)
\begin{dictionary}
\lean{IsOtherCusp} and \lean{HasCuspValue} are the notions of the statement file (Chapter~\ref{ch:16}).  The conclusion is exactly the field
\lean{surjective} of \lean{IsNormalizedHauptmodul}.  The first quoted theorem says that the other cusps of a group commensurable with
$\chnineteenSL$ are the $\rho\infty$, $\rho\in\chnineteenSL$, with $\rho\infty$ not $G$-equivalent to $\infty$; it is how the abstract \lean{IsCusp} of the
statement is converted into the concrete data of the cusp theorems.
\end{dictionary}
\paragraph{Proof on paper.}  Let $V$ be the set of values in question.  $V$ is open: at a value $f(\tau)$ by the open mapping theorem; at a cusp
value by the open mapping theorem for the cusp function at $0$, whose values at $q\ne0$ are values of $f$
(\lean{nhds_le_map_cuspFunction_of_other}).  $V$ is closed, in fact $V=\overline{f(\chnineteenH)}$: cusp values are limits of values of $f$;
conversely if $w\notin V$, every $\tau$ is equivalent to $\rho\sigma$ with $\sigma\in\mathcal D$ and $\rho$ among finitely many representatives; near the cusp
$\rho\infty$, $f\circ\rho$ stays away from $w$ (Theorem~\ref{ch19:thm:index}'s estimate \lean{exists_im_bound_norm_sub_ge}), and on the compact
truncated domain $|f\circ\rho-w|$ has a positive minimum; so $w$ is not a limit of values.  $V\ne\emptyset$ and $\C$ is connected, so $V=\C$. \qed
\begin{physics}Every complex number is the twisted partition function of some (possibly degenerate) twisted torus: together with property
P, the twisted partition function is a bijection from the compactified moduli space to the Riemann sphere.\end{physics}
\begin{status}Hypotheses as displayed; standard axioms (\leanpath{MoonshineCNGenusZeroAudit3.lean}).  Supporting:
\lean{isOtherCusp_smul_infty}, \lean{nhds_le_map_cuspFunction_of_other}.\end{status}

\subsection{The main theorem}

\begin{theorem}[Normalised Hauptmodul of the symmetry group]\label{ch19:thm:haupt}
If $\mathfrak c$ is $J$-dominated, $f=\mathtt{qSeriesFun}\ \mathfrak c$ has property P, and $\Gamma_1(M)\le G(f)$ for some $M\ge1$, then $f$ is a normalised Hauptmodul
of $G(f)$ (all eight fields of \lean{IsNormalizedHauptmodul}).
\end{theorem}
\begin{leancode}
theorem isNormalizedHauptmodul_symmetryGroup {c : ℕ → ℂ} (hc : IsJDominated c)
    (hP : ∀ z w : ℍ, qSeriesFun c z = qSeriesFun c w →
      ∃ γ ∈ symmetryGroup (qSeriesFun c), γ • z = w)
    {M : ℕ} (hM : 0 < M) (hΓ1 : gamma1GL M ≤ symmetryGroup (qSeriesFun c)) :
    IsNormalizedHauptmodul (symmetryGroup (qSeriesFun c)) (qSeriesFun c)
\end{leancode}
(\leanpath{VOA/Moonshine/MoonshineCNGenusZeroMain.lean}.)
\begin{dictionary}
\textbf{Recap.}  \lean{IsNormalizedHauptmodul Γ f} says: $f$ is holomorphic and $\Gamma$-invariant with $f-q^{-1}\to0$ at $i\infty$, has
a finite value at every other cusp, and the induced map from $\Gamma\backslash(\chnineteenH\cup\{\text{cusps}\})$ to $\chnineteenP(\C)$ (sending
$\infty$ to $\infty$) is injective and surjective.  \textbf{Field by field} (Chapter~\ref{ch:16}, definition of \lean{IsNormalizedHauptmodul}): \lean{holomorphic} and \lean{normalized} are
Theorem~\ref{ch19:thm:analytic}; \lean{invariant} is the definition of $G(f)$; \lean{injective} is property P; \lean{cuspFinite},
\lean{cuspValue_not_attained}, \lean{cuspValue_injective} are Theorem~\ref{ch19:thm:cusps}(c), (d), (e), applied at $\delta=\rho\in\chnineteenSL$ with
$h=M$ after converting an abstract other cusp into $\rho\infty$ (Theorem~\ref{ch19:thm:surj}, first part, using commensurability from
Theorem~\ref{ch19:thm:index}); \lean{surjective} is Theorem~\ref{ch19:thm:surj}.  Nothing in the eight fields is weakened: the theorem
concludes exactly the statement-file predicate.
\end{dictionary}
\paragraph{Proof on paper.}  As in the dictionary; the order matters: commensurability (Theorem~\ref{ch19:thm:index}) must be proved before
the cusp fields, because \lean{IsOtherCusp} only provides a parabolic element fixing the cusp, and that the cusp is $\rho\infty$ for some
$\rho\in\chnineteenSL$ comes from commensurability. \qed
\begin{physics}The twisted partition function is a global coordinate on the compactified moduli space of twisted tori, which is
therefore a sphere.\end{physics}
\begin{status}Hypotheses as displayed; standard axioms (\leanpath{MoonshineCNGenusZeroAudit3.lean}).\end{status}

\begin{theorem}[The genus-zero theorem for replicable families]\label{ch19:thm:gz}
Let $N\ge1$ and $F$ a replicable family of period $N$ satisfying the prime Hecke identity for every member, with no member injective modulo
translations.  For every $k$ there are $\Gamma\le\chnineteenGLR$ and $N'\ge1$ with $\Gamma\le\chnineteenSLR$, $\Gamma_1(N')\le\Gamma$, $\Gamma$ commensurable
with $\chnineteenSL$, and $f_k$ a normalised Hauptmodul of $\Gamma$.  (The proof takes $\Gamma=G(f_k)$ and also obtains $N\mid N'$;
neither fact is part of the Lean type.)
\end{theorem}
\begin{leancode}
theorem ReplicableFamily.exists_hauptmodul_gammaOne {N : ℕ} (hN : 0 < N)
    (F : ReplicableFamily N) (hB1 : ∀ k, PrimeHeckeIdentity F k)
    (hD7 : F.NotInjectiveModTranslations) (k : ℕ) :
    ∃ (Γ : Subgroup (GL (Fin 2) ℝ)) (N' : ℕ), 0 < N' ∧ Γ.HasDetOne ∧ gamma1GL N' ≤ Γ ∧
      Subgroup.Commensurable Γ 𝒮ℒ ∧ IsNormalizedHauptmodul Γ (qSeriesFun (F.c k))
\end{leancode}
(\leanpath{VOA/Moonshine/MoonshineCNGenusZeroMain.lean}, namespace \lean{VOA.Moonshine.ConwayNorton}.)
\begin{dictionary}
\textbf{Hypotheses.}  \lean{hN : 0 < N}: positive period.  \lean{hB1 : ∀ k, PrimeHeckeIdentity F k}: discharged for every family by
\lean{M8.ReplicableFamily_primeHeckeIdentity} (Theorem~\ref{ch19:thm:primehecke}).  \lean{hD7 : F.NotInjectiveModTranslations}: discharged
for every family of positive period by \lean{M8.ReplicableFamily_notInjectiveModTranslations} (Theorem~\ref{ch19:thm:nontrivial}).  So the
theorem holds for every replicable family of positive period, although no single Lean declaration states that composite for abstract
families (the composition is made for the McKay--Thompson families in \lean{M8.conwayNorton_of_finOrder_of_completelyReplicable} and
\lean{conwayNortonGammaOne_holds}).

\textbf{Conclusion.}  The clauses are those of \lean{ConwayNortonGammaOne} (Chapter~\ref{ch:16}) with $f_k$ in place of $T_g$.
\lean{Γ.HasDetOne} is $\Gamma\le\chnineteenSLR$.  The existential is witnessed by $G(f_k)$; by the remark in Definition~\ref{ch19:def:interface} this is
the only possible choice up to $\pm1$.

\textbf{How the statement could have been wrong.}  (i) A degenerate witness: \lean{IsNormalizedHauptmodul} has the field \lean{injective},
which pins $\Gamma$ to the full invariance group, so one cannot take a tiny $\Gamma$; and \lean{gamma1GL N' ≤ Γ} with commensurability rules out
$\Gamma$ being only the translations.  (ii) A vacuous hypothesis: the hypotheses are satisfied by the McKay--Thompson families (which are Lean
instances \lean{mckayThompsonReplicableFamily hrep g hg}, Chapter~\ref{ch:16}, given \lean{hrep : CompletelyReplicable}, now the
theorem \lean{completelyReplicable_holds} of Chapter~\ref{ch:22}), and on paper already by $J$ itself (the constant family $c(k,n)=c_J(n)$ of period $1$ is replicable by
Theorem~\ref{ch19:thm:Jrep}; no Lean declaration packages it as a \lean{ReplicableFamily 1}).  (iii) The modular fictions $q^{-1}+\zeta q$ satisfy
replication and are not Hauptmoduln of any group commensurable with $\chnineteenSL$; they are excluded by the clause \lean{zero_eq_J} of
\lean{ReplicableFamily}, which enters the proof exactly once, in the tower (Theorem~\ref{ch19:thm:tower}(b)).
\end{dictionary}
\paragraph{Proof on paper.}  Property P (Theorem~\ref{ch19:thm:propP}), $\Gamma_1(N')\le G$ (Theorem~\ref{ch19:thm:level}), commensurability
(Theorem~\ref{ch19:thm:index}) and Theorem~\ref{ch19:thm:haupt}. \qed
\begin{physics}For every cyclic family of twisted sectors satisfying the covering relations and attached to the moonshine partition
function, each twisted partition function is a Hauptmodul: the genus-zero property is a consequence of the string-theoretic Hecke
relations.\end{physics}
\begin{status}Hypotheses \lean{hN}, \lean{hB1}, \lean{hD7}; standard axioms (\leanpath{MoonshineCNGenusZeroAudit3.lean}; the audit's
docstring notes that \texttt{\#print axioms} does not list \lean{hB1}, \lean{hD7}).  Both are theorems; see the dictionary.\end{status}

\subsection{The Conway--Norton consequences}

\begin{theorem}[Conway--Norton from complete replicability]\label{ch19:thm:cnconsequence}
\begin{enumerate}[label=(\alph*)]
\item For every $g\in\Aut(\Vnat)$ of finite order, $T_g=f_1$ for the McKay--Thompson family of $g$.
\item \lean{CompletelyReplicable}, together with the prime Hecke identity and nontriviality for the McKay--Thompson families, implies
  \lean{ConwayNortonGammaOne}.
\item If moreover every automorphism of $\Vnat$ has finite order, \lean{ConwayNortonConjecture} holds; the same with the two inputs
  stated for all replicable families.
\end{enumerate}
\end{theorem}
\begin{leancode}
lemma mckayThompsonFun_eq_qSeriesFun_family (hrep : CompletelyReplicable) (g : AutVnat)
    (hg : IsOfFinOrder g) :
    mckayThompsonFun g = qSeriesFun ((mckayThompsonReplicableFamily hrep g hg).c 1)
theorem conwayNortonGammaOne_of_completelyReplicable (hrep : CompletelyReplicable)
    (hB1 : ∀ (g : AutVnat) (hg : IsOfFinOrder g) (k : ℕ),
      PrimeHeckeIdentity (mckayThompsonReplicableFamily hrep g hg) k)
    (hD7 : ∀ (g : AutVnat) (hg : IsOfFinOrder g),
      (mckayThompsonReplicableFamily hrep g hg).NotInjectiveModTranslations) :
    ConwayNortonGammaOne
theorem conwayNorton_of_completelyReplicable (hfin : ∀ g : AutVnat, IsOfFinOrder g)
    (hrep : CompletelyReplicable)
    (hB1 : ∀ (g : AutVnat) (hg : IsOfFinOrder g) (k : ℕ),
      PrimeHeckeIdentity (mckayThompsonReplicableFamily hrep g hg) k)
    (hD7 : ∀ (g : AutVnat) (hg : IsOfFinOrder g),
      (mckayThompsonReplicableFamily hrep g hg).NotInjectiveModTranslations) :
    ConwayNortonConjecture
theorem conwayNorton_of_completelyReplicable_of_families (hfin : ∀ g : AutVnat, IsOfFinOrder g)
    (hrep : CompletelyReplicable)
    (hB1 : ∀ (N : ℕ) (F : ReplicableFamily N) (k : ℕ), PrimeHeckeIdentity F k)
    (hD7 : ∀ (N : ℕ), 0 < N → ∀ F : ReplicableFamily N, (∀ k, PrimeHeckeIdentity F k) →
      F.NotInjectiveModTranslations) :
    ConwayNortonConjecture
\end{leancode}
(\leanpath{VOA/Moonshine/MoonshineCNGenusZeroMain.lean}.)
\begin{dictionary}
\lean{mckayThompsonReplicableFamily hrep g hg} is the replicable family $c(k,n)=\tr(g^k|\Vnat_{n+1})$ of period $o(g)$ (Chapter~\ref{ch:16}).
In (c) the passage from $\Gamma_1$ to $\Gamma_0$ is the bridge \lean{conwayNorton_of_gammaOne} of Chapter~\ref{ch:16}, which needs complete
replicability and finite order but not rationality of traces.  The hypotheses \lean{hB1} and \lean{hD7} are discharged by the theorems in
\lean{M8.conwayNorton_of_finOrder_of_completelyReplicable (hfin) (hrep)} (file \leanpath{MoonshineLRConwayNorton.lean}, Chapter~\ref{ch:20}),
and \lean{hrep}, \lean{hfin} by \lean{completelyReplicable_holds} and \lean{M8.autVnat_isOfFinOrder} (Chapters~\ref{ch:20}, \ref{ch:22}).
\end{dictionary}
\paragraph{Proof on paper.}  (a) \lean{mckayThompsonFun_eq_qSeriesFun} and $g^1=g$.  (b) Theorem~\ref{ch19:thm:gz} for the family of $g$, member
$k=1$.  (c) (b) and the bridge. \qed
\begin{physics}The moonshine theorem: every twisted partition function of the moonshine CFT is a genus-zero Hauptmodul.\end{physics}
\begin{status}\textbf{Conditional} on the displayed hypotheses (all discharged elsewhere: \lean{conwayNorton} has no hypotheses,
Chapter~\ref{ch:16}).  Standard axioms (\leanpath{MoonshineCNGenusZeroAudit3.lean}).\end{status}

%======================================================================
\section{How the genus-zero proof compares with the literature}\label{ch19:sec:compare}
%======================================================================

\paragraph{Borcherds' finite check.}  Borcherds (1992) proved complete replicability of all $T_g$ and then observed that a completely
replicable function is determined by finitely many of its coefficients (Norton; Alexander--Cummins--McKay--Simons): it suffices to compare the
first few coefficients of $T_g$, computed from the character table of $\mathbb M$ and the decomposition of the first few $\Vnat_n$, with
those of the $171$ Hauptmoduln proposed by Conway and Norton, each separately known to be a Hauptmodul of a genus-zero group and
completely replicable.  This needs the character table, the list of groups and a case check; none of it is available in Lean, and the
route was rejected.

\paragraph{Cummins--Gannon and Cummins.}  Cummins and Gannon (1997) gave a uniform argument: a completely replicable function with
suitable growth is a Hauptmodul of a genus-zero group (or a degenerate ``modular fiction'').  Their proof: replication $\Rightarrow$ modular
equations (a symmetric polynomial $\Psi_n(X,Y)$ with $\Psi_n(f(\tau),f(\beta\tau))=0$) $\Rightarrow$ coincidences $f(z_1)=f(z_2)$ produce local
symmetries $\Rightarrow$ local symmetries continue analytically along paths to maximal domains, which are invariant under a dense group of affine
maps and hence equal to $\chnineteenH$ or a sector; sectors are excluded; so local symmetries are elements of $\Aut(\chnineteenH)=\mathrm{PSL}_2(\chnineteenR)$
$\Rightarrow$ the symmetry group is discrete, Hecke-closed, contains a congruence subgroup $\Rightarrow$ genus zero.  Cummins (2002) gave the
group-theoretic step for $\Gamma(N)$ and $\Gamma_1$-type conclusions.  Earlier, Martin (1996) studied modular invariance of completely
replicable functions, and the convergence of replicable functions was treated by Kozlov (as cited by Cummins--Gannon).

\paragraph{The library's route.}  It follows Cummins--Gannon's architecture but replaces almost every
analytic step:
\begin{itemize}
\item \emph{Modular equations} are needed only at primes and only as the \emph{pairing} of values (Theorem~\ref{ch19:thm:pairing}); the modular
  polynomial $\Psi_n$, its degree and symmetry, Newton's identities and composite $n$ are never used.  The prime Hecke identity is a pure
  Hecke-algebra computation valid for every family (Theorem~\ref{ch19:thm:primehecke}).
\item \emph{Analyticity} of $T_g$ on $\chnineteenH$ is free from $|c_g(n)|\le c(n)$, replacing Mahler's and Kozlov's convergence results.
\item \emph{Local symmetries}: no analytic continuation, no maximal domains, no density lemma, no $\Aut(\chnineteenH)=\mathrm{PSL}_2(\chnineteenR)$.  One
  intertwining sequence on a fixed disc gives a real linear ODE (Theorems~\ref{ch19:thm:overlap}, \ref{ch19:thm:ode-alpha}); its four explicit
  solution families extend globally by a boundary argument (Theorems~\ref{ch19:thm:boundary}, \ref{ch19:thm:extension}) and are classified by growth
  and top injectivity (Theorem~\ref{ch19:thm:classification}).  The reality of the ODE coefficients (from the rational, hence real, Hecke data) is
  what forces real M\"obius maps.  This replacement is new mathematics relative to the published proofs.
\item \emph{Critical points}: Cummins--Gannon's construction at critical pairs has a gap (equal multiplicities); the library avoids critical
  targets and uses closedness of the symmetry group (Theorem~\ref{ch19:thm:propP}).
\item \emph{Top injectivity} is a three-line Lipschitz estimate with a height depending only on $J$; Cummins--Gannon's inexplicit version, and the
  explicit injective region $\{\operatorname{Im}>1\}$ of the original plan, are not needed.
\item \emph{Nontriviality} (the symmetry group is not just the translations) is new: a tower argument down to $J$ through the replicates,
  using only that $J$'s coefficients are unbounded (Theorems~\ref{ch19:thm:tower}, \ref{ch19:thm:nontrivial}); it needs neither integrality,
  nor the area theorem, nor Galois descent.  It is the place where the clause $c(0,\cdot)=J$ is used, and it excludes the modular fictions.
\item \emph{The congruence subgroup}: rationality of $a/c$ and $c^2$, an integral element by conjugating a translation, the $p$-adic integrality
  count, and a corrected final step of Cummins' Prop.~2.1 (Theorem~\ref{ch19:thm:level}); Cummins--Gannon's Prop.~5.1 is not valid for $\Gamma_1$.
\item \emph{Genus zero} without Riemann surfaces: the eight elementary fields of \lean{IsNormalizedHauptmodul} are proved directly; commensurability
  by Shimura's finiteness lemma on one orbit (no generic fibres), surjectivity by an open-and-closed argument (no sequences, no finiteness of cusp
  classes).
\item \emph{$\Gamma_0$} is not obtained here from rationality of traces (Cummins--Gannon's $\Gamma_0$ case): the library gets $\Gamma_1(N')$ here and
  upgrades to $\Gamma_0$ by the replication-at-primes bridge of Chapter~\ref{ch:16}; rationality of traces is then a corollary (Chapter~\ref{ch:22}).
\end{itemize}
The paper blueprint is part of the authors' working notes (unrefereed by an outside party); the Lean proof is the machine certificate.

%======================================================================
\section{Recorded traps and dead ends}\label{ch19:sec:traps}
%======================================================================

The dead ends specific to single steps are recorded where they occur: the circularity of deriving the product identity from the Lie
algebra (\S\ref{ch19:sec:denominator}), the critical-pair gap and the local-to-global dead ends (\S\ref{ch19:sec:local}), and the two
published steps that fail as printed (\S\ref{ch19:sec:group}).  The remaining ones are collected here.

\begin{warning}
\textbf{The sign convention for momenta} (a recorded dead end).  Borcherds' root $(m,n)$, with root space
$\Vnat_{1+mn}$, is the repository momentum $(m,-n)$ of the lattice \lean{lorentzianLattice} (norm $+2xy$), whose root space is $\Vnat_{1-xy}$.  A product
$\prod_{x,y\ge1}$ over repository momenta is wrong: those root spaces vanish.  In this chapter the conversion happens once, in
\lean{monsterLie_weightAction_trace} (the root space at $\gamma$ is \lean{monsterLieMomentum ![γ.1, -γ.2]}, and $1-\gamma_1(-\gamma_2)=1+\gamma_1\gamma_2$).
\end{warning}

\begin{warning}
\textbf{Modular fictions satisfy replication} (Chapter~\ref{ch:16}).  For every $a\in\C$, $q^{-1}+aq$ satisfies all
replication identities with replicates $q^{-1}+a^kq$; with $a=\zeta_N$ these form a periodic family, injective modulo $\Z$, whose symmetry group is
only the translations.  The genus-zero theorem therefore cannot hold for replicable families not attached to $J$, and the proof must use
$c(0,\cdot)=J$; it does so exactly in the nontriviality argument (Theorem~\ref{ch19:thm:tower}(b)).
\end{warning}

\begin{warning}
\textbf{Complete replicability over all of $\Aut(\Vnat)$ needs finite order.}  The fixed definition
\lean{CompletelyReplicable} quantifies over all $g$; for an element of infinite order (if there were one) $T_g$ need not converge and the
statement could not be proved.  The library therefore proves \lean{CompletelyReplicableOfFinOrder} first and passes to
\lean{CompletelyReplicable} with \lean{hfin}, discharged in Chapter~\ref{ch:20}.  The statement file was not changed.
\end{warning}

\begin{warning}
\textbf{Parts of the original plan that were not formalised.}  The modular polynomials $\Psi_{g,n}$
(degree $\psi(n)$, symmetry), the Hecke identity for composite $n$, the explicit injective region $\{\operatorname{Im}>1\}$, the exponential identity
$p(f(\sigma)-f(\tau))=(1-p/q)\exp(-\sum R_m(\tau)p^m/m)$, ``large values from the top'' by Rouch\'e, the Parseval step of nontriviality, and the
logarithmic Faber identity were all argued on paper in the authors' route notes and then found unnecessary.  They
are \emph{not} Lean theorems; they should not be taken as formalised.  Likewise the product form of the
twisted identity with roots of unity, and its exponential form, are paper statements only.
\end{warning}

\begin{warning}
\textbf{``Product'' means ``series with the product's logarithmic derivative''.}  \lean{coneProd} is defined through the ODE $pA'=A\cdot
\mathrm{coneLog}$ (Definition~\ref{ch19:def:coneprod}); the comparison with literal finite products, proposed as \lean{coneProd_single}, was
not written.  This is a documentation-level gap, not a mathematical one (the classical product satisfies the same ODE), and it does not affect the
use of \lean{knz_product} in Chapter~\ref{ch:20}, which only compares \lean{coneProd}s.
\end{warning}

%======================================================================
\section{The audit modules}\label{ch19:sec:audit}
%======================================================================

The chapter's audit modules contain only \texttt{\#print axioms} commands.  The targeted builds of each (see Chapter~\ref{ch:15}) reported every line as \lean{[propext, Classical.choice, Quot.sound]} and no
\lean{sorryAx}; this chapter did not re-run them.
\begin{center}\small
\begin{tabular}{p{0.42\textwidth}p{0.08\textwidth}p{0.42\textwidth}}
\toprule
module & reports & covers\\
\midrule
\leanpath{MoonshineCNGenusZeroAudit.lean} & 61 & GenusZero\{Analytic, Interface, Top, Shimizu, Pairing\}\\
\leanpath{MoonshineCNGenusZeroAudit2.lean} & 87 & GenusZero\{Partner, Local, Overlap, ODE, Boundary, Families, Extension, Classification,
  Moebius, PropertyP, Critical, Group\}\\
\leanpath{MoonshineCNGenusZeroAudit3.lean} & 68 & GenusZero\{Discrete, HeckeStep, Level, Cusps, Index, Surjective, Main\}\\
\leanpath{MoonshineLRFaberAudit.lean} & 35 & LRFaber, LRFaberGenerating\\
\leanpath{MoonshineLRM5Audit.lean} & 102 & LRM5\{Hecke, Poly, Invariance, Cosets, CosetMap, Replicable, KNZ, Ode, Product\}\\
\leanpath{MoonshineLRHeckeAudit.lean} & 14 & LRHecke\\
\leanpath{MoonshineLRJUnboundedAudit.lean} & 4 & LRJUnbounded\\
\leanpath{MoonshineLRTowerAudit.lean} & 5 & LRTower\\
\leanpath{MoonshineLRCoeffBoundAudit.lean} & 5 & LRCoeffBound\\
\leanpath{MoonshineLRNontrivialAudit.lean} & 10 & LRHeckeExpansion, LRNontrivial, LRPropertyP\\
\bottomrule
\end{tabular}
\end{center}
(The ``reports'' column counts the \texttt{\#print axioms} \emph{commands} in each file in the release, not
mentions of the phrase in docstrings.)  The files
\texttt{MoonshineCNTwisted*} and \texttt{MoonshineCNReplication*} have no audit module of their own; their declarations are printed in
\leanpath{VOA/Moonshine/MoonshineCNAudit18.lean} (a Chapter~\ref{ch:16} file), which itself contains 188 commands.  Building it replays
the audits it imports (\texttt{MoonshineCNAudit} to \texttt{MoonshineCNAudit17}, the $\Gamma_0$ and genus-zero audits); the build
of that whole chain reported 2,268 reports with standard axioms and 0 \lean{sorryAx} (build record, Chapter~\ref{ch:15}; not re-run here).  The \texttt{MoonshineCNGenusZeroAudit3} docstring states explicitly that its headline theorems are conditional on
\lean{hB1} and \lean{hD7}, which \texttt{\#print axioms} does not list.

%======================================================================
\section{Referee's checklist}\label{ch19:sec:checklist}
%======================================================================

``std.'' means \texttt{\#print axioms} $=$ \lean{[propext, Classical.choice, Quot.sound]}.  Namespaces: \lean{M8.} $=$
\lean{VOA.Moonshine.ConwayNorton.M8}; otherwise \lean{VOA.Moonshine.ConwayNorton}.
\begin{center}\small
\begin{longtable}{p{0.26\textwidth}p{0.31\textwidth}p{0.19\textwidth}p{0.14\textwidth}}
\toprule
declaration & intended meaning & hypotheses; status & check\\
\midrule
\endhead
\chnineteenlean{M8.faberPoly}, \chnineteenlean{M8.faberPoly_monic}, \chnineteenlean{M8.coeff_aeval_faberPoly}, \chnineteenlean{M8.eq_faberPoly_of_coeff} &
  Faber polynomials: monic, $F_m(f)=q^{-m}+O(q)$, unique & none; std. & Thm.~\ref{ch19:thm:faber}\\
\chnineteenlean{M8.faberGen_generating}, \chnineteenlean{M8.faberPoly_recursion}, \chnineteenlean{M8.coeff_one_aeval_faberPoly} & $\log G=-\sum F_mp^m/m$;
  recursion; $[q]F_m(f)=mc(m)$ & none; std. & Thm.~\ref{ch19:thm:fabergen}\\
\chnineteenlean{M8.hasSum_qHeckeSum}, \chnineteenlean{M8.hasSum_mul_eval_qSeriesFun} & expansion of $f|T_M$; polynomials of $q$-series &
  $J$-dominated; std. & Thms.~\ref{ch19:thm:qheckesum}, \ref{ch19:thm:polyq}\\
\chnineteenlean{M8.qHeckeSum_smul} & $J|T_M$ is $\chnineteenSL$-invariant & invariance of $f$; std. & Thm.~\ref{ch19:thm:heckeinv}\\
\chnineteenlean{M8.eq_zero_of_invariant_hasSum} & invariant $q$-series without constant term vanishes & none; std. & Thm.~\ref{ch19:thm:vanish}\\
\chnineteenlean{M8.faberPoly_cJ_eval_eq_qHeckeSum}, \chnineteenlean{M8.coeff_faberPoly_cJ_eq_qHeckeCoeff} & $J$ is replicable & \textbf{none}; std. &
  Thm.~\ref{ch19:thm:Jrep}\\
\chnineteenlean{M8.faberGen_map_eq_mul}, \chnineteenlean{M8.knz_log}, \chnineteenlean{M8.knz_log_J} & factorisation; log form of KNZ & none; std. &
  Thm.~\ref{ch19:thm:knzlog}\\
\chnineteenlean{M8.odeSol_spec}, \chnineteenlean{M8.eq_of_ode}, \chnineteenlean{M8.odeSol_add} & $pA'=AL$: existence, uniqueness, exp law & none; std. &
  Thm.~\ref{ch19:thm:ode}\\
\chnineteenlean{M8.knz_product}, \chnineteenlean{M8.coneProd_add}, \chnineteenlean{M8.coneProd_exponent_unique} & $\prod(1-p^mq^n)^{c(mn)}=1-\sum c(a+b-1)p^aq^b$
  (formal; $\prod$ via log-derivative) & \textbf{none}; std. & Thm.~\ref{ch19:thm:knzproduct}\\
\chnineteenlean{TwistedDenominatorIdentity} & $D\tilde u=(1-\tilde u)\tilde P$ (height coordinates) & definition & Def.~\ref{ch19:def:twistedTD}\\
\chnineteenlean{GradedLetters.wordSeries_mul_one_sub_letterSeries}, \chnineteenlean{GradedLetters.eulerDeriv_msetSeries} & characters of $T$ and
  $\mathrm{Sym}$ & none; std. & Thm.~\ref{ch19:thm:characters}\\
\chnineteenlean{sum_eq_sum_of_diag_equiv}, \chnineteenlean{wordPBWEquiv_diag} & transfer of weighted counts through $T(S)\cong U(L)$ & bijectivity; std. &
  Thms.~\ref{ch19:thm:transfer}, \ref{ch19:thm:wordpbw}\\
\chnineteenlean{BorcherdsData.twistedDenominatorIdentity_of_eigenbasis}, \chnineteenlean{BorcherdsData.twistedDenominatorIdentity_of_finOrder} & twisted
  identity for Borcherds data & finite order, traces; std. & Thm.~\ref{ch19:thm:twistedabstract}\\
\chnineteenlean{monsterLie_weightAction_trace}, \chnineteenlean{twistedDenominatorIdentity_of_simpleRootGeneration} & twisted identities of $\chnineteenm$ &
  \textbf{\chnineteenlean{hgen}}, finite order; std. & Thm.~\ref{ch19:thm:twistedmonster}\\
\chnineteenlean{faberPoly_replication}, \chnineteenlean{homogenizedEval_faberPoly} & formal replication & twisted identity; std. & Thm.~\ref{ch19:thm:repformal}\\
\chnineteenlean{faberPoly_eval_mckayThompsonFun} & analytic replication of $T_g$ & finite order, twisted identity; std. & Thm.~\ref{ch19:thm:repanalytic}\\
\chnineteenlean{completelyReplicableOfFinOrder_of_simpleRootGeneration}, \chnineteenlean{completelyReplicable_of_simpleRootGeneration},
  \chnineteenlean{conwayNorton_of_simpleRootGeneration} & complete replicability; Conway--Norton & \textbf{\chnineteenlean{hgen}}, (\chnineteenlean{hfin}); std. &
  Thm.~\ref{ch19:thm:completerep}\\
\chnineteenlean{heckeOrbitSum}, \chnineteenlean{PrimeHeckeIdentity}, \chnineteenlean{IsInjectiveModTranslations}, \chnineteenlean{symmetryGroup} & interface &
  definitions & Def.~\ref{ch19:def:interface}\\
\chnineteenlean{exists_height_injective}, \chnineteenlean{exists_large_values_from_height} & top injectivity; large values from the top & none /
  $J$-dominated; std. & Thms.~\ref{ch19:thm:topinj}, \ref{ch19:thm:largevalues}\\
\chnineteenlean{M8.ReplicableFamily_primeHeckeIdentity} & prime Hecke identity for every family & \textbf{none}; std. & Thm.~\ref{ch19:thm:primehecke}\\
\chnineteenlean{M8.JShift_coeff_unbounded}, \chnineteenlean{M8.norm_coeff_le_of_hasSum_qPar}, \chnineteenlean{M8.not_bounded_faberRemainderCoeff} & $J$
  unbounded; coefficient bound; tower & none / $0<N$; std. & Thms.~\ref{ch19:thm:Junbounded}--\ref{ch19:thm:tower}\\
\chnineteenlean{M8.ReplicableFamily_notInjectiveModTranslations} & nontriviality & $0<N$; std. & Thm.~\ref{ch19:thm:nontrivial}\\
\chnineteenlean{heckeValues_pairing}, \chnineteenlean{exists_heckeAffine_partner} & Hecke pairing (pointwise, on domains) & prime Hecke identity at
  $p$; std. & Thms.~\ref{ch19:thm:pairing}, \ref{ch19:thm:partner}\\
\chnineteenlean{exists_intertwining_of_not_realAffine}, \chnineteenlean{exists_real_linear_ode}, \chnineteenlean{eqOn_explicit_of_real_linear_ode} & overlap;
  ODE; explicit solutions & prime Hecke identity; std. & Thms.~\ref{ch19:thm:overlap}--\ref{ch19:thm:families}\\
\chnineteenlean{false_of_halfDisc_symmetry}, \chnineteenlean{forall_im_pos_and_eq_of_local}, \chnineteenlean{power_family_cases} & boundary exclusion;
  extension; classification & prime Hecke identity; std. & Thms.~\ref{ch19:thm:boundary}--\ref{ch19:thm:classification}\\
\chnineteenlean{isMoebiusSymmetryOn_of_localSymmetry}, \chnineteenlean{exists_mem_symmetryGroup_smul_eq} & local symmetries are M\"obius; property P &
  prime Hecke identity; std. & Thms.~\ref{ch19:thm:moebius}, \ref{ch19:thm:propP}\\
\chnineteenlean{one_le_abs_lowerLeft_qSeriesFun}, \chnineteenlean{exists_int_smul_eq_add_of_lowerLeft_eq_zero}, \chnineteenlean{finite_bounded_mem_symmetryGroup}
  & Shimizu; stabiliser of $\infty$; finiteness & $J$-dominated; std. & Thms.~\ref{ch19:thm:shimizu}--\ref{ch19:thm:discrete}\\
\chnineteenlean{exists_heckeStep_mem_symmetryGroup}, \chnineteenlean{ReplicableFamily.exists_gamma1GL_le_symmetryGroup} & Hecke step;
  $\Gamma_1(N')\le G$ & prime Hecke identity, nontriviality; std. & Thms.~\ref{ch19:thm:heckestep}, \ref{ch19:thm:level}\\
\chnineteenlean{cuspData_of_other}, \chnineteenlean{qSeriesFun_ne_of_tendsto_comp_smul}, \chnineteenlean{exists_mem_lowerLeft_eq_zero_of_tendsto_eq} &
  cusp values exist, not attained, separate cusps & property P; std. & Thm.~\ref{ch19:thm:cusps}\\
\chnineteenlean{commensurable_symmetryGroup_of_gamma1GL_le}, \chnineteenlean{surjective_of_gamma1GL_le} & commensurability; surjectivity & property P,
  $\Gamma_1(M)\le G$; std. & Thms.~\ref{ch19:thm:index}, \ref{ch19:thm:surj}\\
\chnineteenlean{isNormalizedHauptmodul_symmetryGroup} & $f$ is the normalised Hauptmodul of $G(f)$ & property P, $\Gamma_1(M)\le G$; std. &
  Thm.~\ref{ch19:thm:haupt}\\
\chnineteenlean{ReplicableFamily.exists_hauptmodul_gammaOne} & \textbf{genus zero for replicable families} & $0<N$, \chnineteenlean{hB1},
  \chnineteenlean{hD7} (both theorems); std. & Thm.~\ref{ch19:thm:gz}\\
\chnineteenlean{conwayNortonGammaOne_of_completelyReplicable}, \chnineteenlean{conwayNorton_of_completelyReplicable},
  \chnineteenlean{conwayNorton_of_completelyReplicable_of_families} & $\Gamma_1$ and $\Gamma_0$ Conway--Norton & \textbf{conditional}: \chnineteenlean{hrep},
  (\chnineteenlean{hfin}), \chnineteenlean{hB1}, \chnineteenlean{hD7}; std. & Thm.~\ref{ch19:thm:cnconsequence}\\
\bottomrule
\end{longtable}
\end{center}

\noindent\textbf{Files of this chapter and how they are covered.}  All sixty files of the assignment are accounted for.
\emph{Headline declarations individually, supporting declarations collectively}:
\leanpath{MoonshineLRFaber.lean}, \leanpath{MoonshineLRFaberGenerating.lean} (\S\ref{ch19:sec:faber});
\leanpath{MoonshineLRM5Hecke.lean}, \leanpath{MoonshineLRM5Poly.lean}, \leanpath{MoonshineLRM5Invariance.lean},
\leanpath{MoonshineLRM5Cosets.lean}, \leanpath{MoonshineLRM5CosetMap.lean}, \leanpath{MoonshineLRM5Replicable.lean}
(\S\ref{ch19:sec:Jrep}); \leanpath{MoonshineLRM5KNZ.lean}, \leanpath{MoonshineLRM5Ode.lean}, \leanpath{MoonshineLRM5Product.lean}
(\S\ref{ch19:sec:denominator}); \leanpath{MoonshineCNTwistedTransfer.lean}, \leanpath{MoonshineCNTwistedPBW.lean},
\leanpath{MoonshineCNTwistedBases.lean}, \leanpath{MoonshineCNTwistedCombinatorics.lean}, \leanpath{MoonshineCNTwistedCharacters.lean},
\leanpath{MoonshineCNTwistedDenominator.lean} (\S\ref{ch19:sec:twisted}); \leanpath{MoonshineCNReplicationFormal.lean},
\leanpath{MoonshineCNReplicationAnalytic.lean} (\S\ref{ch19:sec:twisted}, \S\ref{ch19:sec:replication});
\leanpath{MoonshineCNGenusZeroInterface.lean} (\S\ref{ch19:sec:gz}); \leanpath{MoonshineCNGenusZeroAnalytic.lean},
\leanpath{MoonshineCNGenusZeroTop.lean} (\S\ref{ch19:sec:analytic}); \leanpath{MoonshineLRHecke.lean},
\leanpath{MoonshineLRHeckeExpansion.lean}, \leanpath{MoonshineLRJUnbounded.lean}, \leanpath{MoonshineLRCoeffBound.lean},
\leanpath{MoonshineLRTower.lean}, \leanpath{MoonshineLRNontrivial.lean}, \leanpath{MoonshineLRPropertyP.lean} (\S\ref{ch19:sec:lrinputs});
\leanpath{MoonshineCNGenusZeroPairing.lean}, \leanpath{MoonshineCNGenusZeroPartner.lean} (\S\ref{ch19:sec:pairing});
\leanpath{MoonshineCNGenusZeroLocal.lean}, \leanpath{MoonshineCNGenusZeroOverlap.lean}, \leanpath{MoonshineCNGenusZeroODE.lean},
\leanpath{MoonshineCNGenusZeroFamilies.lean}, \leanpath{MoonshineCNGenusZeroBoundary.lean}, \leanpath{MoonshineCNGenusZeroExtension.lean},
\leanpath{MoonshineCNGenusZeroClassification.lean}, \leanpath{MoonshineCNGenusZeroMoebius.lean}, \leanpath{MoonshineCNGenusZeroPropertyP.lean},
\leanpath{MoonshineCNGenusZeroCritical.lean} (\S\ref{ch19:sec:local}); \leanpath{MoonshineCNGenusZeroShimizu.lean},
\leanpath{MoonshineCNGenusZeroGroup.lean}, \leanpath{MoonshineCNGenusZeroDiscrete.lean}, \leanpath{MoonshineCNGenusZeroHeckeStep.lean},
\leanpath{MoonshineCNGenusZeroLevel.lean} (\S\ref{ch19:sec:group}); \leanpath{MoonshineCNGenusZeroCusps.lean},
\leanpath{MoonshineCNGenusZeroIndex.lean}, \leanpath{MoonshineCNGenusZeroSurjective.lean}, \leanpath{MoonshineCNGenusZeroMain.lean}
(\S\ref{ch19:sec:main}).  \emph{As a list} (\S\ref{ch19:sec:audit}): \leanpath{MoonshineCNGenusZeroAudit.lean},
\leanpath{MoonshineCNGenusZeroAudit2.lean}, \leanpath{MoonshineCNGenusZeroAudit3.lean}, \leanpath{MoonshineLRM5Audit.lean},
\leanpath{MoonshineLRHeckeAudit.lean}, \leanpath{MoonshineLRFaberAudit.lean}, \leanpath{MoonshineLRJUnboundedAudit.lean},
\leanpath{MoonshineLRNontrivialAudit.lean}, \leanpath{MoonshineLRTowerAudit.lean}, \leanpath{MoonshineLRCoeffBoundAudit.lean}.

%======================================================================
\section{Notes and references}\label{ch19:sec:notes}
%======================================================================

\paragraph{The denominator identity and the Monster Lie algebra.}  R.~E.~Borcherds, \emph{Monstrous moonshine and monstrous Lie
superalgebras}, Invent.\ Math.\ 109 (1992) 405--444: the Monster Lie algebra, the identity \eqref{ch19:eq:borcherds}, the twisted
denominator identities and complete replicability (\S\S7--9).  R.~E.~Borcherds, \emph{Generalized Kac--Moody algebras}, J.\ Algebra 115 (1988)
501--512, for the Weyl--Kac--Borcherds character formula.  V.~G.~Kac, \emph{Infinite Dimensional Lie Algebras}, 3rd ed., Cambridge University
Press, 1990, \S\S10--11, for the Weyl--Kac denominator formula.  E.~Jurisich, \emph{Generalized Kac--Moody Lie algebras, free Lie algebras and the
structure of the Monster Lie algebra}, J.\ Pure Appl.\ Algebra 126 (1998) 233--266, for the free subalgebra used in Chapter~\ref{ch:18}.  The
product formula for $j(p)-j(q)$ was found independently by M.~Koike, S.~P.~Norton and D.~Zagier in the 1980s (hence ``KNZ'').

\paragraph{Replication.}  S.~P.~Norton, \emph{More on moonshine}, in Computational Group Theory (Durham, 1982), Academic Press, 1984,
185--193 (replication formulae).  M.~Koike, \emph{On replication formula and Hecke operators}, Nagoya University preprint (mid-1980s),
relating replication to Hecke operators.  K.~Mahler, \emph{On a class of non-linear functional equations connected with modular functions},
J.\ Austral.\ Math.\ Soc.\ 22 (1976) 65--118 (the recursions for the coefficients of $j$).  D.~Alexander, C.~Cummins, J.~McKay, C.~Simons,
\emph{Completely replicable functions}, in Groups, Combinatorics and Geometry (Durham, 1990), LMS Lecture Note Ser.\ 165, Cambridge University
Press, 1992, 87--98.

\paragraph{Genus zero from replicability.}  C.~J.~Cummins and T.~Gannon, \emph{Modular equations and the genus zero property of moduli
functions}, Invent.\ Math.\ 129 (1997) 413--443 (MPIM preprint 96-72): the design source of \S\S\ref{ch19:sec:gz}--\ref{ch19:sec:main}.
C.~J.~Cummins, \emph{Modular equations and discrete, genus-zero subgroups of $\mathrm{SL}(2,\mathbb R)$ containing $\Gamma(N)$}, Canad.\ Math.\
Bull.\ 45 (2002) 36--45: the group-theoretic proposition (Prop.~2.1) used, in corrected form, in Theorem~\ref{ch19:thm:level}.  Y.~Martin, \emph{On
modular invariance of completely replicable functions}, in Moonshine, the Monster, and Related Topics (South Hadley, 1994), Contemp.\ Math.\
193, AMS, 1996, 263--286.  D.~N.~Kozlov, \emph{On completely replicable functions}, M.Sc.\ thesis, Carleton University (1994), as cited by
Cummins--Gannon for the convergence of replicable functions; this chapter does not depend on it.

\paragraph{Discrete groups.}  H.~Shimizu, \emph{On discontinuous groups operating on the product of the upper half planes}, Ann.\ of Math.\ 77
(1963) 33--71 (Shimizu's lemma).  G.~Shimura, \emph{Introduction to the Arithmetic Theory of Automorphic Functions}, Princeton University Press,
1971, Chapter~1 (discrete subgroups of $\mathrm{SL}_2(\mathbb R)$, cusps, Shimizu's lemma and the finiteness lemma).  F.~Diamond and J.~Shurman,
\emph{A First Course in Modular Forms}, GTM 228, Springer, 2005 (congruence subgroups, Hecke operators).  Dirichlet's theorem on primes in
arithmetic progressions is Mathlib's \lean{Nat.forall_exists_prime_gt_and_eq_mod}.

\paragraph{Generalized moonshine.}  S.~Carnahan, \emph{Generalized moonshine I: genus-zero functions}, Algebra Number Theory 4 (2010) 649--679
(arXiv:0812.3440); \emph{Generalized moonshine II: Borcherds products}, Duke Math.\ J.\ 161 (2012) 893--950 (arXiv:0908.4223); \emph{Generalized
moonshine IV: monstrous Lie algebras}, arXiv:1208.6254: twisted denominator identities and Hecke operators for the twisted-sector Lie algebras,
and a replicability-based route to genus zero.  T.~Gannon, \emph{Moonshine beyond the Monster}, Cambridge University Press, 2006, the standard survey of
replication, Hauptmoduln and generalized moonshine.

\paragraph{Physics.}  P.~Goddard and C.~B.~Thorn, Phys.\ Lett.\ B 40 (1972) 235--238 (no-ghost theorem).  J.~A.~Harvey and G.~W.~Moore,
\emph{Algebras, BPS states, and strings}, Nucl.\ Phys.\ B 463 (1996) 315--368 (arXiv:hep-th/9510182): denominator identities and one-loop string
amplitudes.  R.~Dijkgraaf, G.~Moore, E.~Verlinde, H.~Verlinde, \emph{Elliptic genera of symmetric products and second quantized strings},
Commun.\ Math.\ Phys.\ 185 (1997) 197--209 (arXiv:hep-th/9608096): Hecke operators as second quantisation.  N.~M.~Paquette, D.~Persson,
R.~Volpato, \emph{Monstrous BPS-algebras and the superstring origin of moonshine}, Commun.\ Number Theory Phys.\ 10 (2016) 433--526
(arXiv:1601.05412).  R.~Dijkgraaf, J.~Maldacena, G.~Moore, E.~Verlinde, \emph{A black hole Farey tail}, arXiv:hep-th/0005003.  E.~Witten,
\emph{Three-dimensional gravity revisited}, arXiv:0706.3359 (2007).  J.~F.~R.~Duncan and I.~B.~Frenkel, \emph{Rademacher sums, moonshine and
gravity}, Commun.\ Number Theory Phys.\ 5 (2011) 849--976 (arXiv:0907.4529): the genus-zero property as convergence of Rademacher sums, with
the three-dimensional gravity interpretation.  L.~Dixon, P.~Ginsparg, J.~Harvey, Commun.\ Math.\ Phys.\ 119 (1988) 221--241, for twisted
partition functions of the moonshine orbifold.

\paragraph{The development record.}  The development was recorded in the authors' working notes, which are not part of the release.
The $J$-replicability and product identity, the Faber polynomials, the prime Hecke identity, the unboundedness of $J$'s coefficients and
nontriviality were written on Leonardo Rastelli's side; the twisted identities, replication and the rest of the genus-zero proof on Fedor
Popov's side.  What this chapter states as machine-confirmed rests on the build and audit records of Chapter~\ref{ch:15}.

\endgroup
